\documentclass{amsbook}

%% Load the shared preamble (packages, macros, theorem environments)
%% preamble.tex — Category Theory: A Structural Approach
%% Shared preamble for all chapters and appendices

%% ============================================================
%% PACKAGES
%% ============================================================

%% Mathematics
\usepackage{amsmath}
\usepackage{amssymb}
\usepackage{amsthm}
\usepackage{mathtools}      % extends amsmath

%% Graphics and diagrams
\usepackage{tikz}
\usepackage{tikz-cd}        % commutative diagrams
\usetikzlibrary{arrows.meta, decorations.pathmorphing, positioning, calc, matrix}

%% Layout and typography
\usepackage[margin=1.25in]{geometry}
\usepackage{microtype}      % better justification
\usepackage{graphicx}
\usepackage{booktabs}       % better tables
\usepackage{enumitem}       % fine-grained list control

%% Colors and boxes
\usepackage{xcolor}
\usepackage{tcolorbox}
\tcbuselibrary{skins, breakable}

%% References and hyperlinks
\usepackage[colorlinks=true,
            linkcolor=blue!60!black,
            citecolor=green!50!black,
            urlcolor=blue]{hyperref}
\usepackage{cleveref}

%% Bibliography
\usepackage[backend=biber,
            style=alphabetic,
            maxnames=4,
            doi=false,
            isbn=false]{biblatex}
\addbibresource{bibliography.bib}

%% ============================================================
%% TIKZ-CD SETTINGS
%% ============================================================

\tikzcdset{
  arrow style=math font,
  diagrams={>=Straight Barb},
}
\tikzset{
  every node/.style={font=\small},
}

%% ============================================================
%% THEOREM ENVIRONMENTS
%% ============================================================

\theoremstyle{plain}
\newtheorem{theorem}{Theorem}[chapter]
\newtheorem{proposition}[theorem]{Proposition}
\newtheorem{lemma}[theorem]{Lemma}
\newtheorem{corollary}[theorem]{Corollary}

\theoremstyle{definition}
\newtheorem{definition}[theorem]{Definition}
\newtheorem{example}[theorem]{Example}

\theoremstyle{remark}
\newtheorem{remark}[theorem]{Remark}
\newtheorem{warning}[theorem]{Warning}
\newtheorem{exercise}[theorem]{Exercise}

%% ============================================================
%% TCOLORBOX STYLES
%% ============================================================

%% Concept box (light blue, for key concepts / overviews)
\tcbset{
  conceptbox/.style={
    enhanced,
    colback=blue!7!white,
    colframe=blue!40!black,
    fonttitle=\bfseries,
    breakable,
    before upper={\parindent1.5em},
  }
}

%% ============================================================
%% CUSTOM MATH MACROS
%% ============================================================

%% -- Category names ------------------------------------------
\newcommand{\cat}[1]{\mathbf{#1}}          % generic category

\newcommand{\Set}{\mathbf{Set}}
\newcommand{\Grp}{\mathbf{Grp}}
\newcommand{\Top}{\mathbf{Top}}
\newcommand{\Ab}{\mathbf{Ab}}
\newcommand{\Ring}{\mathbf{Ring}}
\newcommand{\Pos}{\mathbf{Pos}}
\newcommand{\Cat}{\mathbf{Cat}}
\newcommand{\Man}{\mathbf{Man}}
\newcommand{\Vect}{\mathbf{Vect}}

%% -- Morphism operations -------------------------------------
\newcommand{\op}{^{\mathrm{op}}}           % opposite category
\newcommand{\id}{\mathrm{id}}              % identity morphism
\newcommand{\Id}{\mathrm{Id}}              % identity functor
\newcommand{\iso}{\xrightarrow{\sim}}      % isomorphism arrow
\newcommand{\blank}{{-}}                   % placeholder

%% -- Hom and related ----------------------------------------
\DeclareMathOperator{\Hom}{Hom}
\DeclareMathOperator{\End}{End}
\DeclareMathOperator{\Aut}{Aut}
\DeclareMathOperator{\Nat}{Nat}
\DeclareMathOperator{\Fun}{Fun}
\DeclareMathOperator{\ob}{ob}
\DeclareMathOperator{\mor}{mor}
\DeclareMathOperator{\dom}{dom}
\DeclareMathOperator{\cod}{cod}

%% -- Limits and colimits ------------------------------------
\newcommand{\colim}{\operatorname{colim}}
\newcommand{\holim}{\operatorname{holim}}
\newcommand{\hocolim}{\operatorname{hocolim}}

%% -- Kan extensions -----------------------------------------
\newcommand{\Lan}{\operatorname{Lan}}
\newcommand{\Ran}{\operatorname{Ran}}

%% -- Adjunction symbol --------------------------------------
\newcommand{\adj}{\mathrel{\dashv}}        % F ⊣ G

%% -- Yoneda embedding ---------------------------------------
\newcommand{\yo}{\mathbf{y}}

%% -- Tensor product -----------------------------------------
\newcommand{\tensor}{\otimes}

%% -- Number systems ------------------------------------------
\newcommand{\NN}{\mathbb{N}}
\newcommand{\ZZ}{\mathbb{Z}}
\newcommand{\QQ}{\mathbb{Q}}
\newcommand{\RR}{\mathbb{R}}
\newcommand{\CC}{\mathbb{C}}

%% -- Misc ---------------------------------------------------
\newcommand{\defeq}{\coloneqq}            % :=
\newcommand{\eqdef}{\eqqcolon}            % =:

%% ============================================================
%% CLEVEREF NAMES
%% ============================================================

\crefname{theorem}{Theorem}{Theorems}
\crefname{proposition}{Proposition}{Propositions}
\crefname{lemma}{Lemma}{Lemmas}
\crefname{corollary}{Corollary}{Corollaries}
\crefname{definition}{Definition}{Definitions}
\crefname{example}{Example}{Examples}
\crefname{remark}{Remark}{Remarks}
\crefname{exercise}{Exercise}{Exercises}
\crefname{chapter}{Chapter}{Chapters}
\crefname{appendix}{Appendix}{Appendices}
\crefname{section}{Section}{Sections}

%% ============================================================
%% MISC SETTINGS
%% ============================================================

%% Wider display math
\setlength{\jot}{6pt}

%% Penalty settings to reduce widow/orphan lines
\widowpenalty=10000
\clubpenalty=10000


%% ============================================================
%% BOOK METADATA
%% ============================================================

\title{Category Theory:\\A Structural Approach}
\author{Anonymous}
\date{2026}

\begin{document}

\frontmatter

\maketitle

%% ---- Dedication ------------------------------------------
\cleardoublepage
\thispagestyle{empty}
\vspace*{3in}
\begin{center}
  \itshape
  For everyone who wondered why mathematics looks the same everywhere.
\end{center}

%% ---- Preface ---------------------------------------------
\chapter*{Preface}
\addcontentsline{toc}{chapter}{Preface}

This book is an introduction to category theory at the level of a beginning graduate student in mathematics or a mathematically mature advanced undergraduate. Our goal is significant depth: we prove the major theorems in full (or near-full) detail, provide a wide range of examples, and connect the abstract theory to its concrete mathematical origins.

The subject matter divides into two parts. The first six chapters cover the classical core: categories, functors, natural transformations, limits and colimits, and adjoint functors. These topics are essential for anyone who reads modern mathematics, and we treat them with corresponding care. The final four chapters—monads, Kan extensions, enriched category theory, and topos theory—are somewhat more advanced, but are included because they are indispensable for understanding the modern mathematical landscape.

Readers are assumed to be comfortable with proof-writing and to have seen some abstract algebra and topology. No prior category theory is assumed.

We have tried throughout to follow the principle articulated by Mac Lane: show the ideas, not just the proofs. Every major definition is motivated; every major theorem is situated in its broader context. Exercises range from routine verifications to research-adjacent problems.

\bigskip
\noindent\textit{How to use this book.} The first six chapters should be read in order. Chapters 7–10 can be read more selectively: Chapter~7 (Monads) depends on Chapters 5 and 6; Chapter~8 (Kan extensions) depends on Chapters 5 and 6; Chapters 9 and 10 are largely independent of each other and of Chapters 7–8.

\bigskip
\noindent\textit{A note on foundations.} We use set theory informally throughout. Appendix~\ref{app:foundations} discusses the foundational issues (proper classes vs.\ universes) for readers who want precision.

\bigskip
\noindent\textit{Acknowledgments.} The standard references are Mac Lane \cite{MacLane1998}, Riehl \cite{Riehl2016}, Awodey \cite{Awodey2010}, and Leinster \cite{Leinster2014}. Our exposition is deeply indebted to all of these.

%% ---- Table of Contents -----------------------------------
\tableofcontents

%% ============================================================
\mainmatter
%% ============================================================

\chapter{Introduction}

\section{What Is Category Theory?}

Category theory is the study of mathematical structure at its most abstract. Where algebra studies sets equipped with operations, and topology studies sets equipped with notions of proximity, category theory steps back further still and asks: what do these mathematical endeavors have in common? What is the bare minimum structure needed to speak of \emph{objects}, \emph{maps between objects}, and \emph{composition of maps}?

The answer is a \emph{category}: a collection of objects, a collection of morphisms between them, and a rule for composing morphisms that is associative and admits identities. That is all. From this sparse beginning, a surprisingly rich theory unfolds—one that has reshaped large parts of algebra, topology, logic, and theoretical computer science.

To a newcomer, category theory can feel like abstraction for its own sake. Why introduce yet another layer of abstraction when we are already comfortable with groups, rings, topological spaces, and so on? The answer is that categorical thinking is not mere abstraction: it is a \emph{precision instrument} for isolating structural content. When a theorem in group theory and an analogous theorem in ring theory have identical proofs, category theory tells us why—both are instances of a single categorical fact. When two constructions in different fields are "the same in some sense," that sense can almost always be made precise as a functor, a natural transformation, or an adjunction.

Category theory also has a distinctive philosophical character. Rather than studying individual objects in isolation, it insists that objects can only be understood through their relationships to other objects—through the morphisms that connect them. An object is, in a precise sense, determined by what maps into it and what maps out of it. This is the spirit of the Yoneda lemma, perhaps the single most important result in elementary category theory.

\section{A Brief History}

Category theory was born in 1945 with the paper \emph{General Theory of Natural Equivalences} by Samuel Eilenberg and Saunders Mac Lane \cite{EilenbergMacLane1945}. Eilenberg and Mac Lane introduced categories and functors not as the primary objects of study, but as scaffolding for a precise definition of \emph{natural transformation}—a concept they needed to articulate the sense in which a construction in algebraic topology was ``natural.'' The famous quip attributed to Mac Lane was that categories were introduced so that functors could be defined, and functors were introduced so that natural transformations could be defined.

The language spread rapidly. By the 1950s, it had been adopted by the Grothendieck school in algebraic geometry. Alexandre Grothendieck's reformulation of algebraic geometry in categorical terms—introducing schemes, toposes, derived categories, and the yoga of six functors—stands as one of the most ambitious programs in twentieth-century mathematics. His work demonstrated that categorical language was not merely convenient but \emph{necessary} for certain questions.

In the 1960s, F.\ William Lawvere brought category theory into contact with logic and foundations, showing that categories could serve as an alternative foundation for mathematics (elementary topos theory) and that universal algebra could be reformulated in categorical terms (algebraic theories). Daniel Kan introduced adjoint functors and Kan extensions. Saunders Mac Lane's landmark textbook \emph{Categories for the Working Mathematician} \cite{MacLane1998}, first published in 1971, consolidated the foundations of the subject.

Since then, category theory has expanded enormously. Enriched category theory, 2-categories and higher categories, monoidal categories, operads, $(\infty,1)$-categories (quasi-categories and complete Segal spaces)—these developments continue to the present day. Applications have spread to theoretical computer science (domain theory, type theory, the categorical semantics of programming languages), quantum physics (topological quantum field theories and categorical quantum mechanics), linguistics, and beyond.

This book focuses on the classical foundations: categories, functors, natural transformations, limits, adjunctions, monads, Kan extensions, enriched categories, and toposes. Mastery of these concepts provides the vocabulary and intuition needed to engage with the modern literature.

\section{How to Read This Book}

\textbf{Prerequisites.} We assume mathematical maturity at the level of a senior undergraduate: comfort with proof writing, basic abstract algebra (groups, rings, modules), and elementary topology. No prior category theory is assumed. Set-theoretic foundations are discussed in Appendix~\ref{app:foundations}; readers may consult that appendix as needed.

\textbf{Notation.} We use the following conventions throughout.
\begin{itemize}
  \item Categories are denoted by bold letters: $\cat{C}$, $\cat{D}$, $\cat{E}$.
  \item Standard categories have fixed names: $\Set$, $\Grp$, $\Top$, $\Vect_k$, $\Ring$, $\Ab$, $\Cat$.
  \item Functors are denoted by Roman letters: $F$, $G$, $H$.
  \item Natural transformations are denoted by lowercase Greek letters: $\alpha$, $\beta$, $\eta$, $\varepsilon$.
  \item For a category $\cat{C}$, we write $\ob(\cat{C})$ for its collection of objects and $\cat{C}(a,b)$ (also written $\Hom_{\cat{C}}(a,b)$) for the set of morphisms from $a$ to $b$.
  \item Composition of $f: a \to b$ and $g: b \to c$ is written $g \circ f: a \to c$ (diagrammatic order: first $f$, then $g$).
  \item We write $F \dashv G$ to indicate that $F$ is left adjoint to $G$.
\end{itemize}

\textbf{Exercises.} Each chapter ends with exercises ranging from routine verifications to more challenging problems. Harder exercises are marked with a star ($\star$). Particularly demanding research-adjacent exercises carry two stars ($\star\star$). Selected problems are collected in Appendix~\ref{app:exercises}.

\textbf{Diagrams.} Commutative diagrams are used extensively. A diagram is said to \emph{commute} if all directed paths with the same source and target yield equal composites. Commutativity of a diagram is often the most efficient way to express a naturality or universal property condition.

\section{Foundational Remarks}

A recurring issue in category theory is set-theoretic size. If we try to form the ``category of all sets,'' its collection of objects is a proper class, not a set. Similarly, the collection of all groups forms a proper class. This creates potential foundational difficulties.

We handle these issues informally throughout the main text, using the words ``small'' and ``large'' as needed, and relying on the reader's good sense to identify when a construction might require care. The reader who wants a more precise treatment should consult Appendix~\ref{app:foundations}, where we discuss the universe axiom (Grothendieck universes) and its role in making precise the notion of ``locally small category.''

The basic stance of this book is that category theory is best learned in close contact with examples, and examples must be concrete enough to be illuminating. We err on the side of mathematical richness over set-theoretic pedantry.

\section{Chapter Overview}

\begin{description}
  \item[Chapter~\ref{ch:categories} — Categories.] We introduce the central definition and populate it with a wide range of examples: sets, groups, topological spaces, posets, monoids. We study special morphisms (isomorphisms, monomorphisms, epimorphisms), distinguished objects (initial, terminal, zero objects), and the basic constructions of slice categories, comma categories, and opposite categories.

  \item[Chapter~\ref{ch:functors} — Functors.] Functors are structure-preserving maps between categories. We study the Yoneda lemma—the assertion that a category embeds fully faithfully into its own presheaf category—and develop the theory of representable functors and universal properties in categorical language.

  \item[Chapter~\ref{ch:nattrans} — Natural Transformations.] We define natural transformations as morphisms between functors, and show that functors between two fixed categories form a category $[\cat{C}, \cat{D}]$. This leads to the 2-categorical structure of $\Cat$ and the full proof of the Yoneda lemma.

  \item[Chapter~\ref{ch:limits} — Limits and Colimits.] Limits and colimits are the categorical abstractions of ``constructions defined by universal properties.'' Products, equalizers, pullbacks, and their colimit duals are all special cases. We prove that a category with all finite limits is already determined by its terminal object, products, and equalizers.

  \item[Chapter~\ref{ch:adjoints} — Adjoint Functors.] Adjoints are one of the most pervasive structures in mathematics. We give four equivalent definitions, prove the adjoint functor theorem, and show that right adjoints preserve limits and left adjoints preserve colimits.

  \item[Chapter~\ref{ch:monads} — Monads.] Every adjunction gives rise to a monad. Conversely, every monad comes from adjunctions in (at least) two canonical ways: the Kleisli and Eilenberg-Moore constructions. Beck's monadicity theorem characterizes precisely which adjunctions arise from their own induced monad.

  \item[Chapter~\ref{ch:kan} — Kan Extensions.] Kan extensions answer the question: given a functor $F: \cat{C} \to \cat{E}$ and $K: \cat{C} \to \cat{D}$, can we ``extend'' $F$ through $K$ to get a functor $\cat{D} \to \cat{E}$? As Mac Lane famously wrote, ``The notion of Kan extensions subsumes all the other fundamental concepts of category theory.''

  \item[Chapter~\ref{ch:enriched} — Enriched Category Theory.] When the hom-sets of a category are themselves objects of some monoidal category $\cat{V}$, we get a $\cat{V}$-enriched category. This framework unifies preadditive categories, metric spaces (Lawvere), simplicial categories, and more.

  \item[Chapter~\ref{ch:topos} — Topos Theory.] Toposes are categories that behave like categories of sets. They arise naturally as categories of sheaves and as the classifying categories for geometric theories. Their internal logic is intuitionistic, reflecting the constructive nature of sheaf semantics.
\end{description}

\chapter{Categories}
\label{ch:categories}

\section{The Definition of a Category}

We begin with the central definition of the subject.

\begin{definition}[Category]
A \textbf{category} $\cat{C}$ consists of:
\begin{enumerate}
  \item A collection $\ob(\cat{C})$, whose elements are called \textbf{objects}.
  \item For each pair of objects $a, b \in \ob(\cat{C})$, a set $\cat{C}(a,b)$, called the \textbf{hom-set} from $a$ to $b$. Elements of $\cat{C}(a,b)$ are called \textbf{morphisms} (or \textbf{arrows}) from $a$ to $b$, and we write $f: a \to b$ to indicate $f \in \cat{C}(a,b)$.
  \item For each triple of objects $a, b, c \in \ob(\cat{C})$, a function
  \[
    \circ \;:\; \cat{C}(b,c) \times \cat{C}(a,b) \;\longrightarrow\; \cat{C}(a,c),
  \]
  called \textbf{composition}. The composite of $f: a \to b$ and $g: b \to c$ is written $g \circ f: a \to c$.
  \item For each object $a \in \ob(\cat{C})$, a morphism $\id_a \in \cat{C}(a,a)$, called the \textbf{identity morphism} on $a$.
\end{enumerate}
These data are required to satisfy:
\begin{description}
  \item[Associativity.] For all $f: a \to b$, $g: b \to c$, $h: c \to d$:
  \[
    h \circ (g \circ f) = (h \circ g) \circ f.
  \]
  \item[Unit laws.] For all $f: a \to b$:
  \[
    f \circ \id_a = f \qquad \text{and} \qquad \id_b \circ f = f.
  \]
\end{description}
\end{definition}

\begin{remark}
We sometimes write $\Hom_{\cat{C}}(a,b)$ for $\cat{C}(a,b)$, especially when the category must be made explicit.  We also write $a \in \cat{C}$ to mean $a \in \ob(\cat{C})$ when no confusion arises.
\end{remark}

The definition looks sparse, but it encodes a great deal. Note that morphisms need not be functions; objects need not be sets. The definition is a purely formal structure. Its power comes entirely from how many mathematical situations it captures.

\begin{remark}[Diagrammatic vs.\ classical order]
Some authors write composition in diagrammatic order: if $f: a \to b$ and $g: b \to c$, then $f;g$ or $f \cdot g$ denotes the composite ``first $f$, then $g$.'' We use the classical (algebraic) order $g \circ f$ throughout, consistent with function composition notation.
\end{remark}

\section{First Examples}
\label{sec:cat-examples}

The best way to absorb the definition is through examples.

\begin{example}[$\Set$]
The category $\Set$ has sets as objects and functions as morphisms. Composition is ordinary function composition, and identities are identity functions. This is the prototypical example; much of category theory is modeled on—and later abstracted from—the category of sets.
\end{example}

\begin{example}[$\Grp$, $\Ab$, $\Ring$]
The category $\Grp$ has groups as objects and group homomorphisms as morphisms. Similarly:
\begin{itemize}
  \item $\Ab$: abelian groups and group homomorphisms.
  \item $\Ring$: rings (with unity) and ring homomorphisms preserving unity.
  \item $\Vect_k$: vector spaces over a field $k$ and $k$-linear maps.
  \item $\mathbf{Mod}_R$: left $R$-modules over a ring $R$ and $R$-module homomorphisms.
\end{itemize}
In each case, composition is composition of the underlying functions, and identities are identity functions (which are automatically homomorphisms).
\end{example}

\begin{example}[$\Top$]
The category $\Top$ has topological spaces as objects and continuous maps as morphisms. The category $\mathbf{hTop}$ (the \emph{homotopy category}) has topological spaces as objects but morphisms are homotopy classes of continuous maps. Composition is well-defined because the composite of homotopic maps is homotopic.
\end{example}

\begin{example}[Posets as categories]
\label{ex:poset-category}
Let $(P, \leq)$ be a partially ordered set. We define a category $\cat{P}$ as follows:
\begin{itemize}
  \item Objects: elements of $P$.
  \item Morphisms: $\cat{P}(a,b) = \{*\}$ (a one-element set, i.e., exactly one morphism) if $a \leq b$, and $\cat{P}(a,b) = \emptyset$ otherwise.
  \item Composition: if $a \leq b$ and $b \leq c$, transitivity gives $a \leq c$, providing the unique composite.
  \item Identities: reflexivity gives $a \leq a$, providing the identity morphism.
\end{itemize}
The categorical axioms then encode exactly the axioms of a partial order: reflexivity (identities), transitivity (composition), and the unit/associativity laws hold automatically since all hom-sets have at most one element. Category theory thus generalizes order theory.
\end{example}

\begin{example}[Monoids as categories]
\label{ex:monoid-category}
Let $(M, \cdot, e)$ be a monoid. We define a category $\mathbf{B}M$ (the \emph{delooping} of $M$) with:
\begin{itemize}
  \item One object, denoted $\bullet$.
  \item Morphisms: $\mathbf{B}M(\bullet, \bullet) = M$.
  \item Composition: the monoid operation $\cdot$.
  \item Identity: the monoid unit $e$.
\end{itemize}
The category axioms for $\mathbf{B}M$ are exactly the monoid axioms for $M$. A monoid is therefore a \emph{one-object category}. Groups correspond to one-object categories in which every morphism is an isomorphism (see Section~\ref{sec:special-morphisms}).
\end{example}

\begin{example}[Discrete categories]
For any set $S$, the \textbf{discrete category} on $S$ has $S$ as its object set and only identity morphisms: $\cat{C}(a,b) = \emptyset$ for $a \neq b$, and $\cat{C}(a,a) = \{\id_a\}$. Discrete categories are in bijection with sets.
\end{example}

\begin{example}[The terminal category $\mathbf{1}$]
The category $\mathbf{1}$ has a single object $\bullet$ and only the identity morphism $\id_\bullet$. It is the terminal object in $\Cat$.
\end{example}

\begin{example}[The arrow category $\mathbf{2}$]
The category $\mathbf{2}$ has two objects $0$ and $1$, identity morphisms on each, and a single non-identity morphism $f: 0 \to 1$.
\end{example}

\begin{example}[Matrix categories]
Fix a commutative ring $R$. Define a category $\mathbf{Mat}_R$ with:
\begin{itemize}
  \item Objects: natural numbers $0, 1, 2, \ldots$
  \item Morphisms: $\mathbf{Mat}_R(m, n)$ is the set of $n \times m$ matrices over $R$.
  \item Composition: matrix multiplication.
  \item Identity: the $n \times n$ identity matrix.
\end{itemize}
This category is equivalent to the category of finitely generated free $R$-modules.
\end{example}

\begin{example}[The fundamental groupoid]
Let $X$ be a topological space. The \textbf{fundamental groupoid} $\Pi_1(X)$ has:
\begin{itemize}
  \item Objects: points of $X$.
  \item Morphisms $x \to y$: homotopy classes (rel endpoints) of paths from $x$ to $y$.
  \item Composition: concatenation of paths (homotopy classes thereof).
  \item Identity at $x$: the constant path at $x$.
\end{itemize}
When $X$ is path-connected and a basepoint $x_0$ is chosen, $\Pi_1(X)(x_0, x_0) = \pi_1(X, x_0)$, the fundamental group.
\end{example}

\section{The Category of Categories}

Small categories themselves form a category.

\begin{definition}
A category $\cat{C}$ is \textbf{small} if $\ob(\cat{C})$ is a set (not a proper class) and each hom-set $\cat{C}(a,b)$ is a set. It is \textbf{locally small} if each hom-set is a set, even if $\ob(\cat{C})$ is a proper class.
\end{definition}

The categories $\Set$, $\Grp$, $\Top$, etc., are locally small but not small. The category $\Cat$ of all small categories and functors between them is locally small but not small.

\begin{definition}[Functor—preliminary definition]
A \textbf{functor} $F: \cat{C} \to \cat{D}$ is a structure-preserving map between categories. (We give the full definition in Chapter~\ref{ch:functors}.)
\end{definition}

\begin{definition}[$\Cat$]
The category $\Cat$ has small categories as objects and functors as morphisms. Composition is composition of functors, and identities are identity functors.
\end{definition}

\section{Opposite Category}

Every category gives rise to a ``mirror image'' category by reversing all arrows.

\begin{definition}[Opposite category]
Let $\cat{C}$ be a category. The \textbf{opposite category} (or \textbf{dual category}) $\cat{C}\op$ is defined by:
\begin{itemize}
  \item $\ob(\cat{C}\op) = \ob(\cat{C})$.
  \item $\cat{C}\op(a,b) = \cat{C}(b,a)$ (morphisms are reversed).
  \item Composition: the composite of $f: a \to b$ and $g: b \to c$ in $\cat{C}\op$ is $f \circ g$ (computed in $\cat{C}$).
  \item Identities: same as in $\cat{C}$.
\end{itemize}
\end{definition}

The opposite category is a powerful conceptual tool: every theorem about categories has a dual theorem about $\cat{C}\op$, obtained by reversing all arrows. This is the \textbf{duality principle}. When we prove a result about limits, we automatically get the dual result about colimits ``for free.''

\begin{remark}
Note that $(\cat{C}\op)\op = \cat{C}$. The operation $(-)^\mathrm{op}$ is an involution on categories.
\end{remark}

\section{Product Categories}

\begin{definition}[Product category]
Given categories $\cat{C}$ and $\cat{D}$, their \textbf{product category} $\cat{C} \times \cat{D}$ has:
\begin{itemize}
  \item Objects: pairs $(c, d)$ with $c \in \cat{C}$ and $d \in \cat{D}$.
  \item Morphisms: $(f,g): (c,d) \to (c',d')$ where $f: c \to c'$ in $\cat{C}$ and $g: d \to d'$ in $\cat{D}$.
  \item Composition: $(f', g') \circ (f, g) = (f' \circ f,\, g' \circ g)$.
  \item Identities: $(\id_c, \id_d)$.
\end{itemize}
\end{definition}

Product categories are used to express the notion of a \emph{bifunctor}—a functor of two variables—which will be central when we discuss the hom-functor and tensor products.

\section{Slice and Coslice Categories}

Slice and coslice categories are ways of viewing a category ``from the perspective of a fixed object.''

\begin{definition}[Slice category]
Let $\cat{C}$ be a category and $c \in \cat{C}$ an object. The \textbf{slice category} $\cat{C}/c$ (read: ``$\cat{C}$ over $c$'') has:
\begin{itemize}
  \item Objects: pairs $(a, f)$ where $a \in \cat{C}$ and $f: a \to c$.
  \item Morphisms $(a, f) \to (a', f')$: morphisms $h: a \to a'$ in $\cat{C}$ such that $f' \circ h = f$, i.e., the triangle
  \[
    \begin{tikzcd}
      a \arrow[rr, "h"] \arrow[dr, "f"'] & & a' \arrow[dl, "f'"] \\
      & c &
    \end{tikzcd}
  \]
  commutes.
  \item Composition and identities: inherited from $\cat{C}$.
\end{itemize}
\end{definition}

\begin{definition}[Coslice category]
The \textbf{coslice category} $c/\cat{C}$ (read: ``$c$ under $\cat{C}$'') has:
\begin{itemize}
  \item Objects: pairs $(f, a)$ where $f: c \to a$.
  \item Morphisms $(f, a) \to (f', a')$: morphisms $h: a \to a'$ such that $h \circ f = f'$.
\end{itemize}
\end{definition}

\begin{example}
The category $\Set_*$ of \emph{pointed sets} (sets with a chosen basepoint) is the coslice $\{*\}/\Set$, where $\{*\}$ is a one-element set and morphisms are basepoint-preserving functions.
\end{example}

\section{Comma Categories}

Comma categories generalize slice categories and serve as the natural habitat for many universal properties.

\begin{definition}[Comma category]
Let $F: \cat{A} \to \cat{C}$ and $G: \cat{B} \to \cat{C}$ be functors. The \textbf{comma category} $(F \downarrow G)$ has:
\begin{itemize}
  \item Objects: triples $(a, b, h)$ where $a \in \cat{A}$, $b \in \cat{B}$, and $h: Fa \to Gb$ in $\cat{C}$.
  \item Morphisms $(a,b,h) \to (a',b',h')$: pairs $(f: a \to a', g: b \to b')$ such that the square
  \[
    \begin{tikzcd}
      Fa \arrow[r, "h"] \arrow[d, "Ff"'] & Gb \arrow[d, "Gg"] \\
      Fa' \arrow[r, "h'"'] & Gb'
    \end{tikzcd}
  \]
  commutes in $\cat{C}$.
  \item Composition and identities: componentwise.
\end{itemize}
\end{definition}

\begin{example}
When $F = \Id_{\cat{C}}: \cat{C} \to \cat{C}$ is the identity and $G$ is the constant functor at $c$, the comma category $(\Id_{\cat{C}} \downarrow G)$ is the slice category $\cat{C}/c$.
\end{example}

Comma categories will reappear prominently when we formulate adjunctions and Kan extensions.

\section{Subcategories}

\begin{definition}[Subcategory]
A category $\cat{D}$ is a \textbf{subcategory} of $\cat{C}$ if:
\begin{itemize}
  \item $\ob(\cat{D}) \subseteq \ob(\cat{C})$.
  \item For all $a, b \in \cat{D}$, $\cat{D}(a,b) \subseteq \cat{C}(a,b)$.
  \item $\cat{D}$ is closed under composition and contains all identity morphisms of its objects.
\end{itemize}
$\cat{D}$ is a \textbf{full subcategory} if $\cat{D}(a,b) = \cat{C}(a,b)$ for all $a, b \in \cat{D}$ (all morphisms between objects of $\cat{D}$ are included). It is a \textbf{wide subcategory} if $\ob(\cat{D}) = \ob(\cat{C})$ (all objects are included).
\end{definition}

\begin{example}
$\Ab$ is a full subcategory of $\Grp$. The category of bijections is a wide subcategory of $\Set$ (same objects, only bijections as morphisms). The category of finite sets $\mathbf{FinSet}$ is a full subcategory of $\Set$.
\end{example}

\section{Special Morphisms}
\label{sec:special-morphisms}

Just as in algebra we distinguish invertible elements, injective functions, and surjective functions, in category theory we have categorical analogues.

\begin{definition}[Isomorphism]
A morphism $f: a \to b$ in a category $\cat{C}$ is an \textbf{isomorphism} if there exists a morphism $g: b \to a$ such that $g \circ f = \id_a$ and $f \circ g = \id_b$. Such $g$ is unique and is called the \textbf{inverse} of $f$, written $f^{-1}$.

Objects $a$ and $b$ are \textbf{isomorphic}, written $a \cong b$, if there exists an isomorphism between them.
\end{definition}

\begin{example}
In $\Set$, isomorphisms are bijections. In $\Grp$, they are group isomorphisms. In $\Top$, they are homeomorphisms. In a poset (Example~\ref{ex:poset-category}), $a \cong b$ iff $a \leq b$ and $b \leq a$, i.e., $a = b$ (since $\leq$ is antisymmetric).
\end{example}

\begin{definition}[Monomorphism]
A morphism $f: a \to b$ is a \textbf{monomorphism} (or \textbf{mono}) if for all objects $c$ and all $g, h: c \to a$,
\[
  f \circ g = f \circ h \implies g = h.
\]
\end{definition}

\begin{definition}[Epimorphism]
A morphism $f: a \to b$ is an \textbf{epimorphism} (or \textbf{epi}) if for all objects $c$ and all $g, h: b \to c$,
\[
  g \circ f = h \circ f \implies g = h.
\]
\end{definition}

\begin{remark}
In $\Set$, monomorphisms are injections and epimorphisms are surjections. In $\Grp$, monomorphisms are injective homomorphisms, but epimorphisms need not be surjective! The inclusion $\mathbb{Z} \hookrightarrow \mathbb{Q}$ is an epimorphism in $\Ring$ (since any two ring homomorphisms $\mathbb{Q} \to R$ agreeing on $\mathbb{Z}$ must agree on all of $\mathbb{Q}$), even though it is not surjective.
\end{remark}

\begin{proposition}
Every isomorphism is both a monomorphism and an epimorphism. The converse fails in general.
\end{proposition}

\begin{proof}
Let $f: a \to b$ be an isomorphism with inverse $f^{-1}$. If $f \circ g = f \circ h$, then $f^{-1} \circ f \circ g = f^{-1} \circ f \circ h$, so $g = h$; thus $f$ is mono. Dually, $f$ is epi. The converse fails: in $\Ring$, the inclusion $\mathbb{Z} \hookrightarrow \mathbb{Q}$ is mono and epi but not an isomorphism.
\end{proof}

\begin{definition}[Section and retraction]
A morphism $s: a \to b$ is a \textbf{section} (or \textbf{split monomorphism}) if there exists $r: b \to a$ with $r \circ s = \id_a$. We call $r$ a \textbf{retraction} (or \textbf{split epimorphism}) of $b$ onto $a$. Every section is a monomorphism; every retraction is an epimorphism.
\end{definition}

\begin{definition}[Endomorphism, automorphism]
A morphism $f: a \to a$ is an \textbf{endomorphism} of $a$. An invertible endomorphism is an \textbf{automorphism}. The set $\End(a) = \cat{C}(a,a)$ forms a monoid under composition; the automorphisms $\Aut(a) \subseteq \End(a)$ form a group.
\end{definition}

\section{Distinguished Objects}

\begin{definition}[Initial and terminal objects]
An object $0 \in \cat{C}$ is \textbf{initial} if for every object $c \in \cat{C}$ there is a unique morphism $0 \to c$.

An object $1 \in \cat{C}$ is \textbf{terminal} (or \textbf{final}) if for every object $c \in \cat{C}$ there is a unique morphism $c \to 1$.
\end{definition}

\begin{proposition}
Initial objects are unique up to unique isomorphism. Dually, terminal objects are unique up to unique isomorphism.
\end{proposition}

\begin{proof}
Suppose $0$ and $0'$ are both initial. Then there exist unique morphisms $f: 0 \to 0'$ and $g: 0' \to 0$. The composite $g \circ f: 0 \to 0$ is a morphism from $0$ to $0$; since $0$ is initial, this composite must equal $\id_0$. Similarly $f \circ g = \id_{0'}$. So $f$ is an isomorphism. The isomorphism is unique because each of $f$ and $g$ is the unique morphism between the respective objects.
\end{proof}

\begin{example}
\leavevmode
\begin{itemize}
  \item In $\Set$: the empty set $\emptyset$ is initial; any one-element set $\{*\}$ is terminal.
  \item In $\Grp$: the trivial group $\{e\}$ is both initial and terminal (a \emph{zero object}).
  \item In $\Ring$: the ring $\mathbb{Z}$ is initial (there is a unique ring homomorphism $\mathbb{Z} \to R$ for any ring $R$); the zero ring $\{0\}$ is terminal.
  \item In a poset category (Example~\ref{ex:poset-category}): initial objects are least elements; terminal objects are greatest elements.
  \item In $\Top$: the empty space is initial; any one-point space is terminal.
\end{itemize}
\end{example}

\begin{definition}[Zero object]
An object that is both initial and terminal is called a \textbf{zero object}. Categories with zero objects include $\Grp$, $\Ab$, $\Vect_k$, and $\mathbf{Mod}_R$.
\end{definition}

When a zero object $0$ exists, for any objects $a, b$ there is a canonical morphism $a \to 0 \to b$ called the \textbf{zero morphism} from $a$ to $b$.

\section{Groupoids}

\begin{definition}[Groupoid]
A \textbf{groupoid} is a category in which every morphism is an isomorphism.
\end{definition}

\begin{example}
\leavevmode
\begin{itemize}
  \item Any group $G$ gives a one-object groupoid $\mathbf{B}G$ (the delooping of $G$).
  \item The fundamental groupoid $\Pi_1(X)$ of a topological space.
  \item A set $S$ gives a discrete groupoid (only identity morphisms, each of which is trivially its own inverse).
  \item The \emph{pair groupoid} on a set $S$ has $\ob = S$ and a unique isomorphism $s \to t$ for every pair $s, t \in S$.
\end{itemize}
\end{example}

Groupoids unify the notions of groups (one-object case) and equivalence relations (when each hom-set has at most one element). They play an important role in algebraic topology and homotopy theory.

\section{The Homotopy Hypothesis and Higher Categories}

\begin{remark}[Looking ahead]
The categories studied in this book are ordinary (``1-'') categories. But the structural perspective of category theory naturally leads to higher-dimensional generalizations. A \textbf{2-category} has objects, morphisms, and \emph{2-morphisms} (morphisms between morphisms). The category $\Cat$ itself carries a 2-categorical structure: functors are the morphisms, and natural transformations (Chapter~\ref{ch:nattrans}) are the 2-morphisms.

The \textbf{homotopy hypothesis} of Grothendieck states that (weak) $\infty$-groupoids (where morphisms at every level are invertible) model homotopy types of spaces. This motivates the modern field of $(\infty,1)$-category theory. We touch on 2-categories in Chapter~\ref{ch:nattrans} but leave the higher theory for other texts.
\end{remark}

\section{Exercises}

\begin{exercise}
Verify that the fundamental groupoid $\Pi_1(X)$ (Example) satisfies the category axioms. Where do you use the axioms of path homotopy?
\end{exercise}

\begin{exercise}
Show that in any category, the identity morphism $\id_a$ is the unique morphism satisfying both unit laws.
\end{exercise}

\begin{exercise}
Let $\cat{C}$ be a category and $f: a \to b$ an isomorphism. Show that $f^{-1}$ is unique.
\end{exercise}

\begin{exercise}
Show that the inclusion $\mathbb{Z} \hookrightarrow \mathbb{Q}$ is an epimorphism in $\Ring$. \emph{Hint:} Show that any ring homomorphism $\mathbb{Q} \to R$ is determined by its values on $\mathbb{Z}$.
\end{exercise}

\begin{exercise}
Describe all initial and terminal objects in the category of fields $\mathbf{Field}$ (with field homomorphisms). What goes wrong?
\end{exercise}

\begin{exercise}
Let $\cat{C}$ be a category with a terminal object $1$. Show that the slice category $\cat{C}/1$ is isomorphic to $\cat{C}$.
\end{exercise}

\begin{exercise}
Prove that if $\cat{C}$ has a zero object $0$, then for any $a,b \in \cat{C}$, the zero morphism $0_{a,b}: a \to 0 \to b$ is the unique morphism factoring through $0$.
\end{exercise}

\begin{exercise}[$\star$]
A \textbf{preorder} is a set $P$ with a relation $\leq$ that is reflexive and transitive (but not necessarily antisymmetric). Show that preorders correspond to categories in which there is at most one morphism between any two objects. How does a poset differ from a preorder, categorically?
\end{exercise}

\begin{exercise}[$\star$]
Define the \textbf{arrow category} $\cat{C}^{\to}$ of a category $\cat{C}$: its objects are the morphisms of $\cat{C}$, and a morphism from $f: a \to b$ to $g: c \to d$ is a commutative square. Show this is indeed a category, and identify it as a functor category $[\mathbf{2}, \cat{C}]$ (see Chapter~\ref{ch:nattrans}).
\end{exercise}

\begin{exercise}[$\star\star$]
Show that the category of relations $\mathbf{Rel}$ (objects: sets; morphisms $A \to B$: subsets $R \subseteq A \times B$; composition: relational composition) is a category. Identify the isomorphisms in $\mathbf{Rel}$. Show there is a functor $\Set \to \mathbf{Rel}$ sending each function to its graph.
\end{exercise}

\chapter{Functors}
\label{ch:functors}

\section{Definition and First Examples}

If categories are the objects of our study, then functors are the morphisms between them. A functor is a map between categories that preserves the categorical structure: it sends objects to objects, morphisms to morphisms, and respects composition and identities.

\begin{definition}[Functor]
Let $\cat{C}$ and $\cat{D}$ be categories. A \textbf{functor} $F: \cat{C} \to \cat{D}$ consists of:
\begin{enumerate}
  \item A function $F: \ob(\cat{C}) \to \ob(\cat{D})$ (we write $Fa$ or $F(a)$).
  \item For each pair $a, b \in \cat{C}$, a function $F_{a,b}: \cat{C}(a,b) \to \cat{D}(Fa, Fb)$ (we write $Ff$ or $F(f)$).
\end{enumerate}
These must satisfy:
\begin{description}
  \item[Preservation of identities.] $F(\id_a) = \id_{Fa}$ for all $a \in \cat{C}$.
  \item[Preservation of composition.] $F(g \circ f) = Fg \circ Ff$ for all composable $f, g$ in $\cat{C}$.
\end{description}
\end{definition}

\begin{remark}
The subscripts on $F_{a,b}$ are usually suppressed; context makes the source and target clear. Note that functors preserve the \emph{domain} and \emph{codomain}: if $f: a \to b$, then $Ff: Fa \to Fb$.
\end{remark}

\begin{example}[Identity functor]
For any category $\cat{C}$, the \textbf{identity functor} $\Id_{\cat{C}}: \cat{C} \to \cat{C}$ sends every object and morphism to itself.
\end{example}

\begin{example}[Forgetful functors]
Many functors ``forget'' structure:
\begin{itemize}
  \item $U: \Grp \to \Set$ sends each group to its underlying set and each homomorphism to the underlying function. This is the prototypical \emph{forgetful functor}.
  \item Similarly, there are forgetful functors $\Ring \to \Grp$ (forgetting multiplicative structure), $\Top \to \Set$ (forgetting topology), $\Vect_k \to \Ab$ (forgetting scalar multiplication).
  \item In general, any time a structure consists of a set (or more structured object) with additional data, there is a forgetful functor to the underlying category.
\end{itemize}
\end{example}

\begin{example}[Free functors]
Free functors are in some sense ``the most efficient way to impose structure.''
\begin{itemize}
  \item $F: \Set \to \Grp$ sends a set $S$ to the free group $F(S)$ on $S$, and a function $f: S \to T$ to the unique group homomorphism $F(f): F(S) \to F(T)$ extending $f$.
  \item $F: \Set \to \Vect_k$ sends $S$ to the free vector space $k^{(S)}$ (formal finite $k$-linear combinations of elements of $S$).
  \item Free functors are always left adjoint to the corresponding forgetful functor (Chapter~\ref{ch:adjoints}).
\end{itemize}
\end{example}

\begin{example}[Power-set functor]
The covariant power-set functor $\mathcal{P}: \Set \to \Set$ sends:
\begin{itemize}
  \item Each set $X$ to its power set $\mathcal{P}(X)$.
  \item Each function $f: X \to Y$ to the direct image function $\mathcal{P}(f): \mathcal{P}(X) \to \mathcal{P}(Y)$, $A \mapsto f(A)$.
\end{itemize}
We will see in Chapter~\ref{ch:monads} that $\mathcal{P}$ carries a monad structure.
\end{example}

\begin{example}[Fundamental group functor]
The assignment $X \mapsto \pi_1(X, x_0)$ (fundamental group at a basepoint) is not quite a functor from $\Top$ to $\Grp$ because basepoints must be preserved. However, $\pi_1: \Top_* \to \Grp$ (pointed spaces and basepoint-preserving continuous maps) is a functor: a continuous map $f: (X, x_0) \to (Y, y_0)$ induces a group homomorphism $f_*: \pi_1(X,x_0) \to \pi_1(Y, y_0)$.
\end{example}

\begin{example}[Homology functors]
For each $n \geq 0$, the $n$-th singular homology group defines a functor $H_n: \Top \to \Ab$. The naturality of the long exact sequence of a pair is the statement that certain natural transformations (Chapter~\ref{ch:nattrans}) exist.
\end{example}

\begin{example}[Constant functor]
For any category $\cat{C}$ and object $d \in \cat{D}$, the \textbf{constant functor} $\Delta_d: \cat{C} \to \cat{D}$ sends every object of $\cat{C}$ to $d$ and every morphism to $\id_d$. We will use this extensively when defining limits and colimits (Chapter~\ref{ch:limits}).
\end{example}

\begin{example}[Diagonal functor]
For a category $\cat{C}$, the \textbf{diagonal functor} $\Delta: \cat{C} \to \cat{C} \times \cat{C}$ sends $c \mapsto (c,c)$ and $f \mapsto (f,f)$.
\end{example}

\section{Contravariant Functors}

\begin{definition}[Contravariant functor]
A \textbf{contravariant functor} $F: \cat{C} \to \cat{D}$ consists of the same data as a functor, except that morphisms are reversed: for $f: a \to b$ in $\cat{C}$, we get $Ff: Fb \to Fa$ in $\cat{D}$. Composition is reversed: $F(g \circ f) = Ff \circ Fg$.
\end{definition}

A contravariant functor $\cat{C} \to \cat{D}$ is equivalently a covariant functor $\cat{C}\op \to \cat{D}$, or a covariant functor $\cat{C} \to \cat{D}\op$. We generally prefer the covariant-functor-on-opposite-category language.

\begin{example}[Contravariant power-set functor]
The contravariant power-set functor $\mathcal{P}: \Set\op \to \Set$ sends $X \mapsto \mathcal{P}(X)$ and $f: X \to Y$ to the preimage function $f^{-1}: \mathcal{P}(Y) \to \mathcal{P}(X)$.
\end{example}

\begin{example}[Dual vector space]
The functor $(-)^*: \Vect_k\op \to \Vect_k$, $V \mapsto V^* = \Hom_k(V, k)$, is a contravariant functor. For a linear map $T: V \to W$, the transpose $T^*: W^* \to V^*$ is defined by $T^*(\phi) = \phi \circ T$.
\end{example}

\begin{example}[Cohomology]
Singular cohomology $H^n: \Top\op \to \Ab$ is a contravariant functor: a continuous map $f: X \to Y$ induces $f^*: H^n(Y) \to H^n(X)$.
\end{example}

\section{Properties of Functors}

\begin{definition}[Full, faithful, essentially surjective]
Let $F: \cat{C} \to \cat{D}$ be a functor.
\begin{itemize}
  \item $F$ is \textbf{full} if each function $F_{a,b}: \cat{C}(a,b) \to \cat{D}(Fa, Fb)$ is surjective.
  \item $F$ is \textbf{faithful} if each function $F_{a,b}$ is injective.
  \item $F$ is \textbf{essentially surjective} (or \textbf{surjective on objects up to isomorphism}) if for every $d \in \cat{D}$, there exists $c \in \cat{C}$ with $Fc \cong d$.
  \item $F$ is a \textbf{fully faithful} if it is both full and faithful.
\end{itemize}
\end{definition}

\begin{remark}
A functor being ``faithful'' does not mean it is injective on objects—two different objects of $\cat{C}$ might map to the same object of $\cat{D}$. Faithfulness only concerns the action on hom-sets.
\end{remark}

\begin{proposition}
A fully faithful functor reflects isomorphisms: if $Ff$ is an isomorphism in $\cat{D}$, then $f$ is an isomorphism in $\cat{C}$.
\end{proposition}

\begin{proof}
If $Ff: Fa \to Fb$ is an isomorphism, let $g' = (Ff)^{-1}: Fb \to Fa$. Since $F$ is full, there exists $g: b \to a$ with $Fg = g'$. Then $F(g \circ f) = Fg \circ Ff = g' \circ Ff = \id_{Fa} = F(\id_a)$. Since $F$ is faithful, $g \circ f = \id_a$. Similarly $f \circ g = \id_b$.
\end{proof}

\begin{definition}[Equivalence of categories]
An \textbf{equivalence of categories} between $\cat{C}$ and $\cat{D}$ is a functor $F: \cat{C} \to \cat{D}$ such that there exist a functor $G: \cat{D} \to \cat{C}$ and natural isomorphisms $GF \cong \Id_{\cat{C}}$ and $FG \cong \Id_{\cat{D}}$. We write $\cat{C} \simeq \cat{D}$ and say $\cat{C}$ and $\cat{D}$ are \textbf{equivalent}.
\end{definition}

\begin{theorem}
A functor $F: \cat{C} \to \cat{D}$ is an equivalence of categories if and only if it is fully faithful and essentially surjective.
\end{theorem}

We defer the proof to after we have defined natural transformations and natural isomorphisms (Chapter~\ref{ch:nattrans}).

\begin{example}
The category of finite-dimensional $k$-vector spaces $\mathbf{FDVect}_k$ is equivalent to the category $\mathbf{Mat}_k$ of matrices over $k$. An equivalence sends each finite-dimensional space to a natural number (its dimension) and uses the fact that every finite-dimensional space is isomorphic to $k^n$. This example shows that equivalent categories can look very different while having the same categorical structure.
\end{example}

\section{The Hom-Functors}
\label{sec:hom-functors}

One of the most important functors associated to a locally small category $\cat{C}$ is the hom-functor, which encodes all morphisms in $\cat{C}$.

\begin{definition}[Covariant hom-functor]
For a fixed object $a \in \cat{C}$, define the functor
\[
  \cat{C}(a, -): \cat{C} \longrightarrow \Set
\]
by:
\begin{itemize}
  \item On objects: $b \mapsto \cat{C}(a, b)$.
  \item On morphisms: for $f: b \to c$, post-composition gives $f_* = \cat{C}(a,f): \cat{C}(a,b) \to \cat{C}(a,c)$, $g \mapsto f \circ g$.
\end{itemize}
\end{definition}

\begin{definition}[Contravariant hom-functor]
For a fixed object $b \in \cat{C}$, define the functor
\[
  \cat{C}(-, b): \cat{C}\op \longrightarrow \Set
\]
by:
\begin{itemize}
  \item On objects: $a \mapsto \cat{C}(a, b)$.
  \item On morphisms: for $f: a \to a'$, pre-composition gives $f^* = \cat{C}(f,b): \cat{C}(a',b) \to \cat{C}(a,b)$, $g \mapsto g \circ f$.
\end{itemize}
\end{definition}

\begin{definition}[Bifunctor hom]
Together, the two variables combine into the \textbf{hom-bifunctor}:
\[
  \cat{C}(-,-): \cat{C}\op \times \cat{C} \longrightarrow \Set,
\]
sending $(a, b) \mapsto \cat{C}(a,b)$ and $(f: a' \to a,\, g: b \to b') \mapsto (h \mapsto g \circ h \circ f)$.
\end{definition}

\section{Representable Functors}
\label{sec:representable}

\begin{definition}[Representable functor]
A functor $F: \cat{C} \to \Set$ is \textbf{representable} if there exists an object $r \in \cat{C}$ and a natural isomorphism $\alpha: \cat{C}(r, -) \xrightarrow{\sim} F$. The pair $(r, \alpha)$ is called a \textbf{representation} of $F$, and $r$ is called the \textbf{representing object}.

Similarly, a contravariant functor $F: \cat{C}\op \to \Set$ is representable if $F \cong \cat{C}(-,r)$ for some $r$.
\end{definition}

\begin{example}[Universal properties as representability]
The notion of a product $a \times b$ in a category $\cat{C}$ can be expressed as representability: the functor
\[
  c \;\mapsto\; \cat{C}(c, a) \times \cat{C}(c, b)
\]
is (contravariantly) representable by $a \times b$ when the latter exists. The representing object $a \times b$ comes with a natural bijection $\cat{C}(c, a\times b) \cong \cat{C}(c,a) \times \cat{C}(c,b)$, which is precisely the universal property of the product.
\end{example}

\begin{example}[Tensor product as representability]
For $R$-modules, the tensor product $M \otimes_R N$ represents the functor $\Hom_R(M, \Hom_R(N, -))$... actually, it represents $\text{Bilinear}(M \times N, -)$. More precisely, the adjunction $M \otimes_R - \dashv \Hom_R(M,-)$ is a statement of representability.
\end{example}

\section{The Yoneda Lemma}
\label{sec:yoneda-lemma}

The Yoneda lemma is arguably the most important theorem in category theory. It says that an object $a$ of a locally small category $\cat{C}$ is completely determined by the functor $\cat{C}(a, -)$ it represents.

\begin{theorem}[Yoneda Lemma]
\label{thm:yoneda}
Let $\cat{C}$ be a locally small category, $a \in \cat{C}$, and $F: \cat{C} \to \Set$ a functor. Then there is a bijection
\[
  \Nat(\cat{C}(a,-),\, F) \;\cong\; Fa,
\]
natural in both $a$ and $F$. Explicitly, the bijection sends a natural transformation $\alpha: \cat{C}(a,-) \Rightarrow F$ to the element $\alpha_a(\id_a) \in Fa$.
\end{theorem}

\begin{proof}
Define $\Phi: \Nat(\cat{C}(a,-), F) \to Fa$ by $\Phi(\alpha) = \alpha_a(\id_a)$.

We construct an inverse $\Psi: Fa \to \Nat(\cat{C}(a,-), F)$. Given $x \in Fa$, define for each $b \in \cat{C}$ the function $\Psi(x)_b: \cat{C}(a,b) \to Fb$ by
\[
  \Psi(x)_b(f) = Ff(x).
\]
We must verify that $\Psi(x)$ is a natural transformation, i.e., that for any $g: b \to c$ in $\cat{C}$, the square
\[
  \begin{tikzcd}
    \cat{C}(a,b) \arrow[r, "\Psi(x)_b"] \arrow[d, "g_*"'] & Fb \arrow[d, "Fg"] \\
    \cat{C}(a,c) \arrow[r, "\Psi(x)_c"'] & Fc
  \end{tikzcd}
\]
commutes. For $f \in \cat{C}(a,b)$: going right-then-down gives $Fg(Ff(x)) = F(g \circ f)(x)$ (by functoriality of $F$); going down-then-right gives $\Psi(x)_c(g \circ f) = F(g \circ f)(x)$. These are equal.

Now we verify $\Phi$ and $\Psi$ are inverse. First, $\Phi(\Psi(x)) = \Psi(x)_a(\id_a) = F(\id_a)(x) = \id_{Fa}(x) = x$, so $\Phi \circ \Psi = \id$.

Second, let $\alpha: \cat{C}(a,-) \Rightarrow F$ be a natural transformation and let $x = \alpha_a(\id_a)$. We claim $\Psi(x) = \alpha$, i.e., for all $b$ and $f: a \to b$, $\Psi(x)_b(f) = \alpha_b(f)$. Indeed, since $\alpha$ is natural, the square
\[
  \begin{tikzcd}
    \cat{C}(a,a) \arrow[r, "\alpha_a"] \arrow[d, "f_*"'] & Fa \arrow[d, "Ff"] \\
    \cat{C}(a,b) \arrow[r, "\alpha_b"'] & Fb
  \end{tikzcd}
\]
commutes. Evaluating at $\id_a$: $\alpha_b(f_*(\id_a)) = Ff(\alpha_a(\id_a))$, i.e., $\alpha_b(f) = Ff(x) = \Psi(x)_b(f)$.

Naturality in $a$ and $F$ is a direct verification we leave as an exercise.
\end{proof}

\begin{corollary}[Yoneda embedding]
\label{cor:yoneda-embedding}
The functor
\[
  \yo: \cat{C} \longrightarrow [\cat{C}\op, \Set], \qquad a \mapsto \cat{C}(-, a)
\]
(where $[\cat{C}\op, \Set]$ is the functor category of presheaves on $\cat{C}$) is fully faithful. It is called the \textbf{Yoneda embedding}.
\end{corollary}

\begin{proof}
By the Yoneda lemma applied to the contravariant version, natural transformations $\cat{C}(-,a) \Rightarrow \cat{C}(-,b)$ are in bijection with $\cat{C}(-,b)$ evaluated at $a$, which is $\cat{C}(a,b)$. Hence $\yo$ induces a bijection $\cat{C}(a,b) \cong [\cat{C}\op,\Set](\yo(a), \yo(b))$, proving full faithfulness.
\end{proof}

The Yoneda embedding shows that every category embeds fully faithfully into a presheaf category, and presheaf categories have very good properties (they are complete and cocomplete, they are toposes, etc.). This is the starting point for many sheaf-theoretic constructions.

\begin{corollary}[Yoneda principle]
Two objects $a, b \in \cat{C}$ are isomorphic if and only if $\cat{C}(-, a) \cong \cat{C}(-, b)$ as functors $\cat{C}\op \to \Set$.
\end{corollary}

\begin{proof}
Immediate from the fact that $\yo$ is fully faithful: an isomorphism in $[\cat{C}\op,\Set]$ between $\yo(a)$ and $\yo(b)$ corresponds to an isomorphism in $\cat{C}$ between $a$ and $b$.
\end{proof}

This is the categorical incarnation of the idea that an object is determined by its ``external behavior''—by how other objects map into it.

\section{Composition of Functors}

\begin{definition}
Given functors $F: \cat{C} \to \cat{D}$ and $G: \cat{D} \to \cat{E}$, their composite $G \circ F: \cat{C} \to \cat{E}$ (also written $GF$) is defined by $(G \circ F)(a) = G(Fa)$ and $(G \circ F)(f) = G(Ff)$.
\end{definition}

Composition of functors is associative and admits identity functors as identities. This makes $\Cat$ into a category (and, as we shall see, a 2-category).

\section{Exercises}

\begin{exercise}
Verify the functor axioms for the covariant power-set functor $\mathcal{P}: \Set \to \Set$.
\end{exercise}

\begin{exercise}
Show that a functor $F: \cat{C} \to \cat{D}$ preserves isomorphisms: if $f: a \to b$ is an isomorphism, then $Ff: Fa \to Fb$ is an isomorphism.
\end{exercise}

\begin{exercise}
Let $F: \cat{C} \to \Set$ be a representable functor with representing object $r$. Show that $r$ is unique up to unique isomorphism.
\end{exercise}

\begin{exercise}
Explicitly describe the natural transformation $\alpha: \cat{C}(r,-) \Rightarrow F$ corresponding (under the Yoneda lemma) to a universal element $u \in Fr$.
\end{exercise}

\begin{exercise}
Prove naturality in $a$ and $F$ for the Yoneda bijection (Theorem~\ref{thm:yoneda}).
\end{exercise}

\begin{exercise}
Show that the assignment $a \mapsto \cat{C}(a,-)$ defines a contravariant functor from $\cat{C}$ to the category $[\cat{C},\Set]$ of functors $\cat{C} \to \Set$. State and prove the corresponding version of the Yoneda embedding.
\end{exercise}

\begin{exercise}[$\star$]
Let $G$ be a group, viewed as a one-object category $\mathbf{B}G$. Describe the functor category $[\mathbf{B}G, \Set]$. What are its objects and morphisms? (These are called \emph{$G$-sets}.)
\end{exercise}

\begin{exercise}[$\star$]
Show that any equivalence of categories is fully faithful and essentially surjective. (The other direction requires the axiom of choice and naturality arguments—prove it after reading Chapter~\ref{ch:nattrans}.)
\end{exercise}

\begin{exercise}[$\star\star$]
The \textbf{Cayley theorem} for groups states that every group embeds in a symmetric group. Prove the categorical analogue: every small category $\cat{C}$ embeds fully faithfully in $[\cat{C}\op, \Set]$ (the Yoneda embedding). Discuss in what sense this generalizes Cayley's theorem (take $\cat{C} = \mathbf{B}G$).
\end{exercise}

\chapter{Natural Transformations}
\label{ch:nattrans}

\section{Definition}

Natural transformations are morphisms between functors. They provide the precise meaning of ``a natural map'' between constructions—a notion that Eilenberg and Mac Lane introduced category theory specifically to capture.

\begin{definition}[Natural transformation]
Let $F, G: \cat{C} \to \cat{D}$ be functors. A \textbf{natural transformation} $\alpha: F \Rightarrow G$ consists of:
\begin{itemize}
  \item For each object $a \in \cat{C}$, a morphism $\alpha_a: Fa \to Ga$ in $\cat{D}$, called the \textbf{component} of $\alpha$ at $a$.
\end{itemize}
These components must satisfy the \textbf{naturality condition}: for every morphism $f: a \to b$ in $\cat{C}$, the following square commutes in $\cat{D}$:
\[
  \begin{tikzcd}
    Fa \arrow[r, "\alpha_a"] \arrow[d, "Ff"'] & Ga \arrow[d, "Gf"] \\
    Fb \arrow[r, "\alpha_b"'] & Gb
  \end{tikzcd}
\]
That is, $Gf \circ \alpha_a = \alpha_b \circ Ff$.
\end{definition}

\begin{remark}
The naturality condition says that the components of $\alpha$ ``commute with the action of morphisms.'' A natural transformation does not just require a map $Fa \to Ga$ for each $a$; it requires this collection of maps to be coherent with the functorial structure.
\end{remark}

\begin{example}[Double dual]
Let $F = \Id_{\Vect_k}$ and $G = (-)^{**}$ (double dual). For each vector space $V$, define $\alpha_V: V \to V^{**}$ by $\alpha_V(v)(\phi) = \phi(v)$ (evaluation). This is a natural transformation $\Id \Rightarrow (-)^{**}$: for any linear map $T: V \to W$, the diagram
\[
  \begin{tikzcd}
    V \arrow[r, "\alpha_V"] \arrow[d, "T"'] & V^{**} \arrow[d, "T^{**}"] \\
    W \arrow[r, "\alpha_W"'] & W^{**}
  \end{tikzcd}
\]
commutes. (When $\dim V < \infty$, $\alpha_V$ is an isomorphism, but the \emph{single} dual $V \to V^*$ cannot be made natural with respect to all linear maps—it requires a choice of basis.)
\end{example}

\begin{example}[Determinant]
The determinant $\det: \mathbf{GL}_n \Rightarrow (-)^\times$ is a natural transformation from the functor sending a commutative ring $R$ to $\mathbf{GL}_n(R)$ (invertible $n \times n$ matrices) to the functor sending $R$ to its unit group $R^\times$.
\end{example}

\begin{example}[Abelianization]
Let $\mathrm{ab}: \Grp \to \Ab$ be the abelianization functor $G \mapsto G/[G,G]$. There is a natural transformation $\eta: \Id_{\Grp} \Rightarrow i \circ \mathrm{ab}$ (where $i: \Ab \hookrightarrow \Grp$ is inclusion), given by the quotient maps $\eta_G: G \to G/[G,G]$.
\end{example}

\begin{definition}[Natural isomorphism]
A natural transformation $\alpha: F \Rightarrow G$ is a \textbf{natural isomorphism} (written $\alpha: F \xrightarrow{\sim} G$ or $F \cong G$) if every component $\alpha_a: Fa \to Ga$ is an isomorphism in $\cat{D}$.
\end{definition}

\begin{lemma}
$\alpha: F \Rightarrow G$ is a natural isomorphism if and only if $\alpha$ is an isomorphism in the functor category $[\cat{C}, \cat{D}]$ (defined below).
\end{lemma}

\section{The Functor Category}

\begin{definition}[Functor category]
Let $\cat{C}$ and $\cat{D}$ be categories. The \textbf{functor category} $[\cat{C}, \cat{D}]$ (also written $\cat{D}^{\cat{C}}$ or $\Fun(\cat{C},\cat{D})$) has:
\begin{itemize}
  \item Objects: functors $F: \cat{C} \to \cat{D}$.
  \item Morphisms: natural transformations $\alpha: F \Rightarrow G$.
  \item Composition: vertical composition (defined below).
  \item Identities: the identity natural transformation $\id_F$, with components $(\id_F)_a = \id_{Fa}$.
\end{itemize}
\end{definition}

\begin{remark}
When $\cat{C}$ is small and $\cat{D}$ is locally small, the functor category $[\cat{C},\cat{D}]$ is locally small. If $\cat{C}$ is not small, there may be set-theoretic issues.
\end{remark}

\begin{example}[Presheaf category]
For a small category $\cat{C}$, the \textbf{presheaf category} $\widehat{\cat{C}} = [\cat{C}\op, \Set]$ is a functor category of fundamental importance. Its objects are contravariant functors $\cat{C}\op \to \Set$ (called \emph{presheaves} on $\cat{C}$), and its morphisms are natural transformations between them. The Yoneda embedding (Corollary~\ref{cor:yoneda-embedding}) embeds $\cat{C}$ fully faithfully into $\widehat{\cat{C}}$.
\end{example}

\section{Vertical Composition}

\begin{definition}[Vertical composition]
Given natural transformations $\alpha: F \Rightarrow G$ and $\beta: G \Rightarrow H$ (all between functors $\cat{C} \to \cat{D}$), their \textbf{vertical composite} $\beta \circ \alpha: F \Rightarrow H$ is defined componentwise:
\[
  (\beta \circ \alpha)_a = \beta_a \circ \alpha_a: Fa \to Ga \to Ha.
\]
\end{definition}

\begin{proposition}
Vertical composition is associative and unital. Hence $[\cat{C},\cat{D}]$ is indeed a category.
\end{proposition}

\begin{proof}
Associativity of vertical composition follows from associativity of composition in $\cat{D}$. The identity natural transformation $\id_F$ (with components $\id_{Fa}$) is the identity for vertical composition. Naturality of the composite is a routine diagram chase.
\end{proof}

\section{Horizontal Composition and Whiskering}

Natural transformations can also be composed in a horizontal direction, combining transformations of adjacent functors.

\begin{definition}[Horizontal composition]
Let $\alpha: F \Rightarrow G$ be a natural transformation between functors $\cat{C} \to \cat{D}$, and $\beta: H \Rightarrow K$ a natural transformation between functors $\cat{D} \to \cat{E}$. The \textbf{horizontal composite} $\beta * \alpha: HF \Rightarrow KG$ has components
\[
  (\beta * \alpha)_a = \beta_{Ga} \circ H(\alpha_a) = K(\alpha_a) \circ \beta_{Fa}.
\]
(Both expressions are equal by naturality of $\beta$.)
\end{definition}

The two-step construction can be visualized as:
\[
  \begin{tikzcd}
    \cat{C} \ar[r, bend left=50, "F", ""{name=F}]
            \ar[r, bend right=50, "G"', ""{name=G}] &
    \cat{D} \ar[r, bend left=50, "H", ""{name=H}]
            \ar[r, bend right=50, "K"', ""{name=K}] &
    \cat{E}
    \ar[Rightarrow, from=F, to=G, "\alpha"]
    \ar[Rightarrow, from=H, to=K, "\beta"]
  \end{tikzcd}
\]

\begin{definition}[Whiskering]
Special cases of horizontal composition:
\begin{itemize}
  \item \textbf{Left whiskering}: for a functor $L: \cat{B} \to \cat{C}$ and a natural transformation $\alpha: F \Rightarrow G$ (between $\cat{C} \to \cat{D}$), define $\alpha L: FL \Rightarrow GL$ with $(\alpha L)_b = \alpha_{Lb}$.
  \item \textbf{Right whiskering}: for $\alpha: F \Rightarrow G$ and a functor $R: \cat{D} \to \cat{E}$, define $R\alpha: RF \Rightarrow RG$ with $(R\alpha)_a = R(\alpha_a)$.
\end{itemize}
\end{definition}

\section{The Interchange Law}

Vertical and horizontal composition interact via the interchange law.

\begin{theorem}[Interchange law]
Let
\[
  \begin{tikzcd}
    \cat{C} \ar[r, bend left=50, "F_1", ""{name=F1}]
            \ar[r, "F_2"', ""{name=F2, above}]
            \ar[r, bend right=50, "F_3"', ""{name=F3}] &
    \cat{D} \ar[r, bend left=50, "G_1", ""{name=G1}]
            \ar[r, "G_2"', ""{name=G2, above}]
            \ar[r, bend right=50, "G_3"', ""{name=G3}] &
    \cat{E}
    \ar[Rightarrow, from=F1, to=F2, "\alpha"]
    \ar[Rightarrow, from=F2, to=F3, "\beta"]
    \ar[Rightarrow, from=G1, to=G2, "\gamma"]
    \ar[Rightarrow, from=G2, to=G3, "\delta"]
  \end{tikzcd}
\]
Then $(\delta \circ \gamma) * (\beta \circ \alpha) = (\delta * \beta) \circ (\gamma * \alpha)$.
\end{theorem}

\begin{proof}
At each object $a \in \cat{C}$, both sides equal $\delta_{F_3 a} \circ G_2(\beta_a) \circ \gamma_{F_2 a} \circ G_1(\alpha_a)$, using naturality of $\delta$ and $\gamma$ at appropriate steps.
\end{proof}

\section{The 2-Category \texorpdfstring{$\Cat$}{Cat}}

The interchange law says that $\Cat$ (or more precisely, $\mathbf{CAT}$, the 2-category of all categories) has the structure of a \textbf{2-category}.

\begin{definition}[2-Category]
A \textbf{2-category} $\mathcal{K}$ consists of:
\begin{itemize}
  \item \textbf{0-cells} (objects): $\cat{C}$, $\cat{D}$, \ldots
  \item \textbf{1-cells} (morphisms): functors $F: \cat{C} \to \cat{D}$.
  \item \textbf{2-cells} (morphisms between morphisms): natural transformations $\alpha: F \Rightarrow G$.
  \item \textbf{Vertical composition} of 2-cells (along shared 1-cells).
  \item \textbf{Horizontal composition} of 2-cells (along shared 0-cells).
  \item The interchange law.
  \item Strict associativity and unit laws at all levels.
\end{itemize}
\end{definition}

$\Cat$ is the archetypal 2-category. Its 0-cells are categories, its 1-cells are functors, and its 2-cells are natural transformations. We will use 2-categorical language sparingly but it provides a useful meta-framework.

\section{Equivalences of Categories Revisited}

We can now make the notion of equivalence precise using natural isomorphisms.

\begin{definition}[Equivalence of categories, revisited]
An \textbf{equivalence of categories} consists of functors $F: \cat{C} \to \cat{D}$ and $G: \cat{D} \to \cat{C}$, together with natural isomorphisms $\eta: \Id_{\cat{C}} \xrightarrow{\sim} GF$ and $\varepsilon: FG \xrightarrow{\sim} \Id_{\cat{D}}$.
\end{definition}

\begin{theorem}
A functor $F: \cat{C} \to \cat{D}$ is part of an equivalence of categories if and only if it is fully faithful and essentially surjective.
\end{theorem}

\begin{proof}[Proof sketch]
($\Rightarrow$) Suppose $F$ is an equivalence via $(F, G, \eta, \varepsilon)$. Full faithfulness: the bijection $\cat{C}(a,b) \cong \cat{D}(Fa, Fb)$ is realized as $f \mapsto Ff$; use $\eta$ and $\varepsilon$ to construct its inverse. Essential surjectivity: for $d \in \cat{D}$, $\varepsilon_d: FG(d) \xrightarrow{\sim} d$ is an isomorphism, so $d \cong F(Gd)$.

($\Leftarrow$) Given $F$ fully faithful and essentially surjective, use essential surjectivity (and the axiom of choice) to select for each $d \in \cat{D}$ an object $Gd$ and an isomorphism $\varepsilon_d: FGd \xrightarrow{\sim} d$. Define $G$ on morphisms using full faithfulness of $F$. Verify $\eta$ and $\varepsilon$ are natural isomorphisms.
\end{proof}

\begin{example}
The skeletal equivalence: every category is equivalent to its \textbf{skeleton} (the full subcategory containing exactly one representative from each isomorphism class of objects).
\end{example}

\section{The Yoneda Lemma, Revisited}

The full strength of the Yoneda lemma is a naturality statement about functor categories.

\begin{theorem}[Yoneda lemma, natural version]
For a locally small category $\cat{C}$, the bijection of Theorem~\ref{thm:yoneda}
\[
  \Nat(\cat{C}(a,-), F) \;\cong\; Fa
\]
is natural in $a \in \cat{C}\op$ and $F \in [\cat{C}, \Set]$. Equivalently, the evaluation functor
\[
  \mathrm{ev}: [\cat{C}, \Set] \times \cat{C} \longrightarrow \Set, \quad (F, a) \mapsto Fa
\]
is naturally isomorphic to the functor $(F, a) \mapsto \Nat(\cat{C}(a,-), F)$.
\end{theorem}

This version makes the Yoneda lemma a statement about a natural isomorphism of bifunctors, and it is the form most useful in practice.

\section{Limits as Representability}
\label{sec:limits-as-rep}

We preview how limits and colimits (Chapter~\ref{ch:limits}) are naturally expressed in terms of representability and natural transformations.

A \textbf{cone} over a diagram $D: \cat{J} \to \cat{C}$ with apex $c$ is a natural transformation $\Delta_c \Rightarrow D$, where $\Delta_c: \cat{J} \to \cat{C}$ is the constant functor at $c$. The \textbf{limit} of $D$, if it exists, is the representing object for the functor
\[
  c \;\mapsto\; \Nat(\Delta_c, D) \;=\; \text{the set of cones over } D \text{ with apex } c.
\]
This makes the definition of limit a clean statement in terms of representable functors.

\section{Exercises}

\begin{exercise}
Verify that vertical composition of natural transformations satisfies the unit and associativity axioms.
\end{exercise}

\begin{exercise}
Show that a natural transformation $\alpha: F \Rightarrow G$ is a natural isomorphism if and only if it has an inverse natural transformation $\alpha^{-1}: G \Rightarrow F$ (with $(\alpha^{-1})_a = (\alpha_a)^{-1}$).
\end{exercise}

\begin{exercise}
Let $F, G: \cat{C} \to \cat{D}$ be naturally isomorphic. Show that for any $H: \cat{D} \to \cat{E}$, we have $HF \cong HG$. Dually, for any $K: \cat{B} \to \cat{C}$, $FK \cong GK$.
\end{exercise}

\begin{exercise}
Show that the functor category $[\mathbf{1}, \cat{D}]$ (where $\mathbf{1}$ is the terminal category) is isomorphic to $\cat{D}$.
\end{exercise}

\begin{exercise}
Show that the functor category $[\mathbf{2}, \cat{C}]$ (where $\mathbf{2}$ is the category $0 \to 1$) is isomorphic to the arrow category $\cat{C}^{\to}$.
\end{exercise}

\begin{exercise}
Prove the interchange law in full detail.
\end{exercise}

\begin{exercise}[$\star$]
Let $\cat{C}$ be a category and $\widehat{\cat{C}} = [\cat{C}\op, \Set]$ the presheaf category. Show that the Yoneda embedding $\yo: \cat{C} \to \widehat{\cat{C}}$ sends limits in $\cat{C}$ to limits in $\widehat{\cat{C}}$. (This requires knowing what limits are; revisit after Chapter~\ref{ch:limits}.)
\end{exercise}

\begin{exercise}[$\star$]
A functor $F: \cat{C} \to \cat{D}$ is an equivalence of categories if and only if there exists a functor $G$ with $GF \cong \Id_{\cat{C}}$ and $FG \cong \Id_{\cat{D}}$. Show that in this case, the same $G$ witnesses $F$ as fully faithful and essentially surjective.
\end{exercise}

\begin{exercise}[$\star\star$]
Prove that for a small category $\cat{C}$, the presheaf category $[\cat{C}\op, \Set]$ is the \textbf{free cocompletion} of $\cat{C}$: every functor $F: \cat{C} \to \cat{D}$ to a cocomplete category $\cat{D}$ extends uniquely (up to natural isomorphism) to a colimit-preserving functor $\bar{F}: [\cat{C}\op,\Set] \to \cat{D}$.
\end{exercise}

\chapter{Limits and Colimits}
\label{ch:limits}

\section{Motivation}

Many standard constructions in mathematics—products, intersections, preimages, completions, and more—share a common formal structure: they are defined by a universal property. Category theory provides a unified framework for all such constructions through the concept of \emph{limit} and its dual, \emph{colimit}.

Informally, a limit of a diagram is the ``best approximation from outside'' to the diagram, while a colimit is the ``best approximation from inside.'' We will make these ideas precise using the language of cones and universal properties.

\section{Diagrams}

\begin{definition}[Diagram]
A \textbf{diagram} of shape $\cat{J}$ in a category $\cat{C}$ is a functor $D: \cat{J} \to \cat{C}$. The category $\cat{J}$ is called the \textbf{index category} (or \textbf{shape}) of the diagram. We often refer to the diagram by specifying the objects $D(j)$ for $j \in \cat{J}$ and the morphisms $D(f): D(j) \to D(k)$ for $f: j \to k$ in $\cat{J}$.
\end{definition}

\begin{example}[Discrete diagram]
When $\cat{J}$ is a discrete category on objects $\{1, \ldots, n\}$ (only identity morphisms), a diagram $D: \cat{J} \to \cat{C}$ is simply a collection of $n$ objects $D_1, \ldots, D_n$ in $\cat{C}$. The limit of such a diagram is a product.
\end{example}

\begin{example}[Parallel pair diagram]
Let $\cat{J}$ be the category with two objects $0, 1$ and two morphisms $f, g: 0 \to 1$ (plus identities). A diagram $D: \cat{J} \to \cat{C}$ picks out two objects $A = D(0)$, $B = D(1)$, and two morphisms $Df, Dg: A \to B$. The limit of this diagram is an equalizer.
\end{example}

\section{Cones}

\begin{definition}[Cone]
Let $D: \cat{J} \to \cat{C}$ be a diagram. A \textbf{cone} over $D$ with \textbf{apex} $c \in \cat{C}$ is a natural transformation $\lambda: \Delta_c \Rightarrow D$, where $\Delta_c: \cat{J} \to \cat{C}$ is the constant functor at $c$.

Explicitly, a cone consists of morphisms $\lambda_j: c \to D(j)$ for each $j \in \cat{J}$, such that for every morphism $f: j \to k$ in $\cat{J}$:
\[
  D(f) \circ \lambda_j = \lambda_k.
\]
\end{definition}

\begin{definition}[Cocone]
A \textbf{cocone} under $D$ with \textbf{nadir} $c$ is a natural transformation $\lambda: D \Rightarrow \Delta_c$, consisting of morphisms $\lambda_j: D(j) \to c$ such that $\lambda_k \circ D(f) = \lambda_j$ for all $f: j \to k$.
\end{definition}

\section{Limits}

\begin{definition}[Limit]
Let $D: \cat{J} \to \cat{C}$ be a diagram. A \textbf{limit} of $D$ is a cone $(\lim D,\, \pi)$ (where $\pi_j: \lim D \to D(j)$) such that for every cone $(c, \lambda)$ over $D$, there is a \emph{unique} morphism $u: c \to \lim D$ in $\cat{C}$ making all triangles commute:
\[
  \pi_j \circ u = \lambda_j \qquad \text{for all } j \in \cat{J}.
\]
Equivalently, $\lim D$ represents the functor $\mathrm{Cone}(-, D): \cat{C}\op \to \Set$, where $\mathrm{Cone}(c,D) = \Nat(\Delta_c, D)$.
\end{definition}

\begin{proposition}
If a limit of $D$ exists, it is unique up to unique isomorphism.
\end{proposition}

\begin{proof}
Suppose $(L, \pi)$ and $(L', \pi')$ are both limits. By universality of $(L, \pi)$ applied to $(L', \pi')$, there is a unique $u: L' \to L$ with $\pi_j \circ u = \pi'_j$. By universality of $(L', \pi')$ applied to $(L, \pi)$, there is a unique $v: L \to L'$ with $\pi'_j \circ v = \pi_j$. Then $u \circ v: L \to L$ satisfies $\pi_j \circ (u \circ v) = \pi_j = \pi_j \circ \id_L$; by uniqueness, $u \circ v = \id_L$. Similarly $v \circ u = \id_{L'}$.
\end{proof}

\section{Key Examples of Limits}

\subsection{Products}

\begin{definition}[Product]
The \textbf{product} of objects $\{A_i\}_{i \in I}$ in $\cat{C}$ is the limit of the discrete diagram $D: I \to \cat{C}$, $i \mapsto A_i$. It is written $\prod_{i \in I} A_i$, with \textbf{projection morphisms} $\pi_i: \prod_j A_j \to A_i$.

The universal property: for any $c$ and morphisms $f_i: c \to A_i$, there is a unique morphism $\langle f_i \rangle: c \to \prod_j A_j$ with $\pi_i \circ \langle f_i \rangle = f_i$.

For two objects: $A \times B$ is the binary product.
\end{definition}

\begin{example}
In $\Set$, the product is the Cartesian product. In $\Grp$, it is the direct product. In $\Top$, it is the product topology. In a poset, the product of $a$ and $b$ is $\inf(a,b)$ (the greatest lower bound).
\end{example}

\subsection{Equalizers}

\begin{definition}[Equalizer]
The \textbf{equalizer} of two morphisms $f, g: A \to B$ is the limit of the parallel pair diagram. It consists of an object $E$ and a morphism $e: E \to A$ such that:
\begin{itemize}
  \item $f \circ e = g \circ e$.
  \item Universal property: for any $h: X \to A$ with $f \circ h = g \circ h$, there is a unique $u: X \to E$ with $e \circ u = h$.
\end{itemize}
\end{definition}

\begin{example}
In $\Set$, the equalizer of $f, g: A \to B$ is $\{a \in A \mid f(a) = g(a)\} \hookrightarrow A$. In $\Grp$, it is the kernel of $f \cdot g^{-1}$ (preimage of the identity under $f \cdot g^{-1}: A \to B$). In $\Top$, it carries the subspace topology.
\end{example}

\begin{proposition}
Equalizers are monomorphisms.
\end{proposition}

\begin{proof}
Let $e: E \to A$ be the equalizer of $f, g$. Suppose $e \circ u = e \circ v$ for $u, v: X \to E$. Then $f \circ (e \circ u) = g \circ (e \circ u)$, so by uniqueness in the universal property, $u = v$.
\end{proof}

\subsection{Pullbacks}

\begin{definition}[Pullback]
Let $f: A \to C$ and $g: B \to C$. The \textbf{pullback} (or \textbf{fibered product}) $A \times_C B$ is the limit of the cospan diagram:
\[
  \begin{tikzcd}
    & A \arrow[d, "f"] \\
    B \arrow[r, "g"'] & C
  \end{tikzcd}
\]
It comes with morphisms $p_A: A \times_C B \to A$ and $p_B: A \times_C B \to B$ such that $f \circ p_A = g \circ p_B$, and is universal with this property.
\end{definition}

\begin{example}
In $\Set$, $A \times_C B = \{(a,b) \in A \times B \mid f(a) = g(b)\}$. Pullbacks in $\Top$ carry the subspace topology of $A \times B$. In $\Grp$, the pullback of $f: H \to G$ and $g: K \to G$ is $\{(h,k) \in H \times K \mid f(h) = g(k)\}$.
\end{example}

\begin{proposition}
Pullback of a monomorphism is a monomorphism.
\end{proposition}

\begin{proof}
Let $m: A \to C$ be mono, and form the pullback $P = B \times_C A$ with projections $p: P \to B$, $q: P \to A$. Suppose $p \circ u = p \circ v$ for $u, v: X \to P$. Then $q \circ u$ and $q \circ v$ both satisfy $m \circ (q \circ u) = g \circ p \circ u = g \circ p \circ v = m \circ (q \circ v)$. Since $m$ is mono, $q \circ u = q \circ v$. Now $p \circ u = p \circ v$ and $q \circ u = q \circ v$ imply $u = v$ by the universal property of the pullback.
\end{proof}

\section{Colimits}

Colimits are the dual notion, obtained by reversing all arrows.

\begin{definition}[Colimit]
The \textbf{colimit} of $D: \cat{J} \to \cat{C}$ is a cocone $(\colim D, \iota)$ (with $\iota_j: D(j) \to \colim D$) that is universal: for every cocone $(c, \lambda)$, there is a unique $u: \colim D \to c$ with $u \circ \iota_j = \lambda_j$.
\end{definition}

\subsection{Coproducts}

\begin{definition}[Coproduct]
The \textbf{coproduct} $\coprod_{i \in I} A_i$ is the colimit of the discrete diagram $i \mapsto A_i$, with \textbf{inclusion morphisms} $\iota_i: A_i \to \coprod_j A_j$.
\end{definition}

\begin{example}
In $\Set$, the coproduct is the disjoint union. In $\Grp$, it is the free product. In $\Ab$ and $\Vect_k$, it is the direct sum (also the product for finite index sets). In $\Top$, it is the disjoint union with the coproduct topology. In a poset, coproduct is the join (least upper bound).
\end{example}

\subsection{Coequalizers}

\begin{definition}[Coequalizer]
The \textbf{coequalizer} of $f, g: A \to B$ is an object $Q$ with a morphism $q: B \to Q$ such that $q \circ f = q \circ g$, universal with this property.
\end{definition}

\begin{example}
In $\Set$, the coequalizer of $f, g: A \to B$ is the quotient $B / {\sim}$, where $\sim$ is the equivalence relation generated by $f(a) \sim g(a)$ for all $a \in A$.
\end{example}

\begin{proposition}
Coequalizers are epimorphisms.
\end{proposition}

\subsection{Pushouts}

\begin{definition}[Pushout]
The \textbf{pushout} of $f: C \to A$ and $g: C \to B$ (a span diagram) is the colimit, written $A \sqcup_C B$, with inclusions $i_A: A \to A \sqcup_C B$ and $i_B: B \to A \sqcup_C B$ such that $i_A \circ f = i_B \circ g$.
\end{definition}

\begin{example}
In $\Set$, the pushout is $(A \sqcup B) / \sim$ where $f(c) \sim g(c)$ for all $c \in C$. In $\Top$, pushouts are used to define CW-complex attachments. In $\Grp$, the pushout is the amalgamated free product $A *_C B$. In $\mathbf{CommRing}$, the pushout is the tensor product $A \otimes_C B$.
\end{example}

\section{Existence of Limits}

\begin{definition}[Complete and cocomplete categories]
A category is \textbf{complete} if it has all small limits (i.e., limits of all diagrams indexed by small categories). It is \textbf{cocomplete} if it has all small colimits.
\end{definition}

\begin{theorem}[Construction from products and equalizers]
A category has all finite limits if and only if it has all finite products and all equalizers.
\end{theorem}

\begin{proof}
($\Rightarrow$) Products and equalizers are special limits.

($\Leftarrow$) Given a finite diagram $D: \cat{J} \to \cat{C}$, we construct its limit as an equalizer of two maps between products. Specifically, form:
\[
  \prod_{j \in \cat{J}} D(j) \;\underset{g}{\overset{f}{\rightrightarrows}}\; \prod_{f: j \to k \text{ in } \cat{J}} D(k)
\]
where $f$ is defined using the component maps $D(f)$ and $g$ uses the projections. The equalizer of $f$ and $g$ is the limit of $D$.

More explicitly, let $E = \text{Eq}(s,t)$ where:
\[
  s\left((x_j)\right)_{(f: j \to k)} = D(f)(x_j), \quad t\left((x_j)\right)_{(f:j\to k)} = x_k.
\]
Then $E \cong \lim D$.
\end{proof}

\begin{corollary}
A category has all small limits if and only if it has all small products and all equalizers.
\end{corollary}

\begin{theorem}
The categories $\Set$, $\Grp$, $\Ab$, $\Vect_k$, $\Ring$, $\Top$ are complete and cocomplete.
\end{theorem}

\section{Limits and Colimits in Functor Categories}

\begin{theorem}
Let $\cat{C}$ be a small category and $\cat{D}$ a complete (resp.\ cocomplete) category. Then the functor category $[\cat{C}, \cat{D}]$ is complete (resp.\ cocomplete), with limits and colimits computed pointwise: $(\lim F)_c = \lim(F_c)$ for each $c \in \cat{C}$.
\end{theorem}

\begin{corollary}
The presheaf category $[\cat{C}\op, \Set]$ is complete and cocomplete for any small $\cat{C}$.
\end{corollary}

\section{Filtered Colimits and Flat Functors}

\begin{definition}[Filtered category]
A category $\cat{J}$ is \textbf{filtered} if:
\begin{enumerate}
  \item It is nonempty.
  \item For any two objects $j, k \in \cat{J}$, there exist an object $l$ and morphisms $j \to l$ and $k \to l$.
  \item For any two parallel morphisms $f, g: j \to k$, there exists a morphism $h: k \to l$ with $h \circ f = h \circ g$.
\end{enumerate}
A \textbf{filtered colimit} is the colimit of a diagram $D: \cat{J} \to \cat{C}$ with $\cat{J}$ filtered.
\end{definition}

\begin{example}
Directed systems (in algebra) index filtered colimits. The colimit is the \emph{direct limit}. For example, the field of fractions $k(x)$ is the filtered colimit of polynomial rings $k[x, f^{-1}]$ as $f$ ranges over nonzero polynomials.
\end{example}

\begin{theorem}[Filtered colimits commute with finite limits in $\Set$]
In $\Set$, filtered colimits commute with finite limits: for a filtered diagram $D: \cat{J} \to [\cat{K}, \Set]$ (with $\cat{K}$ finite),
\[
  \colim_{\cat{J}} \lim_{\cat{K}} D \;\cong\; \lim_{\cat{K}} \colim_{\cat{J}} D.
\]
\end{theorem}

\section{Preservation, Reflection, and Creation}

\begin{definition}
Let $F: \cat{C} \to \cat{D}$ be a functor and $D: \cat{J} \to \cat{C}$ a diagram.
\begin{itemize}
  \item $F$ \textbf{preserves} limits of shape $\cat{J}$ if: whenever $(\ell, \pi)$ is a limit of $D$, $(F\ell, F\pi)$ is a limit of $FD$.
  \item $F$ \textbf{reflects} limits of shape $\cat{J}$ if: whenever $(c, \lambda)$ is a cone and $(Fc, F\lambda)$ is a limit of $FD$, then $(c, \lambda)$ is a limit of $D$.
  \item $F$ \textbf{creates} limits of shape $\cat{J}$ if: for every limit $(d, \mu)$ of $FD$, there exists a unique cone $(c, \lambda)$ over $D$ with $Fc = d$ and $F\lambda = \mu$, and this cone is a limit.
\end{itemize}
\end{definition}

The most important case is that right adjoints preserve limits and left adjoints preserve colimits—this is the content of the next chapter.

\begin{theorem}[Representables preserve limits]
For any $a \in \cat{C}$, the covariant hom-functor $\cat{C}(a,-): \cat{C} \to \Set$ preserves all limits.
\end{theorem}

\begin{proof}
A cone $(c, \lambda)$ over $D: \cat{J} \to \cat{C}$ corresponds, by the universal property of $\lim D$, to a unique morphism $c \to \lim D$. Applying $\cat{C}(a,-)$: a natural transformation from $\Delta_{\cat{C}(a,a')}$ to $\cat{C}(a, D(-))$ consists of maps $\cat{C}(a, D(j))$ forming a cone, corresponding to a map $a \to \lim D$, i.e., an element of $\cat{C}(a, \lim D)$. This bijection is natural, proving $\cat{C}(a, \lim D) \cong \lim \cat{C}(a, D(-))$.
\end{proof}

\section{Exercises}

\begin{exercise}
Show that the pullback of $f: A \to C$ and $g: B \to C$ in $\Set$ is $\{(a,b) \in A \times B \mid f(a) = g(b)\}$.
\end{exercise}

\begin{exercise}
Show that terminal objects are limits of empty diagrams ($\cat{J} = \emptyset$), and initial objects are colimits of empty diagrams.
\end{exercise}

\begin{exercise}
Prove that the colimit of a diagram $D: \cat{J} \to \cat{C}$ can be computed as the coequalizer of two maps between coproducts, dualizing the construction in the proof of the finite-limits theorem.
\end{exercise}

\begin{exercise}
Show that in $\Ab$, finite products and finite coproducts coincide (both are given by the direct sum). Show this fails for infinite index sets.
\end{exercise}

\begin{exercise}
In $\Top$, describe the pushout of two inclusions $A \hookleftarrow C \hookrightarrow B$. How does this relate to the construction of CW complexes?
\end{exercise}

\begin{exercise}[$\star$]
Prove that if $F: \cat{C} \to \cat{D}$ and $G: \cat{D} \to \cat{E}$ both preserve limits of shape $\cat{J}$, then $G \circ F$ also preserves limits of shape $\cat{J}$.
\end{exercise}

\begin{exercise}[$\star$]
Show that limits in $[\cat{C}, \cat{D}]$ are computed pointwise: $(\lim_i F_i)_c \cong \lim_i (F_i(c))$, when $\cat{D}$ has the required limits. Dualize for colimits.
\end{exercise}

\begin{exercise}[$\star$]
A functor $F: \cat{C} \to \cat{D}$ is called \textbf{continuous} if it preserves all small limits, and \textbf{cocontinuous} if it preserves all small colimits. Show that representable functors $\cat{C}(a,-)$ are continuous, and corepresentable functors $\cat{C}(-,b)$ are contravariant-continuous (send colimits to limits).
\end{exercise}

\begin{exercise}[$\star\star$]
Prove that filtered colimits commute with finite limits in $\Set$.
\end{exercise}

\chapter{Adjoint Functors}
\label{ch:adjoints}

\section{Motivation}

Adjoint functors appear everywhere in mathematics. Free groups, tensor products, sheafification, compactification, localization—all are expressed most cleanly as adjoints. Mac Lane wrote that adjoint functors ``arise everywhere.'' Understanding them is perhaps the single most important goal of an introductory category theory course.

The idea is simple: two functors $F: \cat{C} \to \cat{D}$ and $G: \cat{D} \to \cat{C}$ are \emph{adjoint} ($F$ left adjoint to $G$, $G$ right adjoint to $F$) if there is a natural bijection between morphisms $Fc \to d$ and morphisms $c \to Gd$. This says that ``maps out of $Fc$'' correspond naturally to ``maps into $Gd$''—the left adjoint $F$ is the ``free'' or ``most efficient'' way to get from $\cat{C}$ to $\cat{D}$, and $G$ is the corresponding ``forgetful'' functor.

\section{The Hom-Set Definition}

\begin{definition}[Adjunction, hom-set version]
Let $F: \cat{C} \to \cat{D}$ and $G: \cat{D} \to \cat{C}$ be functors. An \textbf{adjunction} $F \dashv G$ (read: ``$F$ is left adjoint to $G$,'' or ``$G$ is right adjoint to $F$'') is a natural bijection
\[
  \phi_{c,d}:\; \cat{D}(Fc, d) \;\xrightarrow{\;\sim\;}\; \cat{C}(c, Gd),
\]
natural in $c \in \cat{C}$ and $d \in \cat{D}$.

Naturality in $c$: for $h: c' \to c$ and $f: Fc \to d$, $\phi_{c',d}(f \circ Fh) = \phi_{c,d}(f) \circ h$.
Naturality in $d$: for $k: d \to d'$ and $f: Fc \to d$, $\phi_{c,d'}(k \circ f) = Gk \circ \phi_{c,d}(f)$.
\end{definition}

\section{Unit and Counit}

Every adjunction gives rise to canonical natural transformations, the \emph{unit} and \emph{counit}, which encode the adjunction in a particularly transparent way.

\begin{definition}[Unit and counit]
Given an adjunction $F \dashv G$ with natural bijection $\phi$, define:
\begin{itemize}
  \item The \textbf{unit} $\eta: \Id_{\cat{C}} \Rightarrow GF$: the component $\eta_c: c \to GFc$ is $\phi_{c,Fc}(\id_{Fc})$.
  \item The \textbf{counit} $\varepsilon: FG \Rightarrow \Id_{\cat{D}}$: the component $\varepsilon_d: FGd \to d$ is $\phi^{-1}_{Gd,d}(\id_{Gd})$.
\end{itemize}
\end{definition}

\begin{theorem}[Triangle identities]
\label{thm:triangle-identities}
The unit and counit of an adjunction $F \dashv G$ satisfy the \textbf{triangle identities}:
\[
  (\varepsilon F) \circ (F\eta) = \id_F \qquad \text{and} \qquad (G\varepsilon) \circ (\eta G) = \id_G.
\]
Componentwise: $\varepsilon_{Fc} \circ F(\eta_c) = \id_{Fc}$ and $G(\varepsilon_d) \circ \eta_{Gd} = \id_{Gd}$.
\[
  \begin{tikzcd}
    Fc \arrow[r, "F\eta_c"] \arrow[dr, "\id_{Fc}"'] & FGFc \arrow[d, "\varepsilon_{Fc}"] \\
    & Fc
  \end{tikzcd}
  \qquad
  \begin{tikzcd}
    Gd \arrow[r, "\eta_{Gd}"] \arrow[dr, "\id_{Gd}"'] & GFGd \arrow[d, "G\varepsilon_d"] \\
    & Gd
  \end{tikzcd}
\]
\end{theorem}

\begin{proof}
For the first identity: $\varepsilon_{Fc} \circ F(\eta_c) = \phi^{-1}(\id_{GFc}) \circ F(\phi(\id_{Fc}))$. By naturality of $\phi$ in $c$ applied to $\eta_c: c \to GFc$, we have $\phi(\varepsilon_{Fc} \circ F\eta_c) = G(\varepsilon_{Fc}) \circ \phi(F\eta_c) = G(\varepsilon_{Fc}) \circ \eta_{GFc} \circ \eta_c$... (we omit the technical verification and refer to \cite{MacLane1998}, Chapter IV).
\end{proof}

\section{The Unit-Counit Definition}

The triangle identities characterize adjunctions entirely in terms of the unit and counit, without reference to the hom-set bijection.

\begin{theorem}[Adjunction, unit-counit version]
A pair of functors $F: \cat{C} \to \cat{D}$ and $G: \cat{D} \to \cat{C}$ together with natural transformations $\eta: \Id_{\cat{C}} \Rightarrow GF$ and $\varepsilon: FG \Rightarrow \Id_{\cat{D}}$ satisfying the triangle identities constitutes an adjunction $F \dashv G$.

The hom-set bijection is recovered by:
\[
  \phi_{c,d}(f) = Gf \circ \eta_c, \qquad \phi^{-1}_{c,d}(g) = \varepsilon_d \circ Fg.
\]
\end{theorem}

\begin{proof}
We verify $\phi$ and $\phi^{-1}$ are inverse: $\phi^{-1}(\phi(f)) = \varepsilon_d \circ F(Gf \circ \eta_c) = \varepsilon_d \circ FGf \circ F\eta_c$. By naturality of $\varepsilon$ at $f: Fc \to d$: $\varepsilon_d \circ FGf = f \circ \varepsilon_{Fc}$. So we get $f \circ \varepsilon_{Fc} \circ F\eta_c = f \circ \id_{Fc} = f$ (using the first triangle identity). The other direction is dual.
\end{proof}

\section{Equivalent Formulations of Adjunctions}

There are (at least) four equivalent ways to define an adjunction. Each illuminates a different aspect.

\begin{theorem}[Four equivalent definitions]
The following are equivalent for functors $F: \cat{C} \to \cat{D}$ and $G: \cat{D} \to \cat{C}$:
\begin{enumerate}
  \item[(i)] A natural bijection $\cat{D}(Fc,d) \cong \cat{C}(c, Gd)$.
  \item[(ii)] Natural transformations $\eta: \Id_{\cat{C}} \Rightarrow GF$ and $\varepsilon: FG \Rightarrow \Id_{\cat{D}}$ satisfying the triangle identities.
  \item[(iii)] A natural transformation $\eta: \Id_{\cat{C}} \Rightarrow GF$ such that for every $c \in \cat{C}$, $d \in \cat{D}$, and $f: c \to Gd$, there is a unique $\bar{f}: Fc \to d$ with $G\bar{f} \circ \eta_c = f$. (Universality of the unit.)
  \item[(iv)] A representation for each functor $\cat{D}(F-,d): \cat{C}\op \to \Set$ by $Gd$ (and, dually, each $\cat{C}(c,G-)$).
\end{enumerate}
\end{theorem}

Formulation (iii) says: $\eta_c: c \to GFc$ is the \emph{universal arrow} from $c$ to $G$.

\section{Examples of Adjunctions}

\begin{example}[Free-forgetful adjunction]
Let $U: \Grp \to \Set$ be the forgetful functor and $F: \Set \to \Grp$ the free group functor. Then $F \dashv U$. The unit $\eta_S: S \to UF(S)$ is the inclusion of the set $S$ into the free group $F(S)$ as generators. The universal property says: any set-function $f: S \to UG$ extends uniquely to a group homomorphism $\bar{f}: F(S) \to G$. The counit $\varepsilon_G: FU(G) \to G$ is the unique homomorphism extending $\id: UG \to UG$, sending each generator to itself.

The same pattern applies to free monoids, free rings, free $k$-algebras, free $R$-modules, etc.
\end{example}

\begin{example}[Product-diagonal-coproduct]
The diagonal functor $\Delta: \cat{C} \to \cat{C} \times \cat{C}$, $c \mapsto (c,c)$, sits in a chain of adjunctions:
\[
  \coprod \;\dashv\; \Delta \;\dashv\; \prod.
\]
That is, $\cat{C} \times \cat{C}((a,b), \Delta c) \cong \cat{C}(a,c) \times \cat{C}(b,c) \cong \cat{C}(c, a \times b)$ (when products exist), and $\cat{C}(\Delta c, (a,b)) \cong \cat{C}(c,a) \times \cat{C}(c,b)$... Actually more precisely: $a \sqcup b \dashv \Delta \dashv a \times b$ in the sense that coproduct is left adjoint and product is right adjoint to the diagonal.
\end{example}

\begin{example}[Currying: Tensor-Hom adjunction]
For a ring $R$ and a fixed right $R$-module $M$, there is an adjunction:
\[
  M \otimes_R - \;\dashv\; \Hom_R(M,-),
\]
between the categories of left $R$-modules. The natural bijection is:
\[
  \Hom_R(M \otimes_R N, P) \;\cong\; \Hom_R(N, \Hom_R(M, P)).
\]
In the case $R = k$ (a field), this specializes to the tensor-hom adjunction $V \otimes_k - \;\dashv\; \Hom_k(V, -)$ for vector spaces. In $\Set$ (with Cartesian product), the corresponding statement is $A \times - \;\dashv\; A \to -$ (currying/exponential adjunction): $\Set(A \times B, C) \cong \Set(B, C^A)$.
\end{example}

\begin{example}[Sheafification]
The inclusion $i: \mathbf{Sh}(X) \hookrightarrow \mathbf{PSh}(X)$ of sheaves into presheaves on a topological space $X$ has a left adjoint, the sheafification functor $a: \mathbf{PSh}(X) \to \mathbf{Sh}(X)$:
\[
  a \;\dashv\; i.
\]
The unit $\eta_P: P \to i(aP)$ is the natural map from a presheaf to its sheafification.
\end{example}

\begin{example}[Stone-Čech compactification]
The inclusion $i: \mathbf{KHaus} \hookrightarrow \Top$ of compact Hausdorff spaces has a left adjoint, the Stone-Čech compactification $\beta: \Top \to \mathbf{KHaus}$: $\beta \dashv i$ in the sense that $\Top(X, iK) \cong \mathbf{KHaus}(\beta X, K)$.
\end{example}

\begin{example}[Galois connections]
Let $(P, \leq)$ and $(Q, \leq)$ be posets, viewed as categories. A functor $F: P \to Q$ is simply an order-preserving map. An adjunction $F \dashv G$ between poset-categories is precisely a \textbf{Galois connection}: for all $p \in P$ and $q \in Q$,
\[
  Fp \leq q \;\iff\; p \leq Gq.
\]
The triangle identities become $p \leq G(Fp)$ (unit) and $F(Gq) \leq q$ (counit). Galois connections pervade order theory, Galois theory, topology (interior/closure), and logic (proof/model).
\end{example}

\begin{example}[Reflective subcategories]
A full subcategory $\cat{D} \hookrightarrow \cat{C}$ is \textbf{reflective} if the inclusion has a left adjoint (the \emph{reflector}). Examples:
\begin{itemize}
  \item Abelian groups in groups: the reflector is abelianization $G \mapsto G/[G,G]$.
  \item Complete metric spaces in metric spaces: the reflector is completion.
  \item Sheaves in presheaves: the reflector is sheafification.
  \item $T_0$, $T_1$, $T_2$ spaces within $\Top$: various separation-quotient reflectors.
\end{itemize}
\end{example}

\section{Adjunctions and Limits}

The most important structural fact about adjunctions is how they interact with limits and colimits.

\begin{theorem}[RAPL and LAPC]
\label{thm:rapl-lapc}
Let $F \dashv G$ be an adjunction with $F: \cat{C} \to \cat{D}$ and $G: \cat{D} \to \cat{C}$.
\begin{enumerate}
  \item[(RAPL)] \textbf{Right Adjoints Preserve Limits:} $G$ preserves all limits that exist in $\cat{D}$.
  \item[(LAPC)] \textbf{Left Adjoints Preserve Colimits:} $F$ preserves all colimits that exist in $\cat{C}$.
\end{enumerate}
\end{theorem}

\begin{proof}
We prove RAPL; LAPC is dual. Let $D: \cat{J} \to \cat{D}$ be a diagram with limit $(L, \pi)$. We show $(GL, G\pi)$ is a limit of $GD$ in $\cat{C}$.

For any cone $(c, \lambda)$ over $GD$—i.e., morphisms $\lambda_j: c \to GD(j)$ natural in $j$—we need a unique $u: c \to GL$ with $G\pi_j \circ u = \lambda_j$.

By the adjunction bijection, each $\lambda_j: c \to GD(j)$ corresponds to $\bar\lambda_j: Fc \to D(j)$, and the naturality of $\lambda$ implies the $\bar\lambda_j$ form a cone over $D$. By universality of $(L, \pi)$, there is a unique $\bar u: Fc \to L$ with $\pi_j \circ \bar u = \bar\lambda_j$. Applying $\phi^{-1}$, this corresponds to a unique $u: c \to GL$ with the required property.

More elegantly: using representability, $\cat{C}(c, G\lim D) \cong \cat{D}(Fc, \lim D) \cong \lim_j \cat{D}(Fc, D(j)) \cong \lim_j \cat{C}(c, GD(j))$, naturally in $c$, which shows $\lim_j GD(j) \cong G \lim_j D(j)$.
\end{proof}

\begin{example}[Consequences]
\begin{itemize}
  \item The forgetful functor $U: \Grp \to \Set$ preserves limits (products, equalizers, etc.) because it has a left adjoint (free group).
  \item The functor $A \times -: \Set \to \Set$ preserves colimits (disjoint unions, coequalizers, pushouts) because it has a right adjoint $(-)^A$.
  \item Sheafification preserves colimits.
  \item Stone-Čech compactification preserves finite coproducts (a less obvious fact).
\end{itemize}
\end{example}

\section{Uniqueness of Adjoints}

\begin{theorem}
Right adjoints are unique up to unique natural isomorphism: if $F \dashv G$ and $F \dashv G'$, then $G \cong G'$ by a unique natural isomorphism. Similarly, left adjoints are unique up to unique natural isomorphism.
\end{theorem}

\begin{proof}
$\cat{C}(c, Gd) \cong \cat{D}(Fc,d) \cong \cat{C}(c, G'd)$ naturally in $c$ and $d$. By the Yoneda lemma (applied to each $d$), the bijection $\cat{C}(c,Gd) \cong \cat{C}(c,G'd)$ natural in $c$ gives a unique isomorphism $Gd \cong G'd$ for each $d$; naturality in $d$ makes this a natural isomorphism.
\end{proof}

\section{Adjoint Functor Theorems}

The adjoint functor theorems give general conditions under which a functor has a left or right adjoint.

\begin{theorem}[General Adjoint Functor Theorem, GAFT]
\label{thm:gaft}
Let $\cat{C}$ be a complete, locally small category, and $G: \cat{C} \to \cat{D}$ a functor. Then $G$ has a left adjoint if and only if $G$ preserves all small limits and satisfies the \textbf{solution set condition}: for each $d \in \cat{D}$, there exists a set $\{f_i: d \to G(c_i)\}_{i \in I}$ (a \emph{solution set}) such that every morphism $f: d \to Gc$ factors as $f = Gg \circ f_i$ for some $i$ and some $g: c_i \to c$.
\end{theorem}

\begin{proof}[Proof sketch]
For each $d \in \cat{D}$, we seek an initial object in the comma category $(d \downarrow G)$. The solution set ensures we can form a small limit over the solution set to construct this initial object. Completeness of $\cat{C}$ (together with $G$ preserving limits) ensures this limit exists and has the right properties.
\end{proof}

\begin{theorem}[Special Adjoint Functor Theorem, SAFT]
\label{thm:saft}
Let $\cat{C}$ be a complete, locally small, and \textbf{well-powered} category (every object has a set of subobjects) that has a \textbf{cogenerating set} (a set $S \subseteq \ob(\cat{C})$ such that morphisms are detected by maps into elements of $S$). Then any limit-preserving functor $G: \cat{C} \to \cat{D}$ has a left adjoint.
\end{theorem}

\begin{example}
$\Set$ has a cogenerating set $\{2\}$ (the two-element set): maps $f, g: A \to B$ agree iff $h \circ f = h \circ g$ for all $h: B \to 2$. So any limit-preserving functor from $\Set$ has a left adjoint. Similarly, $\Grp$ and $\Ab$ satisfy the hypotheses of SAFT.
\end{example}

\begin{theorem}[Freyd's Adjoint Functor Theorem]
Let $\cat{C}$ be a locally small category with all small limits. A functor $G: \cat{C} \to \cat{D}$ has a left adjoint if and only if for every $d \in \cat{D}$, the comma category $(d \downarrow G)$ has an initial object.
\end{theorem}

This is a tautological reformulation, but it is useful: it reduces the existence of adjoints to finding initial objects in comma categories, which can sometimes be checked directly.

\section{Equivalences as Adjunctions}

\begin{proposition}
An equivalence of categories $F: \cat{C} \simeq \cat{D}$ consists of two adjunctions $F \dashv G$ and $G \dashv F$, where $\eta$ and $\varepsilon$ are natural isomorphisms. Such an adjunction is called an \textbf{adjoint equivalence}.
\end{proposition}

\begin{proposition}
Any equivalence can be promoted to an adjoint equivalence: if $F \dashv G$ with unit and counit that are natural isomorphisms, then it is already an adjoint equivalence.
\end{proposition}

\section{Exercises}

\begin{exercise}
Verify that the free-forgetful adjunction $F \dashv U: \Grp \to \Set$ satisfies the triangle identities.
\end{exercise}

\begin{exercise}
Show that if $F \dashv G$, then $F$ preserves all epimorphisms and $G$ preserves all monomorphisms.
\end{exercise}

\begin{exercise}
Let $F \dashv G$ and $F' \dashv G'$ with $F: \cat{C} \to \cat{D}$ and $F': \cat{D} \to \cat{E}$. Show that $F'F \dashv GG'$.
\end{exercise}

\begin{exercise}
Prove that the solution set condition is necessary for the existence of a left adjoint: if $G$ has a left adjoint $F$, then for each $d \in \cat{D}$, the singleton $\{\eta_d: d \to GFd\}$ is a solution set.
\end{exercise}

\begin{exercise}
Describe the Galois connection between the power sets $\mathcal{P}(A)$ and $\mathcal{P}(B)$ induced by a relation $R \subseteq A \times B$. How does this relate to the Galois connection in classical Galois theory?
\end{exercise}

\begin{exercise}[$\star$]
Prove that a reflective subcategory $i: \cat{D} \hookrightarrow \cat{C}$ with reflector $L \dashv i$ is \textbf{closed under limits} in $\cat{C}$ (limits of diagrams in $\cat{D}$, taken in $\cat{C}$, lie in $\cat{D}$). Is $\cat{D}$ closed under colimits?
\end{exercise}

\begin{exercise}[$\star$]
Prove that any functor $F: \cat{C} \to \cat{D}$ with $\cat{C}$ small and $\cat{D}$ locally small gives a functor $\Lan_F: [\cat{C}, \Set] \to [\cat{D},\Set]$ left adjoint to $F^*: [\cat{D},\Set] \to [\cat{C},\Set]$ (restriction). (This is a preview of Chapter~\ref{ch:kan}.)
\end{exercise}

\begin{exercise}[$\star\star$]
Prove the General Adjoint Functor Theorem in detail, constructing the left adjoint explicitly using the limit of the comma category (d ↓ G) over its solution set.
\end{exercise}

\chapter{Monads}
\label{ch:monads}

\section{From Adjunctions to Monads}

Every adjunction $F \dashv G$ with $F: \cat{C} \to \cat{D}$ and $G: \cat{D} \to \cat{C}$ gives rise to an endofunctor $T = GF: \cat{C} \to \cat{C}$ together with two natural transformations:
\begin{itemize}
  \item The \textbf{unit} $\eta: \Id_{\cat{C}} \Rightarrow GF = T$.
  \item The \textbf{multiplication} $\mu: T^2 = GFGF \Rightarrow GF = T$, defined as $\mu = G\varepsilon F$ (whiskering the counit by $G$ and $F$).
\end{itemize}

These data satisfy axioms analogous to those of a monoid (with $T$ in the role of the set, $\eta$ in the role of unit, and $\mu$ in the role of multiplication), giving the definition of a monad.

\section{Definition of a Monad}

\begin{definition}[Monad]
A \textbf{monad} on a category $\cat{C}$ is a triple $(T, \eta, \mu)$ where $T: \cat{C} \to \cat{C}$ is an endofunctor, $\eta: \Id_{\cat{C}} \Rightarrow T$ and $\mu: T^2 \Rightarrow T$ are natural transformations, satisfying:
\begin{description}
  \item[Associativity.] $\mu \circ T\mu = \mu \circ \mu T$ (as natural transformations $T^3 \Rightarrow T$).
  \item[Unit laws.] $\mu \circ T\eta = \id_T = \mu \circ \eta T$.
\end{description}
Componentwise, for each $c \in \cat{C}$:
\[
  \mu_c \circ T(\mu_c) = \mu_c \circ \mu_{Tc}, \quad
  \mu_c \circ T(\eta_c) = \id_{Tc} = \mu_c \circ \eta_{Tc}.
\]
\end{definition}

The axioms say that $(T, \eta, \mu)$ is a \emph{monoid in the monoidal category of endofunctors} of $\cat{C}$ (with monoidal product given by functor composition and unit $\Id_{\cat{C}}$).

\begin{remark}[Comonad]
A \textbf{comonad} $(S, \varepsilon, \delta)$ is the dual: an endofunctor with $\varepsilon: S \Rightarrow \Id$ (counit) and $\delta: S \Rightarrow S^2$ (comultiplication) satisfying the co-associativity and co-unit laws. Every adjunction gives both a monad on $\cat{C}$ and a comonad on $\cat{D}$.
\end{remark}

\section{Examples of Monads}

\begin{example}[The list monad]
On $\Set$: $T(X) = X^* = \bigsqcup_{n \geq 0} X^n$ (the set of finite lists over $X$). The unit $\eta_X: X \to X^*$ sends $x$ to the singleton list $[x]$. The multiplication $\mu_X: (X^*)^* \to X^*$ flattens a list of lists into a list (concatenation). This monad captures the idea of ``nondeterministic computation.''
\end{example}

\begin{example}[The maybe monad]
On $\Set$: $T(X) = X + \{*\}$ (add a distinguished element ``nothing''). $\eta_X: X \to X + \{*\}$ is the inclusion. $\mu_X: X + \{*\} + \{*\} \to X + \{*\}$ collapses both $\{*\}$'s. Used in programming for partial/exceptional computations.
\end{example}

\begin{example}[The power-set monad]
On $\Set$: $T(X) = \mathcal{P}(X)$. $\eta_X(x) = \{x\}$ (singleton). $\mu_X: \mathcal{P}(\mathcal{P}(X)) \to \mathcal{P}(X)$, $\mathcal{A} \mapsto \bigcup \mathcal{A}$. This is the monad arising from the adjunction between $\Set$ and $\mathbf{Rel}$.
\end{example}

\begin{example}[The free-algebra monads]
For any algebraic theory (groups, rings, modules, etc.), the composite $GF$ of the free functor $F$ and forgetful functor $G$ gives a monad on $\Set$:
\begin{itemize}
  \item Free group monad: $T(X) = $ free group on $X$.
  \item Free monoid monad: $T(X) = X^*$ (lists—the list monad above).
  \item Free $k$-vector space monad: $T(X) = k^{(X)}$.
  \item Free commutative ring monad: $T(X) = \mathbb{Z}[X]$ (polynomial ring).
\end{itemize}
\end{example}

\begin{example}[The state monad]
For a fixed set $S$ (``state''), on $\Set$: $T(X) = (X \times S)^S$ (functions from $S$ to $X \times S$). This monad encapsulates stateful computations.
\end{example}

\begin{example}[The continuation monad]
For a fixed set $R$ (``result type''): $T(X) = R^{R^X}$ (double-dualization with respect to $R$). Unit: $\eta_X(x) = (\phi \mapsto \phi(x))$. This monad encapsulates continuation-passing style.
\end{example}

\section{Kleisli Category}

Given a monad $(T, \eta, \mu)$ on $\cat{C}$, we can construct a new category whose morphisms are ``$T$-computations.''

\begin{definition}[Kleisli category]
The \textbf{Kleisli category} $\cat{C}_T$ of a monad $(T, \eta, \mu)$ has:
\begin{itemize}
  \item Objects: the same as $\cat{C}$.
  \item Morphisms: $\cat{C}_T(a,b) = \cat{C}(a, Tb)$ (a ``Kleisli arrow'' from $a$ to $b$ is a morphism $a \to Tb$ in $\cat{C}$).
  \item Composition: Given $f: a \to Tb$ and $g: b \to Tc$ in $\cat{C}$, their Kleisli composite is
  \[
    g \star f = \mu_c \circ Tg \circ f: a \to Tb \to T^2c \to Tc.
  \]
  \item Identities: $\id^T_a = \eta_a: a \to Ta$.
\end{itemize}
\end{definition}

\begin{proposition}
$\cat{C}_T$ is a category: composition is associative and the Kleisli identities are units.
\end{proposition}

\begin{proof}
Associativity of $\star$ follows from the monad associativity axiom $\mu \circ T\mu = \mu \circ \mu T$. The unit laws follow from $\mu \circ T\eta = \id_T = \mu \circ \eta T$.
\end{proof}

\begin{proposition}[Kleisli adjunction]
There is an adjunction $F_T \dashv U_T$ where:
\begin{itemize}
  \item $F_T: \cat{C} \to \cat{C}_T$: identity on objects, sends $f: a \to b$ to $\eta_b \circ f: a \to Tb$.
  \item $U_T: \cat{C}_T \to \cat{C}$: sends an object $a$ to $Ta$; sends $f: a \to Tb$ (a Kleisli morphism) to $\mu_b \circ Tf: Ta \to T^2b \to Tb$.
\end{itemize}
The monad induced by $F_T \dashv U_T$ is $(T, \eta, \mu)$.
\end{proposition}

\begin{example}
For the list monad, $\cat{C}_T$ has sets as objects and a Kleisli arrow $f: X \to Y$ is a function $f: X \to Y^*$ (a nondeterministic function assigning to each $x \in X$ a list of possible outputs). Kleisli composition is nondeterministic composition.
\end{example}

\section{Eilenberg-Moore Category}

The Kleisli construction gives the ``smallest'' adjunction inducing a given monad. The Eilenberg-Moore construction gives the ``largest.''

\begin{definition}[$T$-algebra]
Let $(T, \eta, \mu)$ be a monad on $\cat{C}$. A \textbf{$T$-algebra} (or \textbf{algebra for $T$}) is a pair $(a, \theta)$ where $a \in \cat{C}$ and $\theta: Ta \to a$ is a morphism in $\cat{C}$, such that:
\[
  \theta \circ \eta_a = \id_a \qquad \text{(unit law)}
\]
\[
  \theta \circ \mu_a = \theta \circ T\theta \qquad \text{(associativity law)}.
\]
The first says $\theta$ is a left inverse of $\eta_a$; the second says:
\[
  \begin{tikzcd}
    T^2 a \arrow[r, "T\theta"] \arrow[d, "\mu_a"'] & Ta \arrow[d, "\theta"] \\
    Ta \arrow[r, "\theta"'] & a
  \end{tikzcd}
  \quad \text{commutes.}
\]
\end{definition}

\begin{example}
For the free-group monad $T$ on $\Set$, a $T$-algebra $(A, \theta)$ is a set $A$ with a map $\theta: F(A) \to A$ from the free group on $A$ to $A$, satisfying the two axioms. This is precisely a group: $\theta$ defines the group structure (multiplication of words in $A$ is determined by $\theta$). The unit law says the identity element maps to the neutral element of $A$, and the associativity law says the group multiplication is associative.
\end{example}

\begin{definition}[Eilenberg-Moore category]
The \textbf{Eilenberg-Moore category} $\cat{C}^T$ has:
\begin{itemize}
  \item Objects: $T$-algebras $(a, \theta)$.
  \item Morphisms: A morphism of $T$-algebras $h: (a, \theta) \to (a', \theta')$ is a morphism $h: a \to a'$ in $\cat{C}$ with $\theta' \circ Th = h \circ \theta$.
  \item Composition and identities: inherited from $\cat{C}$.
\end{itemize}
\end{definition}

\begin{proposition}[Eilenberg-Moore adjunction]
There is an adjunction $F^T \dashv U^T$ where:
\begin{itemize}
  \item $U^T: \cat{C}^T \to \cat{C}$: forgets the algebra structure, $(a,\theta) \mapsto a$.
  \item $F^T: \cat{C} \to \cat{C}^T$: sends $a$ to the \textbf{free $T$-algebra} $(Ta, \mu_a)$.
\end{itemize}
The monad induced by $F^T \dashv U^T$ is $(T, \eta, \mu)$.
\end{proposition}

\begin{proof}
The hom-set bijection: $\cat{C}^T(F^T a, (b, \theta)) = \cat{C}^T((Ta, \mu_a), (b, \theta))$ consists of algebra maps $h: Ta \to b$ with $\theta \circ Th = h \circ \mu_a$. These correspond bijectively to morphisms $\hat{h} = h \circ \eta_a: a \to b$ in $\cat{C}$ via $h \mapsto h \circ \eta_a$ (inverse: $\hat{h} \mapsto \theta \circ T\hat{h}$).
\end{proof}

\section{Comparison Functor and Monadic Adjunctions}

Given any adjunction $F \dashv G$ with induced monad $T = GF$, there is a canonical comparison functor to the Eilenberg-Moore category.

\begin{definition}[Comparison functor]
Given $F \dashv G$ (with $F: \cat{C} \to \cat{D}$, $G: \cat{D} \to \cat{C}$, monad $T = GF$), define the \textbf{comparison functor} $K: \cat{D} \to \cat{C}^T$ by
\[
  K(d) = (Gd, G\varepsilon_d), \quad K(f: d \to d') = Gf.
\]
One verifies $(Gd, G\varepsilon_d)$ is indeed a $T$-algebra: $G\varepsilon_d \circ \eta_{Gd} = G\varepsilon_d \circ G(F\eta_d \circ \eta_d)... $; actually $G\varepsilon_d \circ \eta_{Gd} = \id_{Gd}$ by the triangle identity.
\end{definition}

\begin{definition}[Monadic adjunction]
The adjunction $F \dashv G$ is \textbf{monadic} if the comparison functor $K: \cat{D} \to \cat{C}^T$ is an equivalence of categories.
\end{definition}

\begin{example}
The free-forgetful adjunction for groups is monadic: $\cat{D} = \Grp$ and $\cat{C}^T = \Grp$ (algebras for the free-group monad are exactly groups). Similarly for rings, modules, etc. The free-forgetful adjunction for topological spaces is NOT monadic: topological spaces are not determined by an algebraic structure on their underlying set.
\end{example}

\section{Beck's Monadicity Theorem}

Beck's monadicity theorem gives a precise characterization of when an adjunction is monadic.

\begin{definition}[$U$-split coequalizer]
A coequalizer diagram
\[
  \begin{tikzcd}
    a \arrow[r, bend left, "f"] \arrow[r, bend right, "g"'] & b \arrow[r, "q"] & c
  \end{tikzcd}
\]
is a \textbf{$U$-split coequalizer} (for $U: \cat{D} \to \cat{C}$) if $U$ applied to the diagram is a split coequalizer in $\cat{C}$: there exist sections $s: Uc \to Ub$ and $t: Ub \to Ua$ such that $Uq \circ s = \id_{Uc}$, $Uf \circ t = \id_{Ub}$, and $Ug \circ t = s \circ Uq$.
\end{definition}

\begin{theorem}[Beck's Monadicity Theorem]
\label{thm:beck}
Let $F \dashv U$ be an adjunction with $F: \cat{C} \to \cat{D}$ and $U: \cat{D} \to \cat{C}$. The following are equivalent:
\begin{enumerate}
  \item The adjunction is monadic.
  \item $U$ reflects isomorphisms, and $\cat{D}$ has, and $U$ preserves, coequalizers of $U$-split pairs.
\end{enumerate}
\end{theorem}

\begin{proof}[Proof sketch]
($1 \Rightarrow 2$): If $K: \cat{D} \to \cat{C}^T$ is an equivalence, then $U = U^T \circ K$ reflects isomorphisms (equivalences reflect isos, and $U^T$ creates isos). One shows $\cat{D}$ has coequalizers of $U$-split pairs by transporting them from $\cat{C}^T$.

($2 \Rightarrow 1$): We show $K: \cat{D} \to \cat{C}^T$ is an equivalence. Faithfulness and fullness of $K$ follow from $U$ reflecting isomorphisms and the universal properties. For essential surjectivity, given a $T$-algebra $(a, \theta)$, form the $U$-split coequalizer
\[
  FUFa \;\underset{F\theta}{\overset{\varepsilon_{Fa}}{\rightrightarrows}}\; Fa
\]
(where the splitting is via $\eta_{UFa}$ and $\eta_{Fa}$). The coequalizer $d$ of this diagram in $\cat{D}$ (which exists by hypothesis) satisfies $K(d) \cong (a, \theta)$.
\end{proof}

\begin{example}[Applications of Beck's theorem]
\begin{itemize}
  \item The forgetful functor $\Ab \to \Set$ is monadic (free abelian groups).
  \item Compact Hausdorff spaces over $\Set$: the adjunction $\beta \dashv i$ (Stone-Čech) is monadic: $\mathbf{KHaus}$ is equivalent to algebras for the ultrafilter monad.
  \item Descent theory: Beck's theorem is the categorical foundation for descent theory in algebraic geometry.
\end{itemize}
\end{example}

\section{Distributive Laws}

Monads can sometimes be composed using distributive laws.

\begin{definition}[Distributive law]
Let $(S, \eta^S, \mu^S)$ and $(T, \eta^T, \mu^T)$ be monads on $\cat{C}$. A \textbf{distributive law} of $S$ over $T$ is a natural transformation $\lambda: ST \Rightarrow TS$ satisfying four coherence conditions (compatibility with units and multiplications of both monads).
\end{definition}

\begin{theorem}
Given a distributive law $\lambda: ST \Rightarrow TS$, the composite $TS$ carries the structure of a monad on $\cat{C}$, with unit $\eta^T \eta^S$ and multiplication built from $\mu^T$, $\mu^S$, and $\lambda$.
\end{theorem}

\begin{example}
The distributive law $\lambda: \mathcal{P}\mathcal{P}' \Rightarrow \mathcal{P}'\mathcal{P}$ (where $\mathcal{P}$ is nondeterminism/power set and $\mathcal{P}'$ some other monad) models combinations of computational effects.
\end{example}

\section{Monads in 2-Categories}

The definition of a monad generalizes immediately from $\Cat$ to any 2-category.

\begin{definition}[Monad in a 2-category]
Let $\mathcal{K}$ be a 2-category and $C$ an object of $\mathcal{K}$. A \textbf{monad} on $C$ in $\mathcal{K}$ is a 1-cell $T: C \to C$ together with 2-cells $\eta: \Id_C \Rightarrow T$ and $\mu: T \circ T \Rightarrow T$ satisfying the monad axioms.
\end{definition}

In $\Cat$, this recovers the usual notion. In the 2-category of rings, bimodules, and bimodule maps, a monad on a ring $R$ is an $R$-algebra. This level of generality unifies many constructions.

\section{Exercises}

\begin{exercise}
Verify the monad axioms for the list monad on $\Set$.
\end{exercise}

\begin{exercise}
Show that the Kleisli composition $g \star f = \mu_c \circ Tg \circ f$ is associative using the monad axioms.
\end{exercise}

\begin{exercise}
Describe the Eilenberg-Moore algebras for the power-set monad $\mathcal{P}$. What kind of structure do they capture? (They are \emph{complete semilattices}.)
\end{exercise}

\begin{exercise}
Let $(T, \eta, \mu)$ be a monad on $\cat{C}$. Show that $\cat{C}_T$ (Kleisli) and $\cat{C}^T$ (Eilenberg-Moore) both induce the monad $T$ on $\cat{C}$. In what sense are these the ``initial'' and ``terminal'' adjunctions inducing $T$?
\end{exercise}

\begin{exercise}
Show that there is a functor $L: \cat{C}_T \to \cat{C}^T$ sending a Kleisli object $a$ to the free algebra $(Ta, \mu_a)$, and identify this as the composition $\cat{C}_T \to \cat{C} \xrightarrow{F^T} \cat{C}^T$.
\end{exercise}

\begin{exercise}[$\star$]
Prove that the comparison functor $K: \cat{D} \to \cat{C}^T$ commutes with the forgetful functors: $U^T \circ K = G$ (not just up to isomorphism, but on the nose).
\end{exercise}

\begin{exercise}[$\star$]
A monad $(T, \eta, \mu)$ is called \textbf{idempotent} if $\mu$ is a natural isomorphism. Show that for an idempotent monad, the comparison functor $K$ is fully faithful (so monadic adjunctions for idempotent monads are reflective subcategory inclusions).
\end{exercise}

\begin{exercise}[$\star\star$]
Prove one direction of Beck's monadicity theorem: assuming the adjunction is monadic (comparison functor is an equivalence), deduce that $U$ reflects isomorphisms and that $\cat{D}$ has coequalizers of $U$-split pairs preserved by $U$.
\end{exercise}

\chapter{Kan Extensions}
\label{ch:kan}

\section{The Fundamental Question}

Consider functors $K: \cat{C} \to \cat{D}$ and $F: \cat{C} \to \cat{E}$. Can we ``extend'' $F$ through $K$ to obtain a functor $\cat{D} \to \cat{E}$? That is, can we fill in the dashed arrow:
\[
  \begin{tikzcd}
    \cat{C} \arrow[r, "K"] \arrow[dr, "F"'] & \cat{D} \arrow[d, dashed] \\
    & \cat{E}
  \end{tikzcd}
\]

In general, no extension exists on the nose. But we can ask for the \emph{best approximation}. There are two canonical choices: the \emph{left Kan extension} (which approximates from below, via colimits) and the \emph{right Kan extension} (which approximates from above, via limits). Kan extensions subsume virtually every other universal construction in category theory—limits, colimits, adjunctions, and more.

As Mac Lane famously wrote: ``The notion of Kan extensions subsumes all the other fundamental concepts of category theory.''

\section{Right Kan Extensions}

\begin{definition}[Right Kan extension]
Let $K: \cat{C} \to \cat{D}$ and $F: \cat{C} \to \cat{E}$ be functors. A \textbf{right Kan extension} of $F$ along $K$ is a pair $(\Ran_K F, \varepsilon)$ where:
\begin{itemize}
  \item $\Ran_K F: \cat{D} \to \cat{E}$ is a functor.
  \item $\varepsilon: (\Ran_K F) \circ K \Rightarrow F$ is a natural transformation (the \textbf{counit}).
\end{itemize}
satisfying the universal property: for any functor $G: \cat{D} \to \cat{E}$ and natural transformation $\alpha: GK \Rightarrow F$, there is a unique natural transformation $\bar\alpha: G \Rightarrow \Ran_K F$ such that $\varepsilon \circ (\bar\alpha K) = \alpha$.
\[
  \begin{tikzcd}
    \cat{C} \arrow[r, "K"] \arrow[dr, "F"'] & \cat{D} \arrow[d, "\Ran_K F"] \arrow[dl, "G"', dashed, bend right] \\
    & \cat{E}
  \end{tikzcd}
  \quad \varepsilon: (\Ran_K F)K \Rightarrow F, \quad \bar\alpha: G \Rightarrow \Ran_K F
\]
Equivalently, $\Ran_K F$ represents the functor
\[
  G \;\mapsto\; \Nat(GK, F),
\]
giving a natural bijection $\Nat(G, \Ran_K F) \cong \Nat(GK, F)$.
\end{definition}

\section{Left Kan Extensions}

\begin{definition}[Left Kan extension]
A \textbf{left Kan extension} of $F: \cat{C} \to \cat{E}$ along $K: \cat{C} \to \cat{D}$ is a pair $(\Lan_K F, \eta)$ where:
\begin{itemize}
  \item $\Lan_K F: \cat{D} \to \cat{E}$ is a functor.
  \item $\eta: F \Rightarrow (\Lan_K F) \circ K$ is a natural transformation (the \textbf{unit}).
\end{itemize}
satisfying the universal property: for any $G: \cat{D} \to \cat{E}$ and $\beta: F \Rightarrow GK$, there is a unique $\bar\beta: \Lan_K F \Rightarrow G$ such that $(\bar\beta K) \circ \eta = \beta$.

Equivalently, $\Lan_K F$ corepresents the functor $G \mapsto \Nat(F, GK)$:
\[
  \Nat(\Lan_K F, G) \;\cong\; \Nat(F, GK).
\]
\end{definition}

\begin{remark}[Duality]
Left and right Kan extensions are dual: $\Ran_K F$ in $\cat{E}$ equals $\Lan_K (F\op)$ in $\cat{E}\op$ (roughly). One can always reduce one to the other by dualizing.
\end{remark}

\section{Pointwise Formula for Kan Extensions}

When $\cat{E}$ is complete (resp.\ cocomplete), Kan extensions exist and can be computed by an explicit (co)limit formula.

\begin{theorem}[Pointwise right Kan extension]
\label{thm:pointwise-ran}
Let $\cat{C}$ be small, $K: \cat{C} \to \cat{D}$, $F: \cat{C} \to \cat{E}$. If $\cat{E}$ has all small limits, then $\Ran_K F$ exists and is given, for $d \in \cat{D}$, by:
\[
  (\Ran_K F)(d) = \lim\bigl(K \downarrow d \xrightarrow{\pi} \cat{C} \xrightarrow{F} \cat{E}\bigr),
\]
where $K \downarrow d$ is the comma category (objects: pairs $(c, f: Kc \to d)$, morphisms: maps in $\cat{C}$ over $d$) and $\pi$ is the projection functor.
\end{theorem}

\begin{theorem}[Pointwise left Kan extension]
\label{thm:pointwise-lan}
Under dual hypotheses ($\cat{E}$ cocomplete):
\[
  (\Lan_K F)(d) = \colim\bigl(d \downarrow K \xrightarrow{\pi} \cat{C} \xrightarrow{F} \cat{E}\bigr),
\]
where $d \downarrow K$ has objects $(c, f: d \to Kc)$.
\end{theorem}

\begin{proof}[Proof of Theorem~\ref{thm:pointwise-lan}]
We verify the universal property. Given $G: \cat{D} \to \cat{E}$ and $\beta: F \Rightarrow GK$, we need a unique $\bar\beta: \Lan_K F \Rightarrow G$.

At a fixed $d \in \cat{D}$, we need a map $(\Lan_K F)(d) = \colim_{(c,f) \in d\downarrow K} F(c) \to G(d)$. For each $(c, f: d \to Kc)$, the component $\beta_c: F(c) \to GKc$, composed with $G(f): GKc \to Gd$, gives $G(f) \circ \beta_c: F(c) \to Gd$. These are compatible with morphisms in $d \downarrow K$, hence define a cocone, hence a unique map from the colimit. Naturality in $d$ is verified by functoriality of everything involved.

Uniqueness: the unit $\eta_c: F(c) \to (\Lan_K F)(Kc)$ corresponds to the map from $F(c)$ to the colimit indexed by $(c, \id_{Kc}) \in Kc \downarrow K$. Any $\bar\beta: \Lan_K F \Rightarrow G$ with $(\bar\beta K) \circ \eta = \beta$ determines, for each $(c, f: d \to Kc)$, the map $F(c) \xrightarrow{\eta_c} (\Lan_K F)(Kc) \xrightarrow{(\Lan_K F)(f)} (\Lan_K F)(d) \xrightarrow{\bar\beta_d} Gd$; since the colimit is generated by the elements of $F(c)$ over all $(c,f)$, $\bar\beta_d$ is determined.
\end{proof}

\section{Adjunctions as Kan Extensions}

\begin{theorem}
\label{thm:adj-as-kan}
Let $F: \cat{C} \to \cat{D}$ and $G: \cat{D} \to \cat{C}$. Then $F \dashv G$ if and only if $G \cong \Ran_F \Id_{\cat{C}}$ (and the Kan extension is pointwise), or equivalently, $F \cong \Lan_G \Id_{\cat{D}}$.
\end{theorem}

\begin{proof}
$\Nat(G', \Ran_F \Id_{\cat{C}}) \cong \Nat(G'F, \Id_{\cat{C}}) = $ \{natural transformations $G'F \Rightarrow \Id_{\cat{C}}\}$. This is the universal property of the right adjoint: $G$ represents $\Nat(-F, \Id_{\cat{C}})$, i.e., $F \dashv G$ iff $G = \Ran_F \Id_{\cat{C}}$.
\end{proof}

\begin{corollary}
If $F \dashv G$, then for any $H: \cat{C} \to \cat{E}$, $\Ran_F H \cong H \circ G$ (and the extension is absolute, computed without colimits).
\end{corollary}

\begin{proof}
$\Nat(L, \Ran_F H) \cong \Nat(LF, H) \cong \Nat(L, HG)$, so $\Ran_F H \cong HG$.
\end{proof}

\section{The Density Theorem}

The density theorem is a fundamental result relating any functor to its own Kan extension along itself.

\begin{theorem}[Density theorem]
\label{thm:density}
Let $K: \cat{C} \to \cat{D}$ be any functor. Then the left Kan extension of $K$ along itself is (naturally isomorphic to) the identity:
\[
  \Lan_K K \;\cong\; \Id_{\cat{D}}.
\]
More precisely, if $\cat{D}$ is cocomplete and $K$ is small (e.g., $\cat{C}$ is small), then the left Kan extension $\Lan_K K: \cat{D} \to \cat{D}$ is isomorphic to $\Id_{\cat{D}}$, i.e., every object of $\cat{D}$ is a colimit of objects in the image of $K$.
\end{theorem}

\begin{proof}
Using the pointwise formula: $(\Lan_K K)(d) = \colim_{(c, Kc \to d)} Kc$. The colimit is taken over the comma category $(d \downarrow K)\op$, which is the category of objects in the image of $K$ with maps to $d$. The universal property of this colimit gives the identity morphism on $d$.

More abstractly: by Theorem~\ref{thm:adj-as-kan}, $\Lan_K K \cong \Id_\cat{D}$ iff the pair $(\Id_\cat{D}, K)$ constitutes an adjunction—but actually the density theorem is stronger and doesn't require an adjunction. The key point is that every $d \in \cat{D}$ is canonically a colimit of objects $Kc$ via the comma category $d \downarrow K$.
\end{proof}

\begin{corollary}[Yoneda as a density statement]
Every presheaf $F \in [\cat{C}\op, \Set]$ is a canonical colimit of representables:
\[
  F \cong \colim_{(c, \yo(c) \to F)} \yo(c).
\]
(The colimit is over the category of elements of $F$.) This is the density theorem for the Yoneda embedding $\yo: \cat{C} \to [\cat{C}\op, \Set]$.
\end{corollary}

\section{Nerve and Geometric Realization}

A classic example of Kan extensions in topology.

\begin{example}[Nerve and realization]
Let $\Delta$ be the simplex category: objects are the finite ordinals $[n] = \{0,1,\ldots,n\}$ and morphisms are weakly order-preserving maps. The standard simplices form a functor $\Delta^\bullet: \Delta \to \Top$, $[n] \mapsto \Delta^n$ (the standard topological $n$-simplex).

The \textbf{singular nerve} of a topological space $X$ is the presheaf $N(X) = \Top(\Delta^\bullet, X): \Delta\op \to \Set$, i.e., $N(X)_n = \Top(\Delta^n, X)$ (singular $n$-simplices in $X$).

The \textbf{geometric realization} of a simplicial set $S: \Delta\op \to \Set$ is:
\[
  |S| = \Lan_{\yo} (\Delta^\bullet)(S) = \int^{[n]} S_n \times \Delta^n \quad (\text{a coend/colimit formula}).
\]
This is the left Kan extension of $\Delta^\bullet: \Delta \to \Top$ along the Yoneda embedding $\yo: \Delta \to [\Delta\op, \Set]$:
\[
  \begin{tikzcd}
    \Delta \arrow[r, "\yo"] \arrow[dr, "\Delta^\bullet"'] & {[\Delta\op, \Set]} \arrow[d, "|-|"] \\
    & \Top
  \end{tikzcd}
\]
The adjunction $|-| \dashv N$ (geometric realization left adjoint to nerve) is the Kan extension adjunction.
\end{example}

\section{Day Convolution}

\begin{theorem}[Day convolution]
Let $(\cat{C}, \otimes, I)$ be a small symmetric monoidal category. The presheaf category $[\cat{C}\op, \Set]$ carries a symmetric monoidal structure given by \textbf{Day convolution}:
\[
  (F * G)(c) = \int^{a,b \in \cat{C}} \cat{C}(c, a \otimes b) \times Fa \times Gb \;=\; \colim_{c \to a \otimes b} Fa \times Gb,
\]
(a coend/left Kan extension formula). This is the left Kan extension of the monoidal product $\otimes: \cat{C} \times \cat{C} \to \cat{C}$ along $\yo \times \yo: \cat{C} \times \cat{C} \to [\cat{C}\op,\Set] \times [\cat{C}\op,\Set]$, composed with $\yo$.

The Yoneda embedding $\yo: \cat{C} \hookrightarrow [\cat{C}\op,\Set]$ is strong monoidal: $\yo(a) * \yo(b) \cong \yo(a \otimes b)$.
\end{theorem}

Day convolution unifies many constructions: formal power series (polynomial rings), Dirichlet series, convolution algebras of functions, and the tensor product of $\mathbf{Ab}$-enriched presheaves.

\section{Global Kan Extension Functors}

\begin{theorem}
Let $K: \cat{C} \to \cat{D}$ be a functor. The precomposition functor
\[
  K^*: [\cat{D}, \cat{E}] \to [\cat{C}, \cat{E}], \qquad G \mapsto G \circ K,
\]
has a left adjoint $\Lan_K: [\cat{C},\cat{E}] \to [\cat{D},\cat{E}]$ (global left Kan extension) and a right adjoint $\Ran_K: [\cat{C},\cat{E}] \to [\cat{D},\cat{E}]$ (global right Kan extension), whenever $\cat{E}$ is cocomplete and complete respectively.

Symbolically: $\Lan_K \dashv K^* \dashv \Ran_K$.
\end{theorem}

\begin{proof}
For $\Lan_K$: the natural bijection $[\cat{D},\cat{E}]({\Lan_K F}, G) \cong [\cat{C},\cat{E}](F, K^* G) = \Nat(F, GK)$ is the defining universal property of $\Lan_K F$. For $\Ran_K$: similar, with $\Nat(K^*G, F) = \Nat(GK, F) \cong [\cat{D},\cat{E}](G, \Ran_K F)$.
\end{proof}

\begin{corollary}
When $K = \yo: \cat{C} \to [\cat{C}\op,\Set]$ is the Yoneda embedding, $\Lan_\yo: [[\cat{C}\op,\Set], \cat{E}] \to [\cat{C},\cat{E}]$ realizes the free cocompletion universal property: every functor $F: \cat{C} \to \cat{E}$ (with $\cat{E}$ cocomplete) extends uniquely to a cocontinuous $\bar F: [\cat{C}\op,\Set] \to \cat{E}$, and this extension is $\Lan_\yo F$.
\end{corollary}

\section{Absolute Kan Extensions}

\begin{definition}[Absolute Kan extension]
A left Kan extension $(\Lan_K F, \eta)$ is \textbf{absolute} if it is preserved by all functors $H: \cat{E} \to \cat{E}'$: i.e., $H \circ \Lan_K F \cong \Lan_K (H \circ F)$ (with the obvious unit).
\end{definition}

\begin{example}
If $F \dashv G$, then $\Ran_F \Id_{\cat{C}} \cong G$ (Theorem~\ref{thm:adj-as-kan}), and this Kan extension is absolute. Absolute Kan extensions are precisely those computed without using (co)limits—they are ``preserved for free'' by all functors.
\end{example}

\begin{theorem}
A Kan extension is absolute if and only if it is computed by the counit (resp.\ unit) of an adjunction.
\end{theorem}

\section{Kan Extensions in 2-Categories}

The notion of Kan extension generalizes to any 2-category.

\begin{definition}[Kan extension in a 2-category]
In a 2-category $\mathcal{K}$, given 1-cells $K: C \to D$ and $F: C \to E$, a \textbf{left Kan extension} is a pair $(\Lan_K F: D \to E, \eta: F \Rightarrow (\Lan_K F) \circ K)$ satisfying the obvious 2-categorical universal property.
\end{definition}

In $\Cat$, this recovers the usual Kan extensions. In other 2-categories (e.g., the 2-category of rings, bimodules, bimodule maps), it gives Morita-theoretic Kan extensions.

\section{Coends and the Co-Yoneda Lemma}

Kan extensions are intimately connected with coends, a useful computational tool.

\begin{definition}[Coend]
Let $T: \cat{C}\op \times \cat{C} \to \cat{E}$ be a functor (a \textbf{dinatural} situation). The \textbf{coend} $\int^c T(c,c)$ is the colimit of the diagram indexed by the \textbf{twisted arrow category} of $\cat{C}$: objects are morphisms $f: c \to c'$ in $\cat{C}$, with maps $T(f, \id_c): T(c',c) \to T(c,c)$ and $T(\id_{c'}, f): T(c',c) \to T(c',c')$.

Equivalently, $\int^c T(c,c) = \colim_{f: c \to c' \text{ in } \cat{C}} \bigl(T(c',c) \rightrightarrows T(c,c) \sqcup T(c',c')\bigr)$ (coequalizer of action on morphisms).
\end{definition}

\begin{theorem}[Co-Yoneda lemma]
For $F: \cat{C} \to \cat{E}$ and $H \in [\cat{C}\op, \Set]$:
\[
  \int^c Hc \cdot Fc \;\cong\; \Lan_\yo F (H),
\]
where $Hc \cdot Fc$ denotes the \textbf{tensor} (copower): $\bigsqcup_{h \in Hc} Fc$ (or $Hc \times Fc$ in $\Set$).

In particular, for $H = \yo(d) = \cat{C}(-, d)$:
\[
  \int^c \cat{C}(c, d) \cdot Fc \;\cong\; Fd.
\]
\end{theorem}

This ``co-Yoneda'' formula is dual to the Yoneda lemma and shows that any functor value $Fd$ can be reconstructed as a coend.

\section{Exercises}

\begin{exercise}
Show that the right Kan extension along the unique functor $!: \cat{C} \to \mathbf{1}$ is the limit: $\Ran_! F \cong \lim F$. Dualize for left Kan extensions and colimits.
\end{exercise}

\begin{exercise}
Verify that $\Ran_{\Id} F \cong F$ (the right Kan extension along the identity is $F$ itself).
\end{exercise}

\begin{exercise}
Show that if $K: \cat{C} \to \cat{D}$ is fully faithful, then the unit $\eta: F \Rightarrow (\Lan_K F) \circ K$ is a natural isomorphism (the Kan extension truly extends $F$ in this case).
\end{exercise}

\begin{exercise}
Use the pointwise formula to compute $\Lan_K F$ when $K: \{*\} \to \cat{D}$ picks out a single object $d \in \cat{D}$ and $F: \{*\} \to \cat{E}$ picks out $e \in \cat{E}$. What is $(\Lan_K F)(d')$ for varying $d' \in \cat{D}$?
\end{exercise}

\begin{exercise}[$\star$]
Prove the density theorem (Theorem~\ref{thm:density}) for the Yoneda embedding: every presheaf $F \in [\cat{C}\op,\Set]$ is a colimit of representables, with the colimit indexed by the category of elements of $F$.
\end{exercise}

\begin{exercise}[$\star$]
Using Day convolution, describe the monoidal structure on $[\mathbb{N}, \Set]$ (presheaves on the monoid $(\mathbb{N}, +, 0)$ viewed as a one-object category). How does this relate to formal power series?
\end{exercise}

\begin{exercise}[$\star\star$]
Prove that for any adjunction $F \dashv G$, the right Kan extension $\Ran_F H \cong HG$ for any $H: \cat{D} \to \cat{E}$. (This is the ``absolute'' Kan extension claim.) What is the counit of this Kan extension?
\end{exercise}

\begin{exercise}[$\star\star$]
Let $i: \cat{C} \hookrightarrow \cat{D}$ be a full subcategory inclusion with $\cat{D}$ cocomplete. Show that the left Kan extension $\Lan_i F$ for $F: \cat{C} \to \cat{E}$ gives the ``best approximation'' to $F$ on $\cat{D}$, and that if $F$ has a left adjoint on $\cat{C}$, the Kan extension has a left adjoint on $\cat{D}$ (under mild conditions).
\end{exercise}

\chapter{Enriched Category Theory}
\label{ch:enriched}

\section{Motivation}

In an ordinary locally small category, hom-sets $\cat{C}(a,b)$ are sets. But in many mathematical contexts, hom-sets carry additional structure: they are abelian groups (additive categories), vector spaces (linear categories), topological spaces (topological categories), or simplicial sets (simplicial categories). Enriched category theory provides a unified framework for all these situations.

The key idea: replace $\Set$ as the ``ground category'' for hom-sets with a more structured monoidal category $\cat{V}$. Instead of hom-sets, we have hom-\emph{objects} in $\cat{V}$.

\section{Monoidal Categories}

\begin{definition}[Monoidal category]
A \textbf{monoidal category} is a category $\cat{V}$ equipped with:
\begin{itemize}
  \item A \textbf{tensor product} bifunctor $\otimes: \cat{V} \times \cat{V} \to \cat{V}$.
  \item A \textbf{unit object} $I \in \cat{V}$.
  \item Natural isomorphisms:
  \begin{itemize}
    \item \textbf{Associator}: $\alpha_{A,B,C}: (A \otimes B) \otimes C \xrightarrow{\sim} A \otimes (B \otimes C)$.
    \item \textbf{Left unitor}: $\lambda_A: I \otimes A \xrightarrow{\sim} A$.
    \item \textbf{Right unitor}: $\rho_A: A \otimes I \xrightarrow{\sim} A$.
  \end{itemize}
\end{itemize}
satisfying the pentagon axiom (coherence for the associator) and the triangle axiom (coherence for the unitors).

$\cat{V}$ is \textbf{symmetric monoidal} if it additionally has a natural isomorphism $\sigma_{A,B}: A \otimes B \xrightarrow{\sim} B \otimes A$ satisfying the hexagon axioms and $\sigma_{B,A} \circ \sigma_{A,B} = \id$.
\end{definition}

\begin{theorem}[Mac Lane's Coherence Theorem]
In any monoidal category, all diagrams built from associators and unitors commute. Thus we may treat monoidal categories as if they were strictly associative and unital (up to coherent isomorphism).
\end{theorem}

\begin{example}[Standard monoidal categories for enrichment]
\leavevmode
\begin{itemize}
  \item $(\Set, \times, \{*\})$: ordinary categories.
  \item $(\Ab, \otimes_\mathbb{Z}, \mathbb{Z})$: preadditive ($\Ab$-enriched) categories.
  \item $(\Vect_k, \otimes_k, k)$: $k$-linear categories.
  \item $(\mathbf{Cat}, \times, \mathbf{1})$: strict 2-categories.
  \item $(\mathbf{sSet}, \times, \Delta^0)$: simplicially enriched categories.
  \item $([0,\infty], +, 0)$ (Lawvere): extended metric spaces.
  \item $(\{T, F\}, \wedge, T)$ (the Boolean poset): preorders.
\end{itemize}
\end{example}

\section{$\cat{V}$-Enriched Categories}

\begin{definition}[$\cat{V}$-enriched category]
Let $(\cat{V}, \otimes, I)$ be a monoidal category. A \textbf{$\cat{V}$-enriched category} (or \textbf{$\cat{V}$-category}) $\cat{C}$ consists of:
\begin{itemize}
  \item A collection $\ob(\cat{C})$ of objects.
  \item For each pair $a, b \in \ob(\cat{C})$, an object $\cat{C}(a,b) \in \cat{V}$ (the hom-object).
  \item For each triple $a,b,c$, a \textbf{composition morphism} in $\cat{V}$:
  \[
    \circ_{a,b,c}: \cat{C}(b,c) \otimes \cat{C}(a,b) \to \cat{C}(a,c).
  \]
  \item For each object $a$, an \textbf{identity morphism} in $\cat{V}$:
  \[
    j_a: I \to \cat{C}(a,a).
  \]
\end{itemize}
subject to the usual associativity and unit axioms (expressed as commutative diagrams in $\cat{V}$).
\end{definition}

\begin{remark}
An ordinary (locally small) category is precisely a $\Set$-enriched category. The hom-object is the hom-set; composition and identities are functions.
\end{remark}

\begin{example}[$\Ab$-enriched categories (preadditive)]
An $\Ab$-enriched category has abelian groups as hom-objects and bilinear composition. Examples: $\Ab$ itself (as a $\Ab$-enriched category), $\Vect_k$, $\mathbf{Mod}_R$. Any additive category is an $\Ab$-enriched category with a zero object, finite products, and finite coproducts (which coincide).
\end{example}

\begin{example}[Lawvere metric spaces]
A $([0,\infty], +, 0)$-enriched category is precisely a (generalized) metric space in Lawvere's sense: objects are points; $\cat{C}(a,b) \in [0,\infty]$ is the distance; the composition axiom is the triangle inequality $d(a,c) \leq d(b,c) + d(a,b)$; the identity axiom is $d(a,a) \leq 0$, i.e., $d(a,a) = 0$. Non-symmetry and non-Hausdorff metric spaces are naturally captured. Enriched functors are contractive maps; enriched natural transformations are non-expansive morphisms.
\end{example}

\begin{example}[2-categories as $\Cat$-enriched categories]
A $\Cat$-enriched category is a \textbf{2-category}: its hom-objects are categories (whose objects are the 1-cells and whose morphisms are the 2-cells), and composition is a functor (giving horizontal composition of 2-cells).
\end{example}

\begin{example}[Simplicial categories]
A $\mathbf{sSet}$-enriched category (simplicial category) has simplicial sets as hom-objects. These arise in homotopy theory and higher category theory: the mapping spaces encode homotopy types of spaces of morphisms.
\end{example}

\section{$\cat{V}$-Functors and $\cat{V}$-Natural Transformations}

\begin{definition}[$\cat{V}$-functor]
Let $\cat{C}$ and $\cat{D}$ be $\cat{V}$-categories. A \textbf{$\cat{V}$-functor} $F: \cat{C} \to \cat{D}$ consists of:
\begin{itemize}
  \item A function $F: \ob(\cat{C}) \to \ob(\cat{D})$.
  \item For each $a, b \in \cat{C}$, a morphism in $\cat{V}$: $F_{a,b}: \cat{C}(a,b) \to \cat{D}(Fa,Fb)$.
\end{itemize}
preserving composition and identities (as commutative diagrams in $\cat{V}$).
\end{definition}

\begin{definition}[$\cat{V}$-natural transformation]
A \textbf{$\cat{V}$-natural transformation} $\alpha: F \Rightarrow G$ (between $\cat{V}$-functors $F, G: \cat{C} \to \cat{D}$) consists of morphisms $\alpha_a: I \to \cat{D}(Fa, Ga)$ in $\cat{V}$ (``elements'' of the hom-objects) satisfying the enriched naturality condition:
\[
  \begin{tikzcd}
    \cat{C}(a,b) \arrow[r, "F_{a,b}"] \arrow[d, "G_{a,b}"'] & \cat{D}(Fa,Fb) \arrow[d, "(\alpha_b)_*"] \\
    \cat{D}(Ga,Gb) \arrow[r, "(\alpha_a)^*"'] & \cat{D}(Fa,Gb)
  \end{tikzcd}
\]
(where $(\alpha_b)_*$ and $(\alpha_a)^*$ are composition with $\alpha_b$ and $\alpha_a$, using the composition morphisms of $\cat{D}$).
\end{definition}

\begin{remark}
When $\cat{V} = \Set$, a $\cat{V}$-natural transformation reduces to an ordinary natural transformation (since $\alpha_a: \{*\} \to \cat{D}(Fa, Ga)$ picks out a single element of the hom-set).
\end{remark}

$\cat{V}$-categories, $\cat{V}$-functors, and $\cat{V}$-natural transformations form a 2-category $\cat{V}$-$\Cat$.

\section{The Enriched Yoneda Lemma}

\begin{theorem}[Enriched Yoneda Lemma]
Let $\cat{C}$ be a $\cat{V}$-category and $F: \cat{C} \to \cat{V}$ a $\cat{V}$-functor (where $\cat{V}$ is enriched over itself via its closed structure, when $\cat{V}$ is closed monoidal). Then:
\[
  [\cat{C}, \cat{V}](\cat{C}(a,-), F) \;\cong\; Fa
\]
as objects of $\cat{V}$, naturally in $a \in \cat{C}$ and $F \in [\cat{C},\cat{V}]$.
\end{theorem}

The enriched Yoneda lemma requires $\cat{V}$ to be \textbf{closed monoidal}: for each $A \in \cat{V}$, the functor $A \otimes -$ has a right adjoint $[A, -]$ (the \textbf{internal hom}). Key examples:
\begin{itemize}
  \item $(\Set, \times)$: internal hom is $[A,B] = B^A$ (function set).
  \item $(\Ab, \otimes_\mathbb{Z})$: internal hom is $\Hom_\mathbb{Z}(A,B)$ (abelian group of homomorphisms).
  \item $(\Vect_k, \otimes_k)$: internal hom is $\Hom_k(V, W)$.
  \item $(\mathbf{sSet}, \times)$: internal hom is the simplicial mapping space $\mathrm{Map}(X,Y)$.
\end{itemize}

\section{Enriched Limits: Weighted Limits}

Ordinary limits are controlled by shape categories $\cat{J}$. In the enriched setting, the shape data includes a \textbf{weight} functor.

\begin{definition}[Weighted limit]
Let $\cat{C}$ be a $\cat{V}$-category, $D: \cat{J} \to \cat{C}$ a $\cat{V}$-functor, and $W: \cat{J} \to \cat{V}$ a weight functor. The \textbf{$W$-weighted limit} of $D$ is an object $\{W, D\} \in \cat{C}$ representing the $\cat{V}$-functor:
\[
  c \;\mapsto\; [\cat{J}, \cat{V}](W, \cat{C}(c, D-)): \cat{C}\op \to \cat{V}.
\]
That is, $\cat{C}(c, \{W,D\}) \cong [\cat{J},\cat{V}](W, \cat{C}(c, D-))$ in $\cat{V}$, naturally in $c$.
\end{definition}

\begin{remark}
Ordinary (conical) limits correspond to the constant weight $W = \Delta_I$ (the constant functor at the monoidal unit $I$). Weighted limits strictly generalize ordinary limits and are essential in enriched category theory, where conical limits are often insufficient.
\end{remark}

\begin{example}[Ends as weighted limits]
The \textbf{end} $\int_c T(c,c)$ of a $\cat{V}$-functor $T: \cat{C}\op \otimes \cat{C} \to \cat{V}$ (where $\cat{C}\op \otimes \cat{C}$ is the tensor product of $\cat{V}$-categories) is the limit of $T$ weighted by the hom-functor $\cat{C}(-,-): \cat{C}\op \otimes \cat{C} \to \cat{V}$.

The $\cat{V}$-valued natural transformations $\Nat(F,G)$ are an end:
\[
  \Nat(F,G) = \int_c \cat{D}(Fc, Gc).
\]
\end{example}

\section{Change of Base}

\begin{definition}[Change of base]
A \textbf{lax monoidal functor} $\Phi: \cat{V} \to \cat{W}$ (with coherence maps $\phi_{A,B}: \Phi A \otimes_{\cat{W}} \Phi B \to \Phi(A \otimes_{\cat{V}} B)$ and $\phi_I: I_{\cat{W}} \to \Phi(I_{\cat{V}})$) induces a \textbf{change of base} 2-functor $\Phi_*: \cat{V}\text{-}\Cat \to \cat{W}\text{-}\Cat$: apply $\Phi$ to all hom-objects and compose the structural maps.
\end{definition}

\begin{example}
The ``set of connected components'' functor $\pi_0: \mathbf{sSet} \to \Set$ is lax monoidal. Change of base along $\pi_0$ sends a simplicial category to its underlying ordinary category (obtained by taking connected components of each hom-space).

The free abelian group functor $\mathbb{Z}[-]: \Set \to \Ab$ is lax monoidal. Change of base sends ordinary categories to $\Ab$-enriched categories (the ``abelianization'' of a category).
\end{example}

\section{Exercises}

\begin{exercise}
Verify that an $[0,\infty]$-enriched category is precisely a generalized metric space (in the sense of Lawvere): a set $X$ with a function $d: X \times X \to [0,\infty]$ satisfying $d(x,x) = 0$ and the triangle inequality.
\end{exercise}

\begin{exercise}
Show that in any $\Ab$-enriched category, for any two objects $a, b$, the hom-group $\cat{C}(a,b)$ carries an abelian group structure compatible with composition (i.e., composition is bilinear).
\end{exercise}

\begin{exercise}
Define the $\cat{V}$-enriched functor category $[\cat{C},\cat{D}]$ for $\cat{V}$-categories $\cat{C}$ and $\cat{D}$ (assuming $\cat{V}$ is closed monoidal). What are its hom-objects?
\end{exercise}

\begin{exercise}[$\star$]
Let $\cat{C}$ be a $\cat{V}$-category with $\cat{V} = \Ab$. Define the notion of a $\cat{V}$-functor $F: \cat{C} \to \Ab$ and show that $\Ab$-natural transformations between such functors form an abelian group.
\end{exercise}

\begin{exercise}[$\star$]
Prove that for a small $\Ab$-enriched category $\cat{C}$, the functor category $[\cat{C}\op, \Ab]$ is an abelian category (assuming familiarity with abelian categories). Identify the representable functors with projective objects.
\end{exercise}

\begin{exercise}[$\star\star$]
Formulate and prove the enriched Yoneda lemma for $\cat{V}$-functors $F: \cat{C} \to \cat{V}$ where $\cat{V}$ is a symmetric closed monoidal category. Show how it reduces to the ordinary Yoneda lemma when $\cat{V} = \Set$.
\end{exercise}

\chapter{Topos Theory}
\label{ch:topos}

\section{Introduction}

A topos is a category that behaves like the category $\Set$ of sets, but in a generalized and often constructive sense. Toposes arise naturally as categories of sheaves on topological spaces (or more general sites), as the classifying categories of geometric theories, and as models of constructive set theory.

The theory of toposes was initiated by Grothendieck in the context of algebraic geometry (as ``Grothendieck toposes'') and developed by Lawvere and Tierney into elementary topos theory, which uses no large cardinal axioms and has deep connections to logic and type theory.

\section{Subobject Classifiers}

The most striking feature distinguishing $\Set$ from an arbitrary category is the existence of a \emph{truth-value object} $\Omega$: a set $\Omega = \{T, F\}$ such that every subset $S \subseteq A$ is classified by its characteristic function $\chi_S: A \to \Omega$.

\begin{definition}[Subobject classifier]
In a category $\cat{C}$ with a terminal object $1$, a \textbf{subobject classifier} is an object $\Omega$ together with a morphism $\mathrm{true}: 1 \to \Omega$ such that for every monomorphism $m: S \hookrightarrow A$ in $\cat{C}$, there is a unique \textbf{characteristic morphism} $\chi_m: A \to \Omega$ making the square
\[
  \begin{tikzcd}
    S \arrow[r] \arrow[d, "m"', hook] & 1 \arrow[d, "\mathrm{true}"] \\
    A \arrow[r, "\chi_m"'] & \Omega
  \end{tikzcd}
\]
a pullback. In other words, monomorphisms (subobjects) $m: S \hookrightarrow A$ are in natural bijection with morphisms $A \to \Omega$.
\end{definition}

\begin{example}
In $\Set$: $\Omega = \{0,1\}$ (or $\{F,T\}$), $\mathrm{true} = 1 \mapsto 1$. The characteristic function of $S \subseteq A$ is $\chi_S(a) = 1$ iff $a \in S$.
\end{example}

\begin{example}
In the presheaf category $[\cat{C}\op, \Set]$: $\Omega$ is the presheaf of \textbf{sieves}, $\Omega(c) = \{S \mid S \text{ is a sieve on } c\}$ (a sieve on $c$ is a set of morphisms with codomain $c$ closed under precomposition). This is a generalized truth-value object.
\end{example}

\section{Elementary Toposes}

\begin{definition}[Elementary topos]
An \textbf{elementary topos} (or simply \textbf{topos}) is a category $\mathcal{E}$ that:
\begin{enumerate}
  \item Has all finite limits.
  \item Has all finite colimits. (This follows from the other axioms.)
  \item Has a subobject classifier $\Omega$.
  \item Has \textbf{power objects}: for every object $B \in \mathcal{E}$, there is an object $\Omega^B$ (the ``power object'' of $B$) and a natural bijection
  \[
    \mathcal{E}(A \times B, \Omega) \;\cong\; \mathcal{E}(A, \Omega^B).
  \]
  Equivalently, $\mathcal{E}$ is \textbf{locally cartesian closed}: the functor $- \times B: \mathcal{E} \to \mathcal{E}$ has a right adjoint $(-)^B$ for every $B$.
\end{enumerate}
\end{definition}

\begin{remark}
The power object axiom says $\mathcal{E}$ is \textbf{cartesian closed}: it has exponential objects $B^A = \mathcal{E}(A, B)$ representing the functor $\mathcal{E}(- \times A, B)$.

The combination of finite limits, cartesian closure, and a subobject classifier is remarkably powerful: it suffices to interpret all of higher-order intuitionistic logic.
\end{remark}

\begin{theorem}[Basic consequences]
In any topos $\mathcal{E}$:
\begin{enumerate}
  \item $\mathcal{E}$ has all finite colimits.
  \item $\mathcal{E}$ is locally cartesian closed: every slice category $\mathcal{E}/A$ is itself a topos.
  \item The subobject functor $\mathrm{Sub}: \mathcal{E}\op \to \Set$ is representable by $\Omega$ (via: $\mathrm{Sub}(A) \cong \mathcal{E}(A, \Omega)$).
  \item $\mathrm{Sub}(A)$ carries a Heyting algebra structure (the internal logic is intuitionistic).
\end{enumerate}
\end{theorem}

\section{Examples of Toposes}

\begin{example}[$\Set$]
The category of sets $\Set$ is the archetypal topos. Its subobject classifier is $\Omega = \{0,1\}$ and its internal logic is classical Boolean logic.
\end{example}

\begin{example}[Presheaf toposes]
For any small category $\cat{C}$, the presheaf category $\widehat{\cat{C}} = [\cat{C}\op, \Set]$ is a topos. These are called \textbf{presheaf toposes} (or Grothendieck toposes over a trivial topology). The subobject classifier $\Omega(c)$ is the set of sieves on $c$ in $\cat{C}$; finite limits and colimits are pointwise; exponentials are computed by the formula $F^G(c) = [\cat{C}\op/c, \Set](G|_{c}, F|_{c})$.

Special cases:
\begin{itemize}
  \item $\cat{C} = \mathbf{1}$: $[\mathbf{1}\op, \Set] \cong \Set$.
  \item $\cat{C} = BG$ (a group): $[BG\op, \Set] \cong G$-$\Set$ (category of $G$-sets).
  \item $\cat{C} = \Delta\op$ (simplices): $[\Delta, \Set]$ is the category of simplicial sets.
\end{itemize}
\end{example}

\begin{example}[Sheaf toposes]
Let $X$ be a topological space. The category $\mathbf{Sh}(X)$ of sheaves of sets on $X$ is a topos. The subobject classifier is the sheaf $\Omega$ with $\Omega(U) = \{V \mid V \subseteq U, V \text{ open}\}$ (open subsets). The internal logic of $\mathbf{Sh}(X)$ reflects the topology of $X$: classical logic holds iff $X$ is a discrete space; Heyting (intuitionistic) logic holds in general.
\end{example}

\begin{example}[The effective topos]
The \textbf{effective topos} (Hyland 1982) is a topos whose internal logic is the logic of recursive realizability. It provides a categorical model of Kleene's realizability interpretation and contains a natural number object with the property that all its functions are recursive. This topos cannot be a Grothendieck topos.
\end{example}

\section{The Internal Language of a Topos}

Every topos has an \textbf{internal language}: a formal system (Mitchell-Bénabou language) in which one can reason about objects of the topos using logical language, with the topos playing the role of a model.

\begin{definition}[Mitchell-Bénabou language]
Given a topos $\mathcal{E}$, the \textbf{internal language} $L(\mathcal{E})$ is a typed higher-order language where:
\begin{itemize}
  \item \textbf{Types} correspond to objects of $\mathcal{E}$.
  \item \textbf{Terms of type $A$ in context $\Gamma$} correspond to morphisms $\|\Gamma\| \to A$ in $\mathcal{E}$.
  \item \textbf{Propositions (of type $\Omega$)} correspond to morphisms $\|\Gamma\| \to \Omega$.
  \item \textbf{Logical connectives} ($\wedge, \vee, \Rightarrow, \neg, \forall, \exists$) correspond to morphisms constructed from $\Omega, \times, +, (-)^\Omega$, etc.
  \item \textbf{Subsets} $\{x: A \mid \phi(x)\}$ correspond to pullbacks along $\phi: A \to \Omega$ of $\mathrm{true}: 1 \to \Omega$.
\end{itemize}
\end{definition}

\begin{theorem}[Soundness and completeness]
A formula is provable in the internal language (using intuitionistic logic) if and only if it holds in every topos. The soundness theorem says that valid constructions in the internal language correspond to morphisms in $\mathcal{E}$.
\end{theorem}

\begin{example}
In $\mathbf{Sh}(X)$, the internal statement ``every function $f: \mathbb{R} \to \mathbb{R}$ is continuous'' is satisfied (where $\mathbb{R}$ is the sheaf of continuous real-valued functions). This reflects the fact that the internal real numbers of $\mathbf{Sh}(X)$ are continuous functions.
\end{example}

\section{Heyting Algebras and Intuitionistic Logic}

\begin{definition}[Heyting algebra]
A \textbf{Heyting algebra} is a lattice $H$ with a binary operation $\Rightarrow$ (implication) satisfying $a \leq b \Rightarrow c$ iff $a \wedge b \leq c$ (the residuation property). Equivalently, it is a cartesian closed poset.

A \textbf{Boolean algebra} is a Heyting algebra satisfying $a \vee \neg a = 1$ (law of excluded middle), where $\neg a = a \Rightarrow 0$.
\end{definition}

\begin{theorem}
In any topos $\mathcal{E}$, the subobject lattice $\mathrm{Sub}(A)$ is a Heyting algebra for every $A \in \mathcal{E}$. The topos has Boolean logic (satisfies the law of excluded middle internally) if and only if $\Omega \cong \Omega + \Omega$ (the subobject classifier decomposes as a coproduct of two copies of the terminal object).
\end{theorem}

\begin{corollary}
$\Set$ is a Boolean topos. Sheaf toposes $\mathbf{Sh}(X)$ for non-discrete $X$ are non-Boolean. The effective topos is non-Boolean.
\end{corollary}

\section{Geometric Morphisms}

\begin{definition}[Geometric morphism]
Let $\mathcal{E}$ and $\mathcal{F}$ be toposes. A \textbf{geometric morphism} $f: \mathcal{E} \to \mathcal{F}$ is an adjunction $f^* \dashv f_*$ where:
\begin{itemize}
  \item $f^*: \mathcal{F} \to \mathcal{E}$ is the \textbf{inverse image functor} (left adjoint).
  \item $f_*: \mathcal{E} \to \mathcal{F}$ is the \textbf{direct image functor} (right adjoint).
\end{itemize}
and $f^*$ preserves finite limits (is left exact).
\end{definition}

\begin{remark}
The condition that $f^*$ is left exact (preserves finite limits) is the key additional condition that distinguishes geometric morphisms from arbitrary adjunctions.
\end{remark}

\begin{example}
A continuous map $g: X \to Y$ of topological spaces induces a geometric morphism $g: \mathbf{Sh}(X) \to \mathbf{Sh}(Y)$ via:
\begin{itemize}
  \item $g^*(\mathcal{F})(U) = \colim_{V \supseteq g(U), V \text{ open}} \mathcal{F}(V)$ (sheaf-theoretic inverse image).
  \item $g_*(\mathcal{G})(V) = \mathcal{G}(g^{-1}(V))$ (direct image).
\end{itemize}
\end{example}

\begin{example}[The global sections geometric morphism]
For any topos $\mathcal{E}$, there is a unique geometric morphism $\Gamma: \mathcal{E} \to \Set$ (the \textbf{global sections morphism}):
\begin{itemize}
  \item $\Gamma^*: \Set \to \mathcal{E}$ is the \textbf{constant sheaf functor} (sends a set to the constant presheaf, sheafified if necessary).
  \item $\Gamma_*: \mathcal{E} \to \Set$, $\mathcal{F} \mapsto \mathcal{E}(1, \mathcal{F})$ is the \textbf{global sections functor}.
\end{itemize}
\end{example}

Geometric morphisms are the correct notion of ``morphism'' between Grothendieck toposes, and they form the objects and morphisms of the 2-category of toposes.

\section{Lawvere-Tierney Topologies}

Inside a topos, one can define an internal notion of ``coverage'' or ``topology'' using the subobject classifier.

\begin{definition}[Lawvere-Tierney topology]
A \textbf{Lawvere-Tierney topology} (or \textbf{local operator}) on a topos $\mathcal{E}$ is a morphism $j: \Omega \to \Omega$ satisfying:
\begin{enumerate}
  \item $j \circ \mathrm{true} = \mathrm{true}$ (truth is preserved).
  \item $j \circ j = j$ (idempotence).
  \item $j \circ \wedge = \wedge \circ (j \times j)$ (meets are preserved).
\end{enumerate}
where $\wedge: \Omega \times \Omega \to \Omega$ is the conjunction morphism.
\end{definition}

\begin{theorem}
Given a Lawvere-Tierney topology $j$ on $\mathcal{E}$, the full subcategory of \textbf{$j$-sheaves} (objects $F$ for which $\mathcal{E}(-, F)$ inverts $j$-dense monomorphisms) is itself a topos $\mathbf{Sh}_j(\mathcal{E})$, and the inclusion $\mathbf{Sh}_j(\mathcal{E}) \hookrightarrow \mathcal{E}$ has a left adjoint (the \textbf{$j$-sheafification functor}), making it a geometric morphism.
\end{theorem}

\begin{example}
The double-negation topology $j = \neg\neg: \Omega \to \Omega$ on $\Set$ (or any Boolean topos) is trivial—every presheaf is a sheaf. For presheaf toposes $[\cat{C}\op,\Set]$, the double-negation sheaves form a Boolean topos.
\end{example}

\section{Classifying Toposes}

One of the most profound aspects of topos theory is the correspondence between geometric theories and classifying toposes.

\begin{definition}[Geometric theory]
A \textbf{geometric theory} $\mathbb{T}$ is a theory (in the sense of logic) whose axioms can be expressed using $\top$, $\bot$, $\wedge$ (finite conjunctions), $\vee$ (arbitrary disjunctions), $\exists$ (existential quantification), and $=$ (equality)—but no $\forall$ or $\Rightarrow$.
\end{definition}

\begin{theorem}[Classifying topos]
For any geometric theory $\mathbb{T}$, there exists a topos $\mathbf{Set}[\mathbb{T}]$ (the \textbf{classifying topos} of $\mathbb{T}$) such that for any Grothendieck topos $\mathcal{E}$,
\[
  \{\text{geometric morphisms } \mathcal{E} \to \mathbf{Set}[\mathbb{T}]\} \;\cong\; \{\mathbb{T}\text{-models in } \mathcal{E}\},
\]
naturally in $\mathcal{E}$. The topos $\mathbf{Set}[\mathbb{T}]$ ``classifies'' models of $\mathbb{T}$.
\end{theorem}

\begin{example}
\begin{itemize}
  \item The classifying topos of the theory of objects (no axioms) is $[\mathbf{FinSet}\op, \Set]$ (presheaves on finite sets).
  \item The classifying topos of the theory of local rings is the Zariski topos, relevant to algebraic geometry.
  \item The classifying topos of the theory of linear orders is related to the simplex category $\Delta$.
\end{itemize}
\end{example}

\section{Connections to Logic and Type Theory}

The \textbf{Curry-Howard-Lambek correspondence} unifies:
\begin{center}
\begin{tabular}{ccc}
\toprule
Logic & Type Theory & Category Theory \\
\midrule
Proposition & Type & Object \\
Proof & Term & Morphism \\
Conjunction & Product type & Product \\
Disjunction & Sum type & Coproduct \\
Implication & Function type & Exponential \\
True / False & Unit / Empty & Terminal / Initial \\
Universal & Dependent product $\Pi$ & Right adjoint \\
Existential & Dependent sum $\Sigma$ & Left adjoint \\
\bottomrule
\end{tabular}
\end{center}

\begin{theorem}
The internal language of a topos (equipped with a natural number object) models Martin-Löf dependent type theory (in its impredicative form). Conversely, the category of types and terms of a suitable dependent type theory forms a topos.
\end{theorem}

This correspondence is the foundation of \textbf{homotopy type theory} (HoTT) and \textbf{univalent foundations}, where the category-theoretic perspective (via $(\infty,1)$-toposes) provides semantic models for the type theory.

\section{Exercises}

\begin{exercise}
Show that the subobject classifier $\Omega$ in $[\cat{C}\op, \Set]$ is the functor sending each object $c \in \cat{C}$ to the set of sieves on $c$. Verify the universal property: subobjects of $F$ correspond to natural transformations $F \to \Omega$.
\end{exercise}

\begin{exercise}
Show that in any topos, the subobject lattice $\mathrm{Sub}(A)$ has joins: the join $m_1 \vee m_2$ of two subobjects $m_1: S_1 \hookrightarrow A$ and $m_2: S_2 \hookrightarrow A$ is the image of the induced morphism $S_1 \sqcup S_2 \to A$.
\end{exercise}

\begin{exercise}
Show that every topos $\mathcal{E}$ has a natural number object $N$ (an initial algebra for the functor $1 + -$), and that $N$ supports primitive recursion.
\end{exercise}

\begin{exercise}[$\star$]
Describe the subobject classifier for the topos of $G$-sets $[BG\op, \Set]$ (for a group $G$). What are the ``truth values'' in this topos? (They are the subgroups of $G$.)
\end{exercise}

\begin{exercise}[$\star$]
Prove that a geometric morphism $f: \mathcal{E} \to \mathcal{F}$ with $f^*$ exact (preserving finite limits) induces a morphism of Heyting algebras $\Omega_\mathcal{F} \to \Omega_\mathcal{E}$ (via $f^*(\Omega_\mathcal{F}) \cong \Omega_\mathcal{E}$... actually prove that $f^*(\Omega_\mathcal{F})$ is a subobject classifier in $\mathcal{E}$).
\end{exercise}

\begin{exercise}[$\star\star$]
Let $j: \Omega \to \Omega$ be a Lawvere-Tierney topology on a presheaf topos $[\cat{C}\op, \Set]$. Show that $j$ corresponds to a Grothendieck topology on $\cat{C}$, and that the $j$-sheaves are exactly the sheaves for this Grothendieck topology.
\end{exercise}


%% ============================================================
\appendix
%% ============================================================

\chapter{Set-Theoretic Foundations}
\label{app:foundations}

Category theory, like all of mathematics, requires some set-theoretic foundation. The main text uses set theory informally, appealing to intuition and the reader's common sense. This appendix discusses the foundational issues that arise when categories become ``large,'' and presents two standard resolutions: the von Neumann-Bernays-Gödel (NBG) approach using classes, and Grothendieck's universe approach.

\section{Naive Set Theory and Its Limits}

In naive set theory, one works with sets freely, forming them by comprehension: $\{x \mid \phi(x)\}$ for any property $\phi$. Russell's paradox (1901) shows this is inconsistent: let $R = \{x \mid x \notin x\}$; then $R \in R \iff R \notin R$.

Zermelo-Fraenkel set theory with Choice (ZFC) avoids this paradox by restricting comprehension and laying down explicit axioms.

\begin{definition}[The axioms of ZFC (informal)]
\leavevmode
\begin{enumerate}
  \item \textbf{Extensionality.} Sets with the same elements are equal: $\forall z(z \in x \iff z \in y) \Rightarrow x = y$.
  \item \textbf{Empty set.} There exists a set with no elements: $\exists x \forall y (y \notin x)$.
  \item \textbf{Pairing.} For any $a, b$, the set $\{a,b\}$ exists.
  \item \textbf{Union.} For any set $\mathcal{A}$, the union $\bigcup \mathcal{A}$ exists.
  \item \textbf{Power set.} For any set $x$, the power set $\mathcal{P}(x) = \{y \mid y \subseteq x\}$ exists.
  \item \textbf{Separation (Restricted Comprehension).} For any set $x$ and formula $\phi$, $\{y \in x \mid \phi(y)\}$ exists.
  \item \textbf{Replacement.} If $\phi(x,y)$ defines a function on a set, the image is a set.
  \item \textbf{Infinity.} There exists an infinite set (containing $\emptyset$ and closed under $x \mapsto x \cup \{x\}$).
  \item \textbf{Foundation (Regularity).} Every nonempty set $x$ has an element disjoint from $x$. (No infinite descending $\in$-chains.)
  \item \textbf{Choice (AC).} Every family of nonempty sets has a choice function.
\end{enumerate}
\end{definition}

In ZFC, all mathematical objects are encoded as sets. The natural numbers are $0 = \emptyset$, $1 = \{\emptyset\}$, $2 = \{\emptyset, \{\emptyset\}\}$, etc.\ (von Neumann ordinals). Ordered pairs are Kuratowski pairs: $(a,b) = \{\{a\}, \{a,b\}\}$. Functions are sets of ordered pairs.

\section{The Size Problem in Category Theory}

A fundamental problem arises: the collection of all sets is not a set (Russell's paradox applies). Therefore $\Set$ does not have an object set—it is a \emph{large} category. Similarly, $\Grp$, $\Top$, etc.\ are large.

But the main concern for category theory is: can we form the collection of all small categories? Can we form the functor category $[\cat{C}\op, \Set]$ when $\cat{C}$ is small? Can we compose large functors?

Two standard approaches resolve this.

\section{Classes: The NBG and MK Approach}

\textbf{von Neumann-Bernays-Gödel set theory (NBG)} and \textbf{Morse-Kelley set theory (MK)} extend ZFC by a primitive notion of \textbf{class}: a class is any collection, and a class that is itself an element of another class is called a \textbf{set}. Classes that are not sets are \textbf{proper classes}.

\begin{itemize}
  \item In NBG/MK, $V$ (the class of all sets) and $\mathbf{On}$ (the class of all ordinals) are proper classes.
  \item The ``class of all groups'' $\text{ob}(\Grp)$ is a proper class; the hom-set $\Grp(G,H)$ is a set.
  \item A \textbf{large category} has a proper class of objects; a \textbf{small category} has a set of objects (and a set of morphisms). A \textbf{locally small} category has a set of morphisms between any two objects.
\end{itemize}

The drawback: in NBG, you cannot form the ``class of all classes'' (there is no class of all large categories). This means $\mathbf{CAT}$ (the category of all large categories) is not well-defined in NBG.

\section{Universes: The Grothendieck Approach}

Grothendieck introduced a cleaner solution using \textbf{universes}.

\begin{definition}[Grothendieck universe]
A \textbf{Grothendieck universe} is a set $\mathcal{U}$ satisfying:
\begin{enumerate}
  \item $x \in \mathcal{U}$ and $y \in x$ implies $y \in \mathcal{U}$ (transitivity).
  \item $x, y \in \mathcal{U}$ implies $\{x,y\} \in \mathcal{U}$ (pairing).
  \item $x \in \mathcal{U}$ implies $\mathcal{P}(x) \in \mathcal{U}$ (power set).
  \item If $I \in \mathcal{U}$ and $\{x_i\}_{i \in I}$ is a family with all $x_i \in \mathcal{U}$, then $\bigcup_i x_i \in \mathcal{U}$ (union).
  \item $\omega \in \mathcal{U}$ (the natural numbers are in $\mathcal{U}$).
\end{enumerate}
A set $x$ is \textbf{$\mathcal{U}$-small} if $x \in \mathcal{U}$.
\end{definition}

\begin{remark}
A Grothendieck universe is essentially a model of ZFC inside ZFC. The existence of a Grothendieck universe is equivalent (over ZFC) to the existence of an \textbf{inaccessible cardinal}: an uncountable cardinal $\kappa$ that is a limit cardinal and regular ($\text{cf}(\kappa) = \kappa$). The existence of inaccessible cardinals cannot be proved in ZFC (by Gödel's incompleteness theorems).
\end{remark}

\textbf{The universe axiom} (Grothendieck's convention) asserts: every set is contained in some Grothendieck universe.

With this axiom (or with a fixed universe $\mathcal{U}_0 \in \mathcal{U}_1 \in \mathcal{U}_2 \in \cdots$), we can:
\begin{itemize}
  \item Call a category \textbf{small} if its objects and morphisms are $\mathcal{U}_0$-small.
  \item Call a category \textbf{large} (or \textbf{locally small}) if its hom-sets are $\mathcal{U}_0$-small.
  \item Form the ``category of all small categories'' $\Cat$ (which is $\mathcal{U}_1$-small).
  \item Form functor categories $[\cat{C}\op, \Set]$ for small $\cat{C}$, which are locally small.
\end{itemize}

Grothendieck's universe approach is now standard in algebraic geometry and topos theory.

\section{The Axiom of Choice}

The Axiom of Choice (AC) is used extensively in category theory:
\begin{itemize}
  \item Choosing skeleta of categories (one representative from each isomorphism class).
  \item Showing that a fully faithful essentially surjective functor is an equivalence (requires choosing the equivalence inverse).
  \item The General Adjoint Functor Theorem (GAFT).
  \item The existence of enough injectives and projectives.
\end{itemize}

The connection between AC and category theory is deep: many categorical theorems are equivalent to AC or weaker choice principles. In constructive mathematics and intuitionistic logic (as in topos theory), AC is not assumed, and many results require reformulation.

\section{Cardinal Arithmetic}

We collect some basic facts about cardinals used in the text.

\begin{definition}
A \textbf{cardinal} is an ordinal that cannot be put in bijection with any smaller ordinal. For a set $X$, $|X|$ denotes the cardinality (the unique cardinal in bijection with $X$, assuming AC).

\begin{itemize}
  \item $\aleph_0 = |\mathbb{N}|$ (countably infinite).
  \item $\aleph_1$ = the first uncountable cardinal.
  \item $\mathfrak{c} = 2^{\aleph_0} = |\mathbb{R}|$ (the continuum).
\end{itemize}

A set is \textbf{finite} if $|X| < \aleph_0$, \textbf{countable} if $|X| \leq \aleph_0$, \textbf{uncountable} otherwise.
\end{definition}

\begin{theorem}[Cantor]
$|X| < |\mathcal{P}(X)|$ for all sets $X$. In particular, there is no ``largest'' cardinal.
\end{theorem}

\begin{definition}
A cardinal $\kappa$ is \textbf{regular} if it cannot be expressed as a sum $\sum_{i < \lambda} \kappa_i$ with $\lambda < \kappa$ and each $\kappa_i < \kappa$. $\aleph_0$ and every successor cardinal $\aleph_{\alpha+1}$ are regular; limit cardinals may or may not be.

A cardinal $\kappa > \aleph_0$ is \textbf{(strongly) inaccessible} if it is regular and $2^\lambda < \kappa$ for all $\lambda < \kappa$.
\end{definition}

Every Grothendieck universe $\mathcal{U}$ corresponds to an inaccessible cardinal $\kappa = |\mathcal{U}|$, and $\mathcal{U} = V_\kappa$ (the $\kappa$-th stage of the cumulative hierarchy).

\section{Summary of Conventions}

Throughout this book, we use the following conventions:
\begin{itemize}
  \item \textbf{Small category}: objects and morphisms form sets.
  \item \textbf{Locally small category}: each hom-class is a set.
  \item \textbf{Large category}: objects may form a proper class (or a set outside $\mathcal{U}_0$).
  \item We assume a suitable universe axiom so that all functor categories we form are legitimate.
  \item We use the Axiom of Choice without explicit mention.
  \item Set-theoretic subtleties beyond local smallness are noted but not belabored.
\end{itemize}

\chapter{Selected Solutions and Hints}
\label{app:exercises}

This appendix provides hints and selected solutions to the starred exercises from each chapter.

\section*{Chapter 2: Categories}

\begin{exercise}[2.$\star$: Preorders as categories]
A preorder $(P, \leq)$ gives a category $\cat{P}$ with $\cat{P}(a,b) = \emptyset$ or $\{*\}$ according to $a \leq b$ or not. The category axioms correspond exactly to:
\begin{itemize}
  \item Identity: reflexivity ($a \leq a$).
  \item Composition: transitivity (if $a \leq b$ and $b \leq c$, then $a \leq c$).
  \item Uniqueness of the morphism: automatically satisfied since each hom-set has at most one element.
\end{itemize}
A poset additionally satisfies antisymmetry: if $a \leq b$ and $b \leq a$, then $a = b$. Categorically, this means that if $\cat{P}(a,b) \neq \emptyset$ and $\cat{P}(b,a) \neq \emptyset$ (i.e., there is a morphism in each direction), then $a = b$. A preorder differs from a poset in that two distinct objects may be isomorphic.
\end{exercise}

\begin{exercise}[2.$\star$: Arrow category]
Define $\cat{C}^{\to}$ with:
\begin{itemize}
  \item Objects: morphisms $f: a \to b$ in $\cat{C}$.
  \item Morphisms from $f: a \to b$ to $g: c \to d$: pairs $(u: a \to c, v: b \to d)$ in $\cat{C}$ with $v \circ f = g \circ u$.
\end{itemize}

\noindent\textit{Identification as $[\mathbf{2}, \cat{C}]$.} The category $\mathbf{2}$ has two objects $\{0,1\}$ and one non-identity morphism $\iota: 0 \to 1$. A functor $F: \mathbf{2} \to \cat{C}$ is determined by $F(0)$, $F(1)$, and $F(\iota): F(0) \to F(1)$—i.e., by a morphism in $\cat{C}$. A natural transformation $\alpha: F \Rightarrow G$ consists of components $\alpha_0: F(0) \to G(0)$ and $\alpha_1: F(1) \to G(1)$ with $G(\iota) \circ \alpha_0 = \alpha_1 \circ F(\iota)$, which is exactly a morphism in $\cat{C}^{\to}$.
\end{exercise}

\begin{exercise}[2.$\star\star$: Relations]
The category $\mathbf{Rel}$ has:
\begin{itemize}
  \item Objects: sets.
  \item Morphisms $A \to B$: subsets $R \subseteq A \times B$.
  \item Composition: $(R: A \to B, S: B \to C)$ composes to $S \circ R = \{(a,c) \in A \times C \mid \exists b \in B, (a,b) \in R \text{ and } (b,c) \in S\}$.
  \item Identity: $\Delta_A = \{(a,a) \mid a \in A\} \subseteq A \times A$.
\end{itemize}
The isomorphisms in $\mathbf{Rel}$ are the relations $R: A \to B$ such that $R^{-1} \circ R = \Delta_A$ and $R \circ R^{-1} = \Delta_B$—these are exactly the bijections (functions that are both injective and surjective, viewed as their graphs).

The functor $\Set \to \mathbf{Rel}$ sends $f: A \to B$ to its graph $\Gamma_f = \{(a, f(a)) \mid a \in A\}$.
\end{exercise}

\section*{Chapter 3: Functors}

\begin{exercise}[3.$\star$: $G$-sets]
A group $G$ viewed as a category $\mathbf{B}G$ has one object $\bullet$ with $\mathbf{B}G(\bullet,\bullet) = G$. A functor $F: \mathbf{B}G \to \Set$ consists of a set $X = F(\bullet)$ and, for each $g \in G$, a function $F(g): X \to X$, satisfying $F(gh) = F(g) \circ F(h)$ and $F(e) = \id_X$. This is exactly a (left) $G$-set (a set $X$ with a group action $G \times X \to X$, $g \cdot x = F(g)(x)$).

A natural transformation $\alpha: F \Rightarrow G'$ between two $G$-sets $F, G'$ has a single component $\alpha_\bullet: X \to Y$ satisfying $G'(g)(\alpha_\bullet(x)) = \alpha_\bullet(F(g)(x))$ for all $g \in G, x \in X$—this is exactly a $G$-equivariant map.

So $[\mathbf{B}G, \Set] \cong G$-$\Set$ (the category of $G$-sets and equivariant maps).
\end{exercise}

\begin{exercise}[3.$\star\star$: Categorical Cayley Theorem]
Let $\cat{C} = \mathbf{B}G$. The Yoneda embedding $\yo: \mathbf{B}G \to [\mathbf{B}G\op, \Set]$ sends the single object $\bullet$ to the presheaf $\mathbf{B}G(-,\bullet) = G$ viewed as a right $G$-set (action by right multiplication). The Yoneda embedding is fully faithful, hence $G$ embeds into the automorphism group of the functor $\yo(\bullet)$, which is the group of natural automorphisms of $G$ as a $G$-set. These natural automorphisms are exactly the left multiplications by elements of $G$. So $G$ embeds into $\mathrm{Aut}(G\text{-set}) \cong G$ (left translations)—recovering Cayley's theorem ($G$ embeds into $\mathrm{Sym}(G)$).
\end{exercise}

\section*{Chapter 4: Natural Transformations}

\begin{exercise}[4.$\star\star$: Free cocompletion]
We sketch the proof. Given $F: \cat{C} \to \cat{D}$ with $\cat{D}$ cocomplete, define $\bar{F}: [\cat{C}\op,\Set] \to \cat{D}$ by $\bar{F}(H) = \Lan_\yo F(H)$, the left Kan extension of $F$ along the Yoneda embedding. By the pointwise formula:
\[
  \bar{F}(H) = \colim_{(c, \yo(c) \to H) \in \int H} F(c),
\]
where $\int H$ is the category of elements of $H$. Since $\bar{F}$ is defined as a colimit, it preserves colimits (left adjoints preserve colimits; alternatively, $\bar{F}$ is a left Kan extension along a fully faithful functor and hence is the unique cocontinuous extension).

For uniqueness: if $G: [\cat{C}\op,\Set] \to \cat{D}$ preserves colimits and $G \circ \yo \cong F$, then $G(H) = G(\colim_{\int H} \yo(c)) \cong \colim_{\int H} G(\yo(c)) \cong \colim_{\int H} F(c) = \bar{F}(H)$, using the fact that $H \cong \colim_{\int H} \yo(c)$ (the co-Yoneda lemma).
\end{exercise}

\section*{Chapter 5: Limits and Colimits}

\begin{exercise}[5.$\star\star$: Filtered colimits commute with finite limits in $\Set$]
We prove that for filtered $\cat{J}$ and finite $\cat{K}$, $\colim_j \lim_k D(j,k) \cong \lim_k \colim_j D(j,k)$ in $\Set$.

Elements of $\colim_j \lim_k D(j,k)$ are equivalence classes $[(j, (x_k)_{k \in \cat{K}})]$ where $(x_k)_{k \in \cat{K}} \in \lim_k D(j,k)$. Elements of $\lim_k \colim_j D(j,k)$ are families $([j_k, y_k])_{k \in \cat{K}}$ in $\prod_k \colim_j D(j,k)$, compatible with the morphisms in $\cat{K}$.

Construct the natural map $\phi: \colim_j \lim_k \to \lim_k \colim_j$ by $[(j, (x_k))] \mapsto ([(j, x_k)])_k$.

Injectivity: If $\phi([(j, (x_k))]) = \phi([(j', (x'_k))$ then for each $k$, $[(j, x_k)] = [(j', x'_k)]$ in $\colim_j D(-,k)$, so there exist $j_k$ and morphisms $j \to j_k \leftarrow j'$ with $D_{j \to j_k}(x_k) = D_{j' \to j_k}(x'_k)$. Since $\cat{K}$ is finite and $\cat{J}$ is filtered, we can find a single $j''$ dominating all $j_k$, and then $[(j,(x_k))] = [(j', (x'_k))]$ in $\colim_j \lim_k$.

Surjectivity: Given a compatible family $[(j_k, y_k)]_k$, use filteredness of $\cat{J}$ to find $j''$ dominating all $j_k$, then use compatibility and finiteness of $\cat{K}$ to find a single element in $\lim_k D(j'',k)$ mapping to the given family.
\end{exercise}

\section*{Chapter 6: Adjoint Functors}

\begin{exercise}[6.$\star$: Reflective subcategories and limits]
Let $i: \cat{D} \hookrightarrow \cat{C}$ be a reflective subcategory with $L \dashv i$. We show $\cat{D}$ is closed under limits in $\cat{C}$.

Let $D: \cat{J} \to \cat{D} \hookrightarrow \cat{C}$ be a diagram with limit $(\ell, \pi)$ in $\cat{C}$. The unit $\eta_\ell: \ell \to i(L\ell)$ is a morphism in $\cat{C}$. Since $i$ preserves limits (it's a right adjoint), $i(L\ell)$ is the limit of $iLD$ in $\cat{C}$... Actually we argue directly: consider the cone $(\ell, \pi_j: \ell \to D(j) = i(D'(j)))$ in $\cat{C}$. The unit gives $\eta_\ell: \ell \to iL\ell$; composing with projections gives a cone from $iL\ell$ to $D$. Since $L\dashv i$, maps $L\ell \to D'(j)$ in $\cat{D}$ correspond to maps $\ell \to i(D'(j))$ in $\cat{C}$. The limit condition in $\cat{C}$ and the adjunction bijection show $L\ell$ is the limit of $D'$ in $\cat{D}$. But then $i(L\ell) \cong \ell$ (since $\ell$ is the limit in $\cat{C}$ and $i$ creates limits for reflective inclusions), so $\ell \in \cat{D}$.

For colimits: $\cat{D}$ is not necessarily closed under colimits in $\cat{C}$ (the colimit in $\cat{C}$ may not be a $\cat{D}$-object). But the reflector $L$ can be applied: the colimit in $\cat{D}$ is $L(\colim_\cat{C} D)$.
\end{exercise}

\begin{exercise}[6.$\star\star$: GAFT]
\textit{Construction of the left adjoint.}

Given $G: \cat{C} \to \cat{D}$ limit-preserving with solution sets $\{d \to G(c_i)\}_{i \in I_d}$ for each $d$, we construct $Fd \in \cat{C}$ as follows.

Form the small category $\mathcal{S}_d$ whose objects are the solution set pairs $(i, c)$ where $c \in \cat{C}$ and the morphism $d \to G(c_i) \to G(c)$ factors through $G$. More precisely, $\mathcal{S}_d$ is the full subcategory of the comma category $(d \downarrow G)$ generated by the solution set. Let $Fd = \lim_{\mathcal{S}_d\op} \pi$ where $\pi: \mathcal{S}_d \to \cat{C}$ is the projection. The limit exists by completeness of $\cat{C}$.

The canonical cone from $d$ to $G(Fd) = G(\lim) \cong \lim G$ (using $G$ preserves limits) together with the solution set property shows $Fd$ is the desired initial object in $(d \downarrow G)$, and the collection $d \mapsto Fd$ assembles into the left adjoint functor $F$.
\end{exercise}

\section*{Chapter 7: Monads}

\begin{exercise}[7.$\star\star$: Beck's theorem, one direction]
\textit{Assuming the adjunction is monadic (K is an equivalence), show U reflects isomorphisms.}

Let $f: d \to d'$ be a morphism in $\cat{D}$ such that $Uf$ is an isomorphism in $\cat{C}$. Since $K: \cat{D} \to \cat{C}^T$ is an equivalence and $U^T \circ K = U$, we have $U^T(K(f))$ is an isomorphism. Since $U^T$ creates isomorphisms (algebras maps that are isomorphisms on underlying objects are isomorphisms of algebras), $K(f)$ is an isomorphism in $\cat{C}^T$. Since $K$ is an equivalence (hence reflects isomorphisms), $f$ is an isomorphism in $\cat{D}$.

\textit{Existence of coequalizers of U-split pairs.} Given $f, g: d \to d'$ in $\cat{D}$ with sections $s, t$ in $\cat{C}$ as in the definition of a $U$-split coequalizer, transport the coequalizer along $K$: the corresponding data in $\cat{C}^T$ gives a coequalizer there (since $\cat{C}^T$ has all reflexive coequalizers: the coequalizer of a $T$-algebra map pair that is split as a $\cat{C}$-map always exists). Pulling back along $K$ gives the coequalizer in $\cat{D}$.
\end{exercise}

\section*{Chapter 8: Kan Extensions}

\begin{exercise}[8.$\star$: Density theorem for Yoneda]
By the co-Yoneda lemma, for any presheaf $F: \cat{C}\op \to \Set$:
\[
  F \cong \int^c F(c) \cdot \yo(c) = \colim_{(c, x) \in \int F} \yo(c),
\]
where $\int F$ is the \textbf{category of elements} of $F$: objects are pairs $(c, x)$ with $x \in Fc$; morphisms $(c,x) \to (c',x')$ are morphisms $f: c \to c'$ in $\cat{C}$ with $(Ff)(x') = x$ (note the contravariancy).

The colimit is over the canonical diagram $\int F \to [\cat{C}\op,\Set]$, $(c,x) \mapsto \yo(c)$. The colimiting cocone is given by the Yoneda-corresponding maps: for each $(c,x) \in \int F$, the element $x \in Fc$ corresponds (by Yoneda) to a natural transformation $\yo(c) \Rightarrow F$; these are compatible with morphisms in $\int F$, giving a cocone over $F$ that is universal.
\end{exercise}

\begin{exercise}[8.$\star\star$: Absolute Kan extensions from adjoints]
Let $F \dashv G$ with $F: \cat{C} \to \cat{D}$, $G: \cat{D} \to \cat{C}$, and $H: \cat{D} \to \cat{E}$.

Claim: $\Ran_F H \cong HG$ (absolutely, i.e., for any functor applied to the result).

Proof: $\Nat(L, \Ran_F H) \cong \Nat(LF, H)$ by the universal property of $\Ran_F$. We have $\Nat(LF, H) \cong \Nat(L, HG)$ by the adjunction $F \dashv G$ (since morphisms of functors $LF \Rightarrow H$ correspond to morphisms $L \Rightarrow HG$ via: $\alpha \mapsto (H\varepsilon \circ \alpha F)$, where the bijection uses the unit/counit). By Yoneda, $\Ran_F H \cong HG$.

The counit of this Kan extension is $\varepsilon: HGF \Rightarrow H$ (compose the counit of the adjunction with $H$). This is absolute because the bijection $\Nat(L, HG) \cong \Nat(LF, H)$ is natural in $H$: for any $K: \cat{E} \to \cat{E}'$, $\Nat(L, KHG) \cong \Nat(LF, KH)$ trivially, showing $K(HG) = (KH)G \cong \Ran_F(KH)$.
\end{exercise}

\section*{Chapter 9: Enriched Category Theory}

\begin{exercise}[9.$\star\star$: Enriched Yoneda Lemma]
Let $(\cat{V}, \otimes, I)$ be a closed symmetric monoidal category with internal hom $[-,-]: \cat{V}\op \times \cat{V} \to \cat{V}$. Let $\cat{C}$ be a $\cat{V}$-category and $a \in \cat{C}$.

The enriched Yoneda lemma states: for any $\cat{V}$-functor $F: \cat{C} \to \cat{V}$,
\[
  [\cat{C}, \cat{V}](\cat{C}(a,-), F) \cong Fa \quad \text{in } \cat{V},
\]
naturally in $a$ and $F$. The left side is the $\cat{V}$-object of $\cat{V}$-natural transformations, defined as:
\[
  [\cat{C},\cat{V}](F,G) = \int_c [Fc, Gc] \quad (\text{an end in } \cat{V}).
\]

Proof sketch: The natural bijection (on global elements $I \to [\cat{C},\cat{V}](\cat{C}(a,-),F)$, i.e., ordinary natural transformations $\cat{C}(a,-) \Rightarrow F$) is exactly the ordinary Yoneda lemma at the level of global elements. The full enriched statement follows by an enriched version of the same argument using the closed structure of $\cat{V}$.

When $\cat{V} = \Set$ and $I = \{*\}$, global elements are just elements, and $[\cat{C},\Set](\cat{C}(a,-), F) = \Nat(\cat{C}(a,-), F)$ is exactly the ordinary set of natural transformations, recovering Theorem~\ref{thm:yoneda}.
\end{exercise}

\section*{Chapter 10: Topos Theory}

\begin{exercise}[10.$\star\star$: LT topologies and Grothendieck topologies]
A Grothendieck topology on $\cat{C}$ is a function $J$ assigning to each $c \in \cat{C}$ a collection $J(c)$ of sieves on $c$ (called \textbf{covering sieves}), satisfying the maximality, stability, and transitivity axioms.

Given a Lawvere-Tierney operator $j: \Omega \to \Omega$ on $[\cat{C}\op, \Set]$ (where $\Omega(c) = \{$ sieves on $c \}$), define $J(c) = \{S \in \Omega(c) \mid j_c(S) = \text{max sieve on } c\}$—the covering sieves are those sent to the maximal sieve by $j$. The idempotence, unit, and meet-preservation of $j$ translate to the axioms of a Grothendieck topology.

Conversely, a Grothendieck topology $J$ gives a LT operator $j: \Omega \to \Omega$ by $j_c(S) = \{f: d \to c \mid f^*(S) \in J(d)\}$ (the ``$J$-closure'' of $S$). One verifies this satisfies the LT axioms.

These constructions are inverse, giving a bijection between LT topologies on $[\cat{C}\op,\Set]$ and Grothendieck topologies on $\cat{C}$.
\end{exercise}


%% ============================================================
\backmatter
%% ============================================================

%% ---- Bibliography ----------------------------------------
\printbibliography[heading=bibintoc, title={Bibliography}]

%% ---- Index (optional) ------------------------------------
%% \printindex

\end{document}
